\documentclass[12pt]{article}

% --- replication kit: resolve generated tables/figures from ../../output/ ---
\makeatletter\def\input@path{{../../output/tables/}{../../output/figures/}}\makeatother

\usepackage{comment}
\usepackage{placeins}
%\usepackage{natbib}
\usepackage{natbib,natbibspacing}
% \bibliographystyle{econometrica}
\bibliographystyle{jpe}
%\bibliographystyle{aer}
% \usepackage{biblatex}
%\bibliographystyle{qje}
% \newcounter{forroman}
% \usepackage{ulem}
% \newcommand{\asroman}[1]{
%    \setcounter{forroman}{#1}
%    \Roman{forroman}
%}

%\bibliographystyle{kluwer}
%\bibliographystyle{cje}
%\bibliographystyle{agsm}




\usepackage{amsmath,amssymb,longtable,appendix,graphicx,subfigure,multirow}
\graphicspath{{../../output/figures/}}   % after graphicx (replication kit)

\renewcommand{\baselinestretch}{1.5}

\frenchspacing
\sloppy

\usepackage{caption}
\usepackage{rotating}
\usepackage{nicefrac}
\usepackage{ctable}
\usepackage{float}

\usepackage{setspace}
\newcommand{\1}{\mathbf{1}}
\newcommand{\E}{\mathbb{E}}
\newcommand{\RR}{\mathbb{R}}
\newcommand{\Prob}{\mathbb{P}}
\newcommand{\C}{\mathcal{C}}
\newcommand{\wb}{\bar{w}}
\newcommand{\RP}{\mathbb{R}_+}
\newcommand{\VV}{\mathcal{V}}
\newcommand{\FF}{\mathcal{F}}
\newcommand{\JJ}{\mathcal{J}}

\newtheorem{theorem}{Theorem}[section]
\newtheorem{conjecture}{Conjecture}
\newtheorem{lemma}{Lemma}
\newtheorem{proposition}{Proposition}
\newtheorem{corollary}[theorem]{Corollary}
\newtheorem{assumption}{Assumption}
\newtheorem{aprime}{Assumption}


\newcommand{\beq}{\begin{eqnarray}}
\newcommand{\eeq}{\end{eqnarray}}
\newcommand{\beqn}{\begin{eqnarray*}}
\newcommand{\eeqn}{\end{eqnarray*}}


\newenvironment{proof}[1][Proof]{\begin{trivlist}
\item[\hskip \labelsep {\bfseries #1}]}{\end{trivlist}}
\newenvironment{definition}[1][Definition]{\begin{trivlist}
\item[\hskip \labelsep {\bfseries #1}]}{\end{trivlist}}
\newenvironment{example}[1][Example]{\begin{trivlist}
\item[\hskip \labelsep {\bfseries #1}]}{\end{trivlist}}
\newenvironment{remark}[1][Remark]{\begin{trivlist}
\item[\hskip \labelsep {\bfseries #1}]}{\end{trivlist}}
\newenvironment{claim}[1][Claim]{\begin{trivlist}
\item[\hskip \labelsep {\bfseries #1}]}{\end{trivlist}}
\newenvironment{changemargin}[2]{%
\begin{list}{}{%
\setlength{\topsep}{0pt}%
\setlength{\leftmargin}{#1}%
\setlength{\rightmargin}{#2}%
\setlength{\listparindent}{\parindent}%
\setlength{\itemindent}{\parindent}%
\setlength{\parsep}{\parskip}%
}%
\item[]}{\end{list}}

\newcommand{\qed}{\nobreak \ifvmode \relax \else
      \ifdim\lastskip<1.5em \hskip-\lastskip
      \hskip1.5em plus0em minus0.5em \fi \nobreak
      \vrule height0.75em width0.5em depth0.25em\fi}


\newcommand{\ccr}{\color{red}}
\newcommand{\ccb}{\color{blue}}
\newcommand{\dat}{\delta^A_t}
\newcommand{\dbt}{\delta^B_t}
\newcommand{\vat}{v^A_t}
\newcommand{\vbt}{v^B_t}
\newcommand{\vatt}{\tilde{v}^A_t}
\newcommand{\vbtt}{\tilde{v}^B_t}
\newcommand{\xat}{x^A_t}
\newcommand{\xbt}{x^B_t}
\newcommand{\tat}{\theta^A_t}
\newcommand{\tbt}{\theta^B_t}
\newcommand{\uat}{u^A_t}
\newcommand{\ubt}{u^B_t}
\newcommand{\uatp}{u^A_{t+1}}
\newcommand{\ubtp}{u^B_{t+1}}
\newcommand{\uatt}{\tilde{u}^A_t}
\newcommand{\ubtt}{\tilde{u}^B_t}
\newcommand{\uattp}{\tilde{u}^A_{t+1}}
\newcommand{\ubttp}{\tilde{u}^B_{t+1}}
\newcommand{\tatt}{\tilde{\theta}^A_t}
\newcommand{\tbtt}{\tilde{\theta}^B_t}
\newcommand{\sbt}{s^B_t}
\newcommand{\sat}{s^A_t}
\newcommand{\ebt}{e^B_t}
\newcommand{\eat}{e^A_t}
\newcommand{\nat}{n^A_t}
\newcommand{\nbt}{n^B_t}
\newcommand{\ust}{u^S_t}
\newcommand{\uct}{u^C_t}
\newcommand{\zst}{z^S_t}
\newcommand{\zct}{z^C_t}
\newcommand{\tst}{\theta^S_t}
\newcommand{\tct}{\theta^C_t}
\newcommand{\ost}{\Omega^S_t}
\newcommand{\oct}{\Omega^C_t}
\newcommand{\pat}{p^A_t}
\newcommand{\pbt}{p^B_t}


\usepackage{titlesec}
\titlespacing\section{0pt}{9pt plus 1pt minus 1pt}{8pt plus 1pt minus 1pt}
\titlespacing\subsection{0pt}{9.0pt plus 1pt minus 1pt}{8pt plus 1pt minus 1pt}
\titlespacing\subsubsection{0pt}{8.0pt plus 1pt minus 1pt}{8pt plus 1pt minus 1pt}

%\titlespacing{command}{left spacing}{before spacing}{after spacing}[right]
%From the titlesec package
% spacing: how to read {12pt plus 4pt minus 2pt}
%           12pt is what we would like the spacing to be
%           plus 4pt means that TeX can stretch it by at most 4pt
%           minus 2pt means that TeX can shrink it by at most 2pt
%%%%%%%%%%%%%%%%%%%%%%%%%%%%


%%%%%%%%%%%%%%%%%%%%%%%%%%%%
% Adjust Spacing Around Equations
%%%%%%%%%%%%%%%%%%%%%%%%%%%%
\expandafter\def\expandafter\normalsize\expandafter{%
    \normalsize
    \setlength\abovedisplayskip{5pt}
    \setlength\belowdisplayskip{5pt}
    \setlength\abovedisplayshortskip{5pt}
    \setlength\belowdisplayshortskip{5pt}
}


\setlength{\topmargin}{-0.7in} \setlength{\oddsidemargin}{-0.1 in}
\setlength{\textwidth}{6.7 in} \setlength{\textheight}{9.3in}



\begin{document}


\renewcommand{\baselinestretch}{1}

\title{{ Unemployment Benefits and Unemployment in the Great Recession: The Role of Macro Effects\footnote{This draft: 11 February 2026.
Marcus Hagedorn, our dear friend, cherished co-author, and one of the leading macroeconomists of our time, passed away in May 2025, shortly after submitting the last major revision of this paper. The remaining co-authors dedicate this paper to his memory.\\
\indent This paper previously circulated under the title ``Unemployment Benefits and Unemployment in the Great Recession: The Role of Equilibrium Effects.''\\
\indent We would like to thank Manuel Arellano, Dmitry Arkhangelsky, Bob Hall, Sam Schulhofer-Wohl, and seminar participants at Bocconi, Bonn, Cambridge, U.S. Census Bureau, Columbia, Edinburgh, EIEF, Erasmus, EUI, USC, Maryland, Penn State, UPenn, Princeton, Pompeu Fabra, Royal Holloway, Stanford, Toronto, Toulouse, UCL, UConn, Wisconsin, CUNY Graduate Center, Greater Stockholm Macro Group, Federal Reserve Banks of Cleveland, New York, and Philadelphia, 2013 conference on Macroeconomics Across Time and Space, 2013 SED, 2013 NBER Summer Institute (EFCE, EFMB, EFRSW groups), 2013 North American Summer Meeting of Econometric Society, 2013 Minnesota Workshop in Macroeconomic Theory, 15th IZA/CEPR European Summer Symposium on Labor Economics, Mannheim conference on ``Financial Frictions and Real Economy,'' 4th Ifo Conference on "Macroeconomics and Survey Data", 2014 ASSA Meetings, 2014 NBER Public Economics Program Meeting, 2014 Cowles Foundation Summer Conference on ``Structural Empirical Microeconomic Models'', the 2018 ERMAS conference, and the 2022 Annual Meeting of the Austrian Economics Association for their comments. We are especially grateful to June Shelp, at The Conference Board, for her help with the HWOL data. The opinions expressed herein are those of the authors and not necessarily those of the Federal Reserve Bank of New York or the Federal Reserve System. Support from the National Science Foundation Grants No. SES-0922406 and SES-1357903, FRIPRO Grant No. 250617 and No.334487, the Ragnar S\"oderbergs stiftelse, and the European Research Council grant No. 759482 under the European Union's Horizon 2020 research and innovation programme is gratefully acknowledged.}}
\author{Marcus Hagedorn\thanks{University of Oslo, Department of Economics. Email: marcus.hagedorn07@gmail.com.} \hspace{0.5cm} Fatih Karahan\thanks{Governor, Central Bank of the Republic of T\"{u}rkiye} 
\hspace{0.5cm} Iourii Manovskii\thanks{University of Pennsylvania, Department of Economics.  Email: manovski@econ.upenn.edu.} 
\hspace{0.5cm} Kurt Mitman\thanks{CEMFI and IIES, Stockholm University. Email: kurt.mitman@iies.su.se.}}}
\date{}
\maketitle
\thispagestyle{empty}


\vspace{-0.5cm}\begin{abstract}
Equilibrium labor market theory suggests that unemployment benefit extensions affect unemployment by impacting the unemployed's job search decisions and employers' job creation decisions. The existing empirical literature focused only on the former effect. We develop a new methodology necessary to incorporate the measurement of the latter effect. Implementing this methodology in the data, we find that benefit extensions raise equilibrium wages and lead to a sharp contraction in vacancy creation and employment and a rise in unemployment.

\vspace{0.6cm} \noindent\textbf{Keywords:} Unemployment insurance, Unemployment, Employment, Vacancies, \\

\vspace{-0.4cm} \hspace{1.5cm} Wages, Search, Matching

\vspace{0.4cm} \noindent\textbf{JEL codes:} E24, J63, J64, J65


\end{abstract}

\newpage

%%%%%%%%%%%%%%%%%%%%%%%%%%%%%%%%%%%%%%%%%%%%%%%%%%%%%%%%%%%%%%%%%%%%%%%%%%%%%%%%%%
%--------------------------------------------------------------------------------%
%%%%%%%%%%%%%%%%%%%%%%%%%%%%%%%%%%%%%%%%%%%%%%%%%%%%%%%%%%%%%%%%%%%%%%%%%%%%%%%%%%
\onehalfspacing
\renewcommand{\baselinestretch}{1.0}

\setcounter{page}{1}
\section{Introduction}

Unemployment in the U.S. rose dramatically during the Great Recession and remained persistently high long after its official end. In response to this crisis, policymakers implemented unprecedented extensions of unemployment benefits, increasing the maximum duration from 26 to 99 weeks. While these measures were lauded for reinforcing the social safety net, critics \citep[e.g.,][]{barro:10,mulligan:12} argued that they disincentivized work and prolonged unemployment by imposing implicit taxes on market activity. However, the existing empirical literature found only modest impacts of extended benefits on job search behavior, seemingly rebutting the arguments that extended benefits could significantly prolong unemployment. These studies, however, focus solely on workers' job search decisions and cannot fully evaluate the policy's impact because they exclude the possibility that this policy could also impact labor demand. In this paper, we develop a novel measurement methodology incorporating employers' vacancy creation decisions to capture the effects of unemployment benefit extensions on labor demand. We apply our estimator and exploit a policy discontinuity at state borders to measure the effect of unemployment benefit extensions on labor market variables in the U.S. during the Great Recession, providing a more complete empirical assessment of the labor market implications of this policy response.\footnote{\citet{feldstein:78} was probably the first to call for the analysis of the impact of policies on equilibrium market-level unemployment rates rather than merely individual transitions out of unemployment.}


The following stylized decomposition helps illustrate the two margins:
\begin{equation}
\label{eq:decomposition}
\text{Job-finding rate}_{it} \quad = \underbrace{s_{it}}_{\text{search intensity}}\times\underbrace{f(\theta_{t})}_{\text{finding rate per unit of } s}
\end{equation}
In words, the probability that an individual $i$ finds a job at time $t$ depends on how hard that individual searches and how selective he is in his acceptance decisions, which is captured by the ``search intensity'' component $s_{it}$. It also depends on the aggregate labor market conditions $\theta_t$ that determine how easy it is to locate jobs by expending a unit of search effort. To use an extreme example, if there are no job vacancies created by employers, $f(\theta_t)=0$, no amount of search effort by an unemployed worker would yield a positive probability of obtaining a job.

Changes in unemployment benefit policies affect both the search intensity of unemployed workers - the \emph{micro effect}, and the aggregate job-finding rate per unit of search effort through \emph{general equilibrium macro effects}. Indeed, in the classic equilibrium search framework of \citet{mortensen+pissarides:94}, the primary analytical device used by economists to study the determination of unemployment, the response of unemployment to changes in benefits is mainly driven by the response of employers’ decisions of whether and how many jobs to create and not by the impact on workers’ job search and acceptance decisions. The logic of the model is simple. Everything else equal, extending unemployment benefits exerts an upward pressure on the equilibrium wage. This lowers the profits employers receive from filled jobs, leading to a decline in vacancy creation. Lower vacancies imply a lower job-finding rate for workers, which leads to an increase in unemployment.

 Starting with \citet{millard+mortensen:97} and \citet{shi+wen:99}, the evidence on the magnitude of equilibrium macro effects is predominantly based on the estimation of structural models based on \citet{mortensen+pissarides:94}. For example, a structural analysis based on this model in \citet{krause+uhlig:12} reveals a large reduction in unemployment and an increase in vacancy creation due to the benefit duration cut (known as Hartz IV reform) in Germany. 

However, the measurement of the equilibrium macro effects relies on the details of the model specification and on the value assigned to notoriously difficult-to-estimate parameters, e.g., the flow utility obtained by unemployed workers \citep[][]{hagedorn+manovskii:08}. Our objective in this paper is to directly measure in the data the impact of unemployment benefits on the labor market variables of interest without having to rely on the estimate of the flow utility of the unemployed and without having to fully specify the model. The empirical strategy we develop is, however, consistent with a fully specified model.

Incorporating the measurement of the macro effect in a policy evaluation requires us to develop a novel empirical methodology that overcomes two key threats to identifying the causal effects of benefits extensions: expectations and policy endogeneity. The macro effect measures the impact of policy on firms’ forward-looking decisions to create jobs, which, like any long-term investment decisions, are affected not only by the current policy but also by the expectation of possible future policy changes. As a result, the full expected sequence of future benefit durations, in addition to the contemporaneous benefit policy, affects current job creation.  Empirical analyses that fail to account for firms' expectations measure an uninterpretable mixture of the true policy effect and the unknown effect of labor market participants' expectations of future policies. The approach in the literature that attempts to identify and exploit surprise policy changes does not fully resolve this problem. For example, \citet{hagedorn+manovskii+mitman:15} evaluate the consequences of a one-time unexpected event when benefit extensions introduced during the Great Recession expired suddenly at the end of 2013. However, while the exact timing of the cut of extended benefits in 2013 was unexpected, people must have expected the benefit extensions to expire eventually, and those expectations must have been incorporated into job creation decisions before the unexpected cut in 2013. Thus, the estimated effect of the surprise benefit cut does not reveal the pure effect of benefit duration on unemployment without the researcher's ability to control for agents' expectations (for which data is not available). 

In contrast, in this paper, we develop a novel estimator that controls for agents' unobserved (by the researcher) expectations of future policies. This not only allows us to obtain directly interpretable estimates of the impact of unemployment benefits but also allows us to utilize the entire panel data on benefit duration across space and time without having to identify any unexpected policy changes.

Our strategy to separate the effects of current policies from expectations of future policies uses a semi-structural approach pioneered by \citet{hansen:82} and \citet{hansen_singelton:82}. Specifically, we only use firms' profit maximization condition together with rational expectations, free entry in vacancy creation, and efficient separations -- the assumptions made in \citeauthor{pissarides:00}' (\citeyear{pissarides:00}) textbook and most of the equilibrium search literature -- to infer the parameter of interest without taking a stand on other features of the (search) model. These fairly parsimonious theoretical restrictions are sufficient to estimate the effect of period $t$ benefit duration on period $t$ unemployment without having access to data on expectations. Specifically, we show that a regression of the appropriately quasi-differenced unemployment rate (as opposed to the unemployment rate itself) on current benefit duration reveals this effect because quasi-differencing eliminates the impact of the unobserved expected future benefits path, identifying the impact of only current benefits on current unemployment. This effect, which can be measured in the data, is the key object of interest. For example, to obtain the effect of persistent changes in benefits, we can simply add up the estimated contemporaneous effects, discounting them using the theoretical restrictions underpinning our estimator. %where the same theoretical restrictions on which our estimator relies deliver the correct discounting rate.
While we employ this methodology to answer a specific empirical question, it seems worth noting that the methodological advance we propose is applicable much more generally. In particular, a version of a quasi-differenced estimator that we develop can be used in any empirical analysis where expectations of future policies affect current decisions.\footnote{This applies to virtually all investment decisions, including investment in physical capital. As another example, in sticky-price New Keynesian models, firms' pricing decisions depend on the expected path of current and future marginal costs, potentially affected by future policies. In this environment, \citet{hagedorn+handbury+manovskii:15} show that quasi-differenced inflation can be used to control for expectations and measure the impact of fiscal policies on current marginal costs to reveal their stimulative effect.}

While economic theory restrictions help reveal the effect of interest that can be estimated without having data on expectations of future policies, its empirical implementation has to overcome the challenge that benefits themselves are endogenous and depend on (past) unemployment in the state and on federal policy. To separate the causal effect of benefits on unemployment from shocks that drive both unemployment and benefits, we use an interactive effects estimator and differencing between counties that border each other but belong to different states.


The interactive effects estimator \citep{bai:09} controls for the heterogeneous impact of various aggregate shocks during the Great Recession (the time period we study in the data) on different counties. For example, if two counties border each other but belong to two different states, one with right-to-work law and the other without, they are likely to be differentially affected by an aggregate shock to manufacturing. At the same time, this shock will likely induce spatial correlation across, e.g., counties in the same state and is likely to be correlated with state-level unemployment benefit duration (this is on top of the direct dependence of benefit duration on aggregate shocks that drive changes in federal unemployment insurance policy). Thus, if such a shock is not controlled for in the regression, it would be included in the error term and would lead to a biased estimate. A common approach in the literature uses Bartik instruments/controls, formed by interacting local industry shares and national industry growth rates, in an attempt to capture the heterogeneous impacts of aggregate shocks induced by differences in preexisting industrial structure across locations. However, this is a very narrow approach that does not control for the heterogeneous exposure of local economies to aggregate shocks through different mechanisms. For example, the aggregate financial crisis potentially had a different impact on states depending on their different foreclosure laws or household wealth distribution. An interactive effects estimator provides a natural, parsimonious, and very flexible way to control for every type of heterogeneous exposure to aggregate shocks that is relevant to local unemployment dynamics in the regression. It identifies latent aggregate shocks (factors) and county-specific exposure to them (factor loadings).
Using border county pairs to overcome the endogeneity of UI policy means that we compare the evolution of unemployment in counties that border each other but belong to different states. The (testable) identifying assumption is that after the heterogeneous exposure to aggregate shocks is controlled for through interactive effects, the remaining shocks affect neighboring counties similarly. At the same time, border counties differ in the level of benefits due to a policy discontinuity that arises since benefits are a function of state unemployment rates, which discontinuously change at the state border.\footnote{A fundamentally similar border county identification strategy was used, among others, by \citet{holmes:98} to identify the impact of right-to-work laws on location of manufacturing industry and by \citet{dube+etal:10} to identify the effect of minimum wage laws on earnings and employment of low-wage workers.} Thus, our empirical specification regresses the difference of quasi-differenced unemployment between pairs of border counties on differences in contemporaneous benefit durations and interactive effects. Our identifying assumption allows border counties to be different in terms of county-specific factors. A sufficient condition for identification is that state-level factors affect the two counties symmetrically so that the difference in state-level shocks does not affect the difference in quasi-differenced unemployment across the two counties.\footnote{For identification, we only need the weaker condition that any remaining difference in state-level factors is uncorrelated with the difference in benefits across the two counties.} In Section \ref{sec:endogeneity_test}, we develop and implement a test of this identifying assumption and find it to be supported by the data.


 
After describing our main data sources in Section \ref{sec:data}, we measure the effects of unemployment benefit extensions on unemployment in Section \ref{sec:benefits_and_unemployment}. We find that unemployment is significantly higher in the border counties belonging to the states that expanded unemployment benefit duration as compared to the counties just across the state border. The quantitative magnitude of this effect is so large that our estimates attribute a prominent role to benefit extensions in accounting for the persistence of high unemployment following the end of the Great Recession in 2009. It is important to note that our estimates do not imply that the large increase in unemployment at the onset of the Great Recession was due to extensions of unemployment benefits. However, we do find that the extensions of benefits contributed significantly to the slow decline of unemployment thereafter. For example, assuming perfect foresight of future benefits, our estimates imply that unemployment in 2011 would have been $2.15$ percentage points lower had benefits not been extended. While our estimates are undoubtedly large, they are smaller than consensus estimates in the existing empirical literature reviewed in Appendix \ref{sec:literature}, although we will argue that the interpretation of the existing estimates is unclear. We also conduct a host of other robustness exercises in Section \ref{sec:benefits_and_unemployment} and find that our results are stable and consistent across different specifications.

In Section \ref{sec:macro_effects} we assess whether the mechanisms embedded in the standard equilibrium labor market search model can provide a coherent rationalization of the effect of unemployment benefit extensions on unemployment that we document. We find that, consistent with the implications of our equilibrium search model, relative to the paired border county, job vacancy rates and employment fall significantly in counties experiencing larger benefit extensions. Moreover, wages of all workers, as well as wages of job stayers and wages of new hires, increase in counties experiencing larger benefit extensions relative to paired counties across the state border. The estimated magnitudes of these changes are also quantitatively consistent with our calibrated model.



Finally, while our estimate of the effects of unemployment benefit extensions is based on the difference across border counties, it is also desirable to be able to use the resulting coefficients to predict the effect of a nationwide extension. A potential concern is that when some states extend benefits more than others, economic activity may reallocate to states with, say, lower benefits. Our estimates would pick up this reallocation but would be absent if policy changes everywhere. Our results in Section \ref{sec:macro_effects} \citetext{and the formal analysis of aggregation in this setting in \citealp{hagedorn+manovskii+mitman:15}} alleviate such concerns. For example, in Section \ref{sec:UI_and_Emp}, we find a similarly large negative effect of unemployment benefit extensions on employment in non-tradable sectors, which are not subject to reallocation as we do for employment in tradable sectors, which are more exposed to potential reallocation. 

Numerous important but less crucial components of the analysis are described in the Online Appendix.

%%%%%%%%%%%%%%%%%%%%%%%%%%%%%%%%%%%%%%%%%%%%%%%%%%%%%%%%%%%%%%%%%%%%%%%%%%%%%%%%%%
%--------------------------------------------------------------------------------%
%--------------------------------------------------------------------------------%
%%%%%%%%%%%%%%%%%%%%%%%%%%%%%%%%%%%%%%%%%%%%%%%%%%%%%%%%%%%%%%%%%%%%%%%%%%%%%%%%%%
\section{Empirical Methodology}
\label{sec:methodology}

In this section, we develop our empirical methodology in three steps. First, we formulate the relationship of interest between unemployment and UI benefits and discuss the two main challenges to identification: expectations and endogeneity of benefits. Second, we develop a semi-structural approach following ideas pioneered by \cite{hansen_singelton:82} to overcome the identification challenges posed by expectations. Specifically, we only build on the premise that firms' job creation decisions are optimal, firms have rational expectations, vacancies are free to enter, and job separations are efficient. We don't take a stand on other model details, rendering our findings robust to potentially misspecifying those details. 
Third, we describe our identification strategy via border counties and the interactive effects estimator of \citet{bai:09}.



\subsection{\setstretch{0.2} \hspace{-.4cm} Estimating the Effect of UI on unemployment}
\label{sec:basic}
Our objective is to estimate the effect of current benefits on contemporaneous unemployment: 
\beq \label{eq:HU}
\log(u_{c,t}) = \alpha \log(b_{c,t}) + \epsilon_{c,t},
\eeq
where $\alpha$ is  the effect of period $t$ UI benefits $b$ in county $c$ on the unemployment rate $u$ in that county in period $t$, keeping fixed all other determinants of unemployment, that is $\epsilon_{c,t}$ is unchanged. Estimating such a regression would lead to a biased estimate because of two identification challenges. First, a simple endogeneity problem. Unemployment benefits are a function of state unemployment rate and federal policy responding to aggregate shocks. Benefits will thus be correlated with idiosyncratic state and aggregate shocks.






Second, $\alpha$ measures the effect of current benefits, keeping benefits in all other periods unchanged. However, any shock to the economy that increases current benefits $b_{c,t}$ might also lead to a different path of expected future benefits. This may affect hiring decisions of forward looking employers with rational expectations, which will impact unemployment today, as in the \cite{pissarides:00} model.

 Implementing the regression (\ref{eq:HU}) then does not identify the effect of current benefits on current unemployment but a combination of the effects of current benefits and all expected future benefits. 

A simple example might be helpful in illustrating the need to control for agents' expectations.  Consider a state that passes a law to permanently extend benefits, but that the extension does not come into effect for one year. In response to the passage of the law, firms expect lower future profits, and will thus reduce vacancy creation, increasing unemployment. By the time the new benefit policy actually takes effect, the labor market adjustment may be virtually complete. Thus, a naive difference-in-difference analysis conducted before and after the implementation of the reform would erroneously suggest no impact of unemployment benefits on unemployment, when, in fact, all of the adjustments would have occurred between passage and implementation. Indeed, the basic optimality of firms' decisions implies no discrete jumps in vacancy posting when anticipated policy changes are implemented. Thus, an observer may conclude that unemployment and vacancies are unrelated to benefits because they change dramatically when benefits do not and do not change when benefits change. It is only by controlling for the movements induced by the changes in expectations that the correct magnitude of the policy effect can be identified.

If we could observe contemporaneous measures of all relevant state and aggregate shocks and expected future benefits and shocks directly, we could add them as observables to (\ref{eq:HU}). However, data on such variables does not exist. In the next two subsections, we discuss how to overcome these identification challenges.


%%%%%%%%%%%%%%%%%%%%%%%%%%%%%%%%%%%%%%%%%%%%%%%%%%%%%%%%%%%%%%%%%%%%%%%%%%%%%%%%%%%%%
%-----------------------------------------------------------------------------------%
%%%%%%%%%%%%%%%%%%%%%%%%%%%%%%%%%%%%%%%%%%%%%%%%%%%%%%%%%%%%%%%%%%%%%%%%%%%%%%%%%%%%%
\subsection{\setstretch{0.2} \hspace{-.4cm} Using Economic Theory to Control for Expectations}
\label{sec:qd}
We start with the second challenge, which amounts to controlling for expectations of future benefits. We show that under a parsimonious subset of assumptions made in \citet{pissarides:00} Ch. 1 (and underlying most of the equilibrium search literature) one can estimate the effect of period $t$ benefits, $b_t$, on period $t$ unemployment, $u_t$, without relying on nonexistent data on expectations (in this subsection, to lighten the notation, we omit the county, $c$, subscript). 



\begin{assumption}\label{ass:firm_optimization}
    Firms maximize profits. 
\end{assumption}
Specifically, the value of a filled job for a firm, $J_t$, is given by the following recursive equation:
\begin{equation} \label{eq:Bellman_profit}
J_t = \pi_t + \beta(1-s_t) E_t J_{t+1},
\end{equation}
where $\pi_t$ is period $t$ profits from the job, $\beta$ is the discount factor, $s_t$ is the probability that the job ends and $E_t$ is the expectation operator using information available at time $t$. As we explain below, it is immaterial for our analysis whether the separation probability $s_t$ is exogenous or endogenous to unemployment benefits.

As in the standard \citet{pissarides:00} model, firms' period $t$ profits from employing a worker are given by the difference between workers' marginal product and the wage. The wage, in turn, is affected by the generosity of unemployment benefits available to the worker.\footnote{\label{foot:eqm_wage} Note that this is the equilibrium wage response to a change in benefits, combining the direct effect of benefits on wages and various indirect effects, which, in our empirical analysis, we can be deliberately agnostic about. For the response of vacancy creation, it is this equilibrium wage response that matters, and consequently, this is what we estimate (equation (\ref{eq:wageequation_estimate})) in Section 
\ref{sec:wages}.} Thus, up to a log-linear approximation with respect to the two state variables of the model, firms' profits from employing a worker are given by:
\begin{equation}\label{eq:profit}
\log(\pi_t) = \gamma_z \log(z_t) - \gamma_b \log(b_t) + \nu^c_t,
\end{equation}
where $\nu^c_t$ is an approximation error,\footnote{Approximation errors, by definition, have mean zero and are uncorrelated with all other variables.} $\log(z_t)$ is workers' mean zero log productivity, $b_t$ are unemployment benefits, and $\gamma_z$ and $\gamma_b$ are unknown coefficients (which the standard theory implies should both be positive but we do not impose such a restriction). As we discussed in Footnote \ref{foot:eqm_wage}, the two state variables $z_t$ and $b_t$ affect a firm's profits through their impact on a number of variables in the model and Equation (\ref{eq:profit}) refers to their relevant total effect.

Note that rearranging Eq. (\ref{eq:Bellman_profit}), we can express unobserved period $t$ profits as a \emph{quasi-differenced} value of the job:
\begin{equation} \label{eq:qd_profit}
\pi_t = J_t - \beta(1-s_t) E_t J_{t+1}.
\end{equation}
Value of the job and its expected future values are also unobserved, but in what follows we will relate them to observable variables.

\begin{assumption}\label{ass:rational_exp}
    Expectations are rational. 
\end{assumption}
This assumption allows us to rewrite Eq. (\ref{eq:qd_profit}) as
\begin{equation} \label{eq:qd_profit_noE}
\pi_t = J_t - \beta(1-s_t) J_{t+1} +\eta_t,
\end{equation}
where $\eta_t$ is a mean zero expectation error.

What is important to note here, is that quasi-differencing reveals current profits and eliminates the effect of future variables. To see this differently, without uncertainty, the period $t$ discounted present value of current profits is 
$J_t =  \pi_t + \beta(1-s_t) \pi_{t+1} + \beta^2 (1-s_t) (1-s_{t+1})  \pi_{t+2} + \ldots$ 
 and the period $t+1$ discounted present value of current profits is 
$J_{t+1} =  \pi_{t+1} + \beta(1-s_{t+1}) \pi_{t+2} + \beta^2 (1-s_{t+1}) (1-s_{t+2})  \pi_{t+3} + \ldots$.
Multiplying $J_{t+1}$ by $\beta(1-s_t)$ and subtracting the product from $J_t$, then eliminates any dependence on $s_t$ or future variables, $J_t - \beta(1-s_t) J_{t+1} = \pi_t $. Note that the presence of $s_t$ in the quasi-difference ensures that the terms multiplied by $s_t$ cancel so that the quasi-difference reveals current profits independent of $s_t$.

\begin{assumption}\label{ass:free_entry}
    Free entry into vacancy posting. 
\end{assumption}
Specifically, this assumption implies that firms create job vacancies until the expected cost of posting a vacancy equals the value of a filled job:
\begin{equation}\label{eq.free_entry}
q(\theta_t) J_t=c,
\end{equation}
where $q(\theta_t)$ is the probability to fill a vacancy, $\theta_t:=v_t/u_t$ is (observed) labor market tightness measured as the ratio of vacancies and unemployment, and $c$ is the cost of posting a vacancy. 

This assumption allows us to relate unobserved value of the job $J_t$ to observed labor market tightness, $\theta_t$. Specifically, using Eq. (\ref{eq.free_entry}) and approximating $q(\theta)$ as a function of $\log(\theta)$, yields
\begin{equation}\label{eq:jc}
\log(\theta_t)=\tilde{\kappa} \log(J_t) + \nu^a_t,
\end{equation}
for an approximation error $\nu^a_t$ and constant $\tilde{\kappa}$.

We now approximate $\log(J_t)$ as a function of $\log(\pi_t)$, $\log(J_{t+1})$, an expectational error $\eta_t$, and approximation error $\nu^b_t$ around the steady state with constant $\pi$, $s$, and $\pi = J (1-\beta (1-s))$, so that the previous equation becomes
\begin{equation}
\log(\theta_t)=\tilde{\kappa} \frac{\pi}{J} \log(\pi_{t}) + \tilde{\kappa} \beta (1-s_t) \log(J_{t+1}) + \nu^a_t + \nu^b_t + \eta_t.
\end{equation}
Using $\pi / J = (1-\beta (1-s))$ and the job creation decision (\ref{eq:jc}) for $t+1$, yields
\begin{equation}
\log(\theta_t)=\tilde{\kappa}  (1-\beta (1-s)) \log(\pi_{t}) + \beta (1-s_t) (\log(\theta_{t+1}) - \nu^a_{t+1})+ \nu^a_t + \nu^b_t + \eta_t.
\end{equation}

Substituting for $\log(\pi_{t})$ using Eq. (\ref{eq:profit}), we can thus isolate the effect of contemporaneous benefits and productivity by \emph{quasi-differencing} log labor market tightness:
\begin{equation} \label{eq:tildetheta}
 \tilde{\theta}_t  :=\log(\theta_t) - \beta (1-s_t) \log(\theta_{t+1}) =\tilde{\kappa}  (1-\beta (1-s)) (\gamma_z \log(z_t) - \gamma_b \log(b_t))  + \tilde{\eta}_t
\end{equation}
as a function of $\log(z_t)$ and $\log(b_t)$ and $ \tilde{\eta}_t = \nu^a_t - \beta (1-s_t) \nu^a_{t+1} + \nu^b_t + \nu^c_t + \eta_t$. We henceforth denote quasi-differenced variables with a tilde: for any observable $x$, $\tilde{x}_t:=\log(x_t)-\beta(1-s_t)\log(x_{t+1})$. Thus, up to an assumption about the value of $\beta$ any quasi-differenced observable can be measured directly in the data since $s_t$ is observable.


Equation (\ref{eq:tildetheta}) can be more compactly written as 
\beq \label{eq:regress}
\tilde{\theta}_t = \tilde{\alpha} \log(b_t) + \tilde{\epsilon}_t,
\eeq  
where $\tilde{\epsilon}_t =  \tilde{\kappa}  (1-\beta (1-s)) \gamma_z \log(z_t) + \tilde{\eta}_t$ and  $\tilde{\alpha}$ equals
\beq
- \gamma_b \tilde{\kappa} (1-\beta (1-s)). \label{eq:alpha1}
\eeq
The coefficient $\tilde{\alpha}$ describes the effect of a change in $\log(b)$ on quasi-differenced tightness, $\tilde{\theta}$. Note that we have not established the necesssary orthogonality conditions which imply that we can recover $\tilde{\alpha}$ from a regression of $\tilde{\theta}$ on $\log(b_t)$. 
The error term $\tilde{\epsilon}_t$ depends on $\log(z_t)$ and we describe our approach to overcome this potential source of endogeneity below in Section \ref{sec:bc}. 

Adding $\beta (1-s_t) \log(\theta_{t+1})$ to both sides of (\ref{eq:regress}) as
\beq
\theta_t = \tilde{\alpha} \log(b_t) + \beta (1-s_t) \log(\theta_{t+1}) + \tilde{\epsilon}_t,
\eeq
shows that $\tilde{\alpha}$ also describes the effect of $\log(b_t)$ on $\theta_t$, keeping $\beta (1-s_t) \log(\theta_{t+1})$ fixed.

Quasi-differencing achieves this controlling for expected benefits and allows recovering $\tilde{\alpha}$:  Equation (\ref{eq:regress}) states that neither future benefits nor separations rates affect quasi-differenced tightness if we know current benefits and $\tilde{\epsilon}_t$, i.e., future benefits would enter with a coefficient of zero at the RHS of (\ref{eq:regress}). Note again that adding future benefits as a regressor 
would not yield a coefficient of zero at this stage since we have not established the necessary orthogonality conditions. This means the coefficient on future benefits is zero only if we control for current benefits and $\tilde{\epsilon}_t$.



In terms of our econometric assumptions on data, the result of these derivations is an orthogonality condition which states that future benefits $b_{t+k+1}$ and separations rates $s_{t+k}$ add no further information once we condition on current benefits and $\tilde{\epsilon}_t$,
\beq
E( \tilde{\theta}_t \mid b_t,  \tilde{\epsilon}_t, E_t b_{t+k+1} ) = E( \tilde{\theta}_t \mid b_t,  \tilde{\epsilon}_t)  = \tilde{\alpha} \log(b_t) + \tilde{\epsilon}_t, \qquad \forall k \geq 0,  \label{cond_ortho1}
\eeq
and
\beq
E( \tilde{\theta}_t \mid b_t,  \tilde{\epsilon}_t, s_{t+k} ) = E( \tilde{\theta}_t \mid b_t,  \tilde{\epsilon}_t) = \tilde{\alpha} \log(b_t) + \tilde{\epsilon}_t, \qquad \forall k \geq 0,  \label{cond_ortho2}
\eeq



 The economic reasoning underlying this condition is a simple consequence of firms' optimal and time-consistent job creation decisions and rational expectations. Note first that assumptions (\ref{cond_ortho1}) and (\ref{cond_ortho2}) follow from 
 the model orthogonality conditions\footnote{The model implies that $Cov ( J_t - \beta(1-s_t) E_t J_{t+1}, E_t b_{t+k+1} \mid b_t) = 0$ and that the expectational error $E_t J_{t+1} - J_{t+1}$ is uncorrelated with expected benefits $E_t b_{t+k+1}$.}
\beq
E( J_t - \beta(1-s_t) J_{t+1} \mid b_t,  \tilde{\epsilon}_t, E_t b_{t+k+1} ) = E( J_t - \beta(1-s_t) J_{t+1} \mid b_t,  \tilde{\epsilon}_t)  \qquad \forall k \geq 0,  
\eeq
and
\beq
E( J_t - \beta(1-s_t) J_{t+1} \mid b_t,  \tilde{\epsilon}_t, s_k ) = E( J_t - \beta(1-s_t) J_{t+1} \mid b_t,  \tilde{\epsilon}_t)\qquad \forall k \geq 0.
\eeq

 
 Suppose now that there is a revision of expectation about next period's benefit level $b_{t+1}$ and that this is the only change and $b_t$ and all other expectations are unchanged. The recursive nature of the firm problem implies that the effect of the revised $b_{t+1}$ on $J_t$ is fully captured by $J_{t+1}$. The only direct effect of $b_{t+1}$ is on period $t+1$ profits, and the present value of period $t$ profits is affected only because it discounts these future profits. As a result quasi-differencing $J_t$, $J_t - \beta(1-s_t) J_{t+1}$, eliminates the effect of $b_{t+1}$ on $J_t$. By the same logic, the quasi-difference eliminates the effect of a change in $b_{t+2}, b_{t+3}, \ldots$
 
 

\subsubsection{\setstretch{0.2} \hspace{-.4cm} From Market Tightness to Unemployment}

It is straightforward to use the derivations for market tightness to obtain the corresponding expressions for unemployment, using 
\begin{assumption}\label{ass:eff_sep}
Workers and firms separate if the match surplus (the sum of values of the worker and the firm in the match minus the sum of values of a vacant job and unemployed worker) is negative.
\end{assumption}

We first notice that unemployment can be written as a function of $J_t$ 
\beq
\log(u_t) = \lambda_J \log(J_t) + \nu^u_t = \frac{\lambda_J}{\tilde{\kappa}} \log(\theta_t) - \frac{\lambda_J}{\tilde{\kappa}}  \nu^a_t + \nu^u_t, \label{eq:uJ}
\eeq

where we use that both the job-finding rate and the endogenous separation rate are functions of $J_t$ (Appendix \ref{app:end_sep} provides the details). 
While conditional on knowing $J_t$, other shocks $\nu^u_t$, e.g. matching function shocks or exogenous separation shocks, can affect the unemployment rate, benefits affect the unemployment rate only through its effects on $J_t$. 
In other words, knowing benefits provides no additional information to $J_t$. The second equality uses  (\ref{eq:jc})  
\begin{equation}
 \log(J_t)  = \frac{\log(\theta_t)-  \nu^a_t}{\tilde{\kappa}},
\end{equation}
implying that $\log(\theta_t)$ is also a sufficient statistic since  $\nu^a_t$ is an approximation error.

Quasi-differenced unemployment is given by
\beq \label{eq:qdunemp}
\tilde{u}_t & := & \log(u_t) - \beta (1-s_t) \log(u_{t+1}) \\ \notag
& = & \{ \frac{\lambda_J}{\tilde{\kappa}} \log(\theta_t) - \frac{\lambda_J}{\tilde{\kappa}}  \nu^a_t + \nu^u_t \} -
\beta (1-s_t) \{\frac{\lambda_J}{\tilde{\kappa}} \log(\theta_{t+1}) - \frac{\lambda_J}{\tilde{\kappa}}  \nu^a_{t+1} + \nu^u_{t+1} \} \\ \notag
&= & \frac{\lambda_J}{\tilde{\kappa}}  \tilde{\theta}_t  - \frac{\lambda_J}{\tilde{\kappa}} \{ \nu^a_t - \beta (1-s_t) \nu^a_{t+1} \} + \nu^u_t   - \beta (1-s_t) \nu^u_{t+1}, 
\eeq
We can further rewrite quasi-differenced unemployment using (\ref{eq:regress}),
\beqn
\tilde{u}_t &= & \frac{\lambda_J}{\tilde{\kappa}}  (\tilde{\alpha} \log(b_t) + \tilde{\epsilon}_t ) - \frac{\lambda_J}{\tilde{\kappa}} \{ \nu^a_t - \beta (1-s_t) \nu^a_{t+1} \} + \nu^u_t   - \beta (1-s_t) \nu^u_{t+1},
\eeqn
which we write more compactly as
\beq \label{eq:alpha}
\tilde{u}_t & = & \alpha \log(b_t) + \epsilon_t, \label{eq:Ureg}
\eeq
with
$$
\epsilon_t =  \frac{\lambda_J}{\tilde{\kappa}}  \tilde{\epsilon}_t  -  \frac{\lambda_J}{\tilde{\kappa}} \{ \nu^a_t - \beta (1-s_t) \nu^a_{t+1} \} + \nu^u_t   - \beta (1-s_t) \nu^u_{t+1}.
$$
and
$$
\alpha =  \frac{\lambda_J}{\tilde{\kappa}} \tilde{\alpha}  =  -  \lambda_J \gamma_b  (1-\beta (1-s)). 
$$
Analogously to tightness, note that $\alpha$ measures the effect of period $t$ benefits on period $t$ unemployment - our coefficient of interest. It is constructed from three terms. $\gamma_b$ is the effect of a permanent increase in unemployment benefits on the value of the filled job $J_t$. Multiplying this by $(1-\beta (1-s))$ yields the effect of an increase in benefits in period $t$ only on $J_t$. Finally, multiplying by $\lambda_J$ reveals the impact of the change in the value of the job from the change in period $t$ benefits on period $t$ unemployment. %We now establish that this coefficient does not depend on the expectation of future policies.


The sufficiency result carries over from $\tilde{\theta}$ to $\tilde{u}$. Equation (\ref{eq:Ureg}) is analogous to 
 equation (\ref{eq:regress}) as it states that neither future benefits nor separations rates affect quasi-differenced unemployment if we know current benefits and $\tilde{\epsilon}_t$, i.e. future benefits  would enter with a coefficient of zero at the RHS of (\ref{eq:Ureg}), with the same caveat as above on this not being a regression coefficient yet. 
 
Defining $\tilde{\nu}^u_t = \nu^u_t   - \beta (1-s_t) \nu^u_{t+1}$, the results for quasi-differenced unemployment in econometric terms mean, analogously to (\ref{cond_ortho1}) and (\ref{cond_ortho2}) for quasi-differenced tightness,
\beq
E( \tilde{u}_t \mid b_t, \tilde{\nu}^u_t, \tilde{\epsilon}_t, E_t b_{t+k+1} ) = E( \tilde{u}_t \mid b_t, \tilde{\nu}^u_t, \tilde{\epsilon}_t)  \qquad \forall k \geq 0,  \label{Ass:U1}
\eeq
and
\beq
E( \tilde{u}_t \mid b_t, \tilde{\nu}^u_t , \tilde{\epsilon}_t, s_k ) = E( \tilde{u}_t \mid b_t, \tilde{\nu}^u_t, \tilde{\epsilon}_t)\qquad \forall k \geq 0. \label{Ass:U2}
\eeq
We again condition on $\tilde{\epsilon}_t$ and $\tilde{\nu}^u_t$ since we have not yet discussed the orthogonality assumptions. 


This implies that we are able to measure the effect of a contemporaneous change in $b_t$ keeping the benefits at future dates unchanged.\footnote{We do not need further assumptions or results for separation rates beyond (\ref{Ass:U2}) which ensures that quasi-difference tightness is uncorrelated with separation rates (conditional on $b_t$) and 
the orthogonality assumptions explained below. We allow for a correlation of the unemployment rate and exogenous separation rate movements.}




\subsection{\setstretch{0.2} \hspace{-.4cm} Identification via Border Counties and Interactive Effects}
\label{sec:bc}
While quasi-differencing has solved the endogeneity problem of future expected variables, we still have to overcome the challenge of the potential contemporaneous correlation between $b_t$ and state or aggregate shocks. For concreteness, we will refer to state or aggregate productivity shocks, but these could be any shocks affecting state or aggregate economy which may induce changes in benefits, including demand shocks. Our strategy to identify the parameter $\alpha$ of equation (\ref{eq:alpha}) features two key elements: considering border counties and using an interactive effects estimator. Consider two counties, $c1$ in state $S1$, and a bordering county $c2$ in state $S2$. Our assumption on the data-generating process for quasi-differenced county unemployment is:
\beq \label{eq:dgp}
\tilde{u}_{c1,t} & = & \alpha \log(b_{S1,t}) + \lambda_{c1}'F_t + Z_{S1,c1,t}+Z_{S2,c1,t}+\epsilon_{c1,t},
\eeq
where $F_t$ is a vector of aggregate shocks, $\lambda_{c1}$ vector of county $c1$-specific loadings, $Z_{S1,c1,t}$ is the effect of productivity in state $S1$, $Z_{S2,c1,t}$ is the effect of productivity in the state $S1$ which borders county $c1$, and $\epsilon_{c1,t}$ is a county-specific error term that contains idiosyncratic county productivity, as well as county-specific approximation and expectation errors. The two challenges are that: 1) the impact of state productivity shocks on counties is unobservable; and 2) the aggregate shocks and their county loadings are also unobservable. We now discuss how we can overcome these two potential threats to identification.

\noindent
\underline{Border County Pairs}

To overcome the threat to identification posed by shocks to state productivity $Z$, we use a border-pair identification strategy. We consider two counties bordering each other but belonging to different states, thus having potentially different benefits available. Denote by $p$ the border-county pair. Differencing (\ref{eq:alpha}) between border counties within a pair yields:
\begin{equation}
\Delta \tilde{u}_{p,t} = \alpha \Delta b_{p,t} + \lambda_{p}'F_t + \Delta {Z}_{p,1,t}+\Delta {Z}_{p,2,t}+\Delta\epsilon_{p,t},
\end{equation}
where $\Delta$ is the difference operator over counties in the same pair. More specifically, if counties $c1$ and $c2$ are in the same border-county pair $p$, then $\Delta \tilde{u}_{p,t}=\tilde{u}_{c1,t}-\tilde{u}_{c2,t}$, $\lambda_p=\lambda_{c1}-\lambda_{c2}$, $\Delta {Z}_{p,1,t}=Z_{S1,c1,t}-Z_{S1,c2,t}$, ${Z}_{p,2,t}=Z_{S2,c1,t}- Z_{S2,c2,t}$, $\Delta \epsilon_{p,t}=\epsilon_{c1,t}-\epsilon_{c2,t}$ and, with a slight abuse of notation, $\Delta b_{p,t}=\log(b_{c1,t})-\log(b_{c2,t})$. 

Our baseline specification can then be written as:
\begin{equation}
\Delta \tilde{u}_{p,t} =  \alpha \Delta b_{p,t} +\lambda_{p}'F_{t}+ \nu_{p,t}. \label{eq:baseline_spec}
\end{equation}
Equation (\ref{eq:baseline_spec}), which will form the basis of our empirical strategy, differs from the standard specification in the literature in that the left-hand-side variable is the quasi-difference $\tilde{u}_{p,t}$ as opposed to simply $u_{p,t}$. As explained above, this is essential in order to take into account the forward-looking nature of vacancy postings, which carries over to unemployment.

The identifying assumption is that
\begin{equation} \label{assu:identify_main}
Corr (\nu_{p,t}, \Delta b_{p,t}) = 0,
\end{equation}

The motivation behind the border county design is that, conditional on identifying the heterogeneous effects of aggregate shocks (discussed below), the remaining shocks to states $S1$ and $S2$ affect the two countries symmetrically,
$$
 \Delta {Z}_{p,1,t}+\Delta {Z}_{p,2,t} = 0.
$$
This implies that the difference in benefits is uncorrelated with the error term, $\nu_{p,t}$, since benefits are a function of state shocks.
In contrast, a county-level regression of $\tilde{u}_{\hat{c},t}$ on $b_{\hat{c},t}$ would be biased since $b_{\hat{c},t}$ could be correlated with $Z_{S1,c1,t}+Z_{S2,c1,t}$. Instead of the previous sufficient condition, for identification we only need the weaker condition that after differencing across border counties any residual impact of state shocks is uncorrelated with the difference in benefits, i.e.:
\begin{equation}
\E\left[ \left(\Delta {Z}_{p,1,t}+\Delta {Z}_{p,2,t}\right)  \Delta b_{p,t} \right] = 0. \label{eq:bc_identify}
\end{equation}



Note that our identifying assumption Eq. (28) does not require border counties to be identical (conditional on the differences accounted for by the factor model) so that $\nu_{p,t}$ is pure measurement error. It is weaker than this as it allows counties to be different in terms of county-specific factors and only requires the difference in state-level factors (which influence state-level benefit policies) to be uncorrelated with the difference in benefits. Based on this insight, in Section \ref{sec:endogeneity_test} we will develop a test of the identifying assumption using observed measures of state-level productivity or demand. Implementing this test will let us conclude that our identifying assumption is satisfied in the data implying that differencing between border counties is indeed effective in eliminating the potential bias due to state-level policy endogeneity.


\noindent
\underline{Interactive Effects}

To overcome the threat to identification caused by aggregate shocks $F_t$, we use an interactive effects estimator developed by \citet{bai:09} and \citet{moon2015linear}. 

As we mentioned in the Introduction, various shocks affected the aggregate economy during the Great Recession. However, the same aggregate shocks may have a heterogeneous impact on different border county pairs. In this case, estimating the panel regression in Equation (\ref{eq:alpha}), perhaps with a set of county pair and time fixed effects, is problematic for inference.\footnote{See \citet{andrews:05} for the discussion of this problem in a cross-sectional regression. \citet{gobillon+magnac:13} establish the superior performance of the interactive effects estimator that we use relative to alternative methods used in the literature.} Fortunately, \citet{bai:09} has shown that consistency and proper inference can be obtained in a panel data context, such as ours, through the use of an interactive effects estimator.\footnote{Similar to \citet{gobillon+magnac:13}, given our identification assumptions, consistency and other asymptotic properties can be derived when $N \to \infty$ and $T \to \infty$.} In particular, we can estimate the latent vector of pair-specific factor loadings ($\lambda_{p}$) and vector of time-specific common factors ($F_{t}$) from our baseline specification (\ref{eq:baseline_spec}). 

As shown in \citet{bai:09}, this model nests additive time and county pair fixed effects as a special case.\footnote{This motivates our test between the interactive and additive effects models in Appendix \ref{app:OLS}. The additive effects model is rejected using both econometric tests proposed in \cite{bai:09}. Nevertheless, we report our baseline estimates using the additive effects estimator in Appendix \ref{app:OLS} and find that our substantive conclusions are largely unaffected.} It is, however, much more general and allows for a very flexible model of the heterogeneous time trends at the county pair level. The key to estimating $\alpha$ consistently is to treat the unobserved factors and factor loadings as parameters to be estimated. Our implementation is based on an iterative two-stage estimator described in Appendix \ref{app:implementation}.

We list the formal identifying assumptions in Appendix \ref{app:moon}, but discuss them briefly here. In addition to the standard assumptions for identification in OLS (finite second moments and orthogonality between benefits and the error term), we also require a noncollinearity assumption (analogous to the standard no-multicollinearity condition in OLS). Specifically, we require that benefits (our regressor) have significant variation across pairs $p$ and over time $t$ after projecting out all of the variation that can be explained by the true factor loadings for arbitrary factors. In other words, we can't have low-rank regressors. This requirement is satisfied in our data as the rank of $b$ is 20, much in excess of the rank of factors we estimate (1 to 4). %Finally, we have to assume a finite number of factors $r$. 


\subsubsection{Estimating the Number of Factors}
To implement the interactive effect estimator, we need to specify the number of factors. \citet{bai+ng:02} have shown that the number of factors in pure factor models can be consistently estimated based on the information criterion approach. \citet{bai:09} shows that their argument can be adapted to panel data models with interactive fixed effects. Thus, we define our criterion $CP$ as a function of the number of factors $k$ as:
\[
CP(k)=\hat{\sigma}^{2}(k)+\hat{\sigma}^{2}(\bar{k})\left[k\left(N+T\right)-k^{2}\right]\frac{\log\left(NT\right)}{NT},
\]
where $\bar{k}\geq r$ is the maximum number of factors, $N$ is the number of pairs, $T$ is the number of time observations, $\hat{\sigma}^{2}(k)$ is the mean squared error, defined as
\[
\hat{\sigma}^{2}(k)=\frac{1}{NT}\sum_{i=1}^{N}\sum_{t=1}^{T}\left(\Delta\tilde{x}_{p,t}-a\Delta b_{p,t}-\lambda_{i}^{'}\left(k\right)F_{t}\left(k\right)\right)^{2},
\]
and $F_{t}\left(k\right)$ and $\lambda_{i}^{'}\left(k\right)$ are the estimated factors and their loadings, respectively, when $k$ factors are estimated. To avoid collinearity, we set $\bar{k}$ to the minimum of seven and $T-1$, one less than the total number of time observations. Our estimator for the number of factors is then given by
\[
\hat{k}=\arg\min_{k\leq\bar{k}}CP(k).
\]

\subsubsection{Standard Errors}
To properly compute standard errors, we need to take into account potential correlation in the residuals. There are two possible sources of correlation. First, the outcomes that we are interested in (unemployment, vacancies, wages, etc.) are highly serially correlated. This aspect of the data may cause serial correlation in the errors. Second, the relevant variables can be correlated among border counties along the same state border because they are exposed to the same state-level policies. Both sources of correlation are reduced, if not eliminated, through our use of the interactive effects model. Should, however, any serial or spatial correlation remain in the residuals, it is accounted for, following \citet{bertrand+etal:04}, by using the block-bootstrap on state border segments to compute standard errors. Specifically, we draw our bootstrap samples from the set of border segments (e.g. the NJ-PA border segment includes all county pairs that lie on the NJ-PA border), not from the set of border county pairs. Moreover, each border segment in the bootstrap sample is drawn for the entire sample period. This preserves any potential spatial correlation across the different county pairs that lie along the same state border as well as the serial correlation. Note further that this is the same level at which treatment takes place, i.e. all pairs along a border segment will have the same value for the difference in log benefits, since they all lie along the same border. Thus, by using a block bootstrap on state border segments instead of a more parametric approach, we allow for an arbitrary form of serial correlation as well as spatial autocorrelation between the county pairs along the same border. The Monte Carlo experiments in Appendix \ref{app:MC_corr} illustrate that serial and spatial correlation as observed in our data does not induce a bias on the estimates when the interactive effects estimator is used.



\subsubsection{Counterfactuals}
While our empirical approach allows identifying the effect of a one-time change in benefits in the data, we can use our theoretical assumptions to infer the effects of persistent and permanent benefit changes. 

Dividing the coefficient $\alpha$ by the measurable factor $(1-\beta (1-s))$ yields the permanent percentage change of $u$ in response to a permanent one percentage change in the policy variable $b$, $ -  \lambda_J \gamma_b$. More generally, the effect of increasing benefit duration from $\omega_1$ to $\omega_2$ weeks for $n$ time periods is given by
\beq
\alpha \times\frac{1-(\beta(1-s))^{n}}{1-\beta(1-s)}\times\left(\log(\omega_2)-\log(\omega_1)\right). \label{eq:partial_sum}
\eeq

\subsubsection{Validating the Empirical Methodology using Model-Generated Data}

To ascertain the accuracy of our empirical methodology, in Appendix \ref{sec:simulation} we compare the predicted permanent effect estimated using the quasi-difference estimator to the actual permanent effect in a calibrated extension of the \citet{mortensen+pissarides:94} model to allow for unemployment benefit expiration as in \citet{mitman+rabinovich:15,mitman+rabinovich:24}.  We find that our empirical methodology is very accurate in model-generated data.
 


%%%%%%%%%%%%%%%%%%%%%%%%%%%%%
%---------------------------%
%---------------------------%
%%%%%%%%%%%%%%%%%%%%%%%%%%%%%
\section{Data}
\label{sec:data}

The paper relies on numerous sources of data that are described when they are used. In this Section we only mention the data sets that play the most significant role in the analysis.

Data on unemployment among the residents in each county\footnote{A Map of U.S. state and county borders can be found in Appendix Figure \ref{fig:Counties_Map}.}  are from the Local Area Unemployment Statistics (LAUS) provided by the Bureau of Labor Statistics.\footnote{ ftp://ftp.bls.gov/pub/time.series/la/} We supplement these data with the corresponding local-level administrative data on unemployment insurance claims and final payments from the state unemployment insurance system. We provide an extensive discussion of the quality and appropriateness of LAUS data for our estimates in Appendix \ref{app:data_quality}.\footnote{Specifically, we show that our results are robust to removing state additivity factors. Further, our substantive conclusions are unchanged when we perform an analysis using administrative data on UI claims. In Appendix \ref{app:endogeneity}, we perform an alternative test of LAUS data quality as proposed in \citet{hall:13} and find that it also supports the appropriateness of the LAUS data for our analysis.}

We use county-level employment (the number of jobs located in a county) data from the Quarterly Census of Employment and Wages (QCEW) provided by the BLS.
County-level data on wages and sectoral private sector employment (the number of jobs located in a county) are from the Quarterly Workforce Indicators (QWI).\footnote{http://lehd.ces.census.gov/datatools/qwiapp.html} QWI is derived from the Local Employment Dynamics, which is a partnership between state labor market information agencies and the Census Bureau. QWI supplies data for all counties except those in Massachusetts. Data availability varies substantially across states until 2004 Q4. Thus, for our main empirical analysis we will restrict attention to quarters beginning with 2005 Q1. To protect confidentially, noise is infused into QWI estimates resulting in some instances in significantly distorted data.\footnote{See ftp://ftp2.census.gov/ces/tp/tp-2006-02.pdf.} We exclude any observations that are flagged as significantly distorted from the analysis.


We obtain county-level vacancy data from the Help Wanted OnLine (HWOL) dataset provided by The Conference Board (TCB). This dataset is a monthly series that covers the universe of unique vacancies advertised on around 16,000 online job boards and online newspaper editions (with duplicate ads identified and removed by TCB). The HWOL database started in May 2005 and replaced the Help-Wanted Advertising Index of print advertising also collected by TCB.\footnote{For detailed information on survey methodology, coverage, and concepts see the Technical Notes at http://www.conference-board.org/data/helpwantedonline.cfm.} For a more detailed description of the data, some of the measurement issues, and a comparison with the well-known Job Openings and Labor Turnover Survey (JOLTS) data,\footnote{http://www.bls.gov/jlt/} see \citet{sahin+etal:2012}. Importantly, our analysis is based only on approximately 93\% of all online vacancies that are uniquely matched by TCB to a county of prospective employment. In other words, we do not use approximately 5\% of HWOL vacancies that are coded as ``statewide'' and 2\% that are coded as ``nationwide.''



To identify the role of unemployment benefit extensions on labor market outcomes, we focus our analysis on a sample of county pairs that are in different states and share a border.\footnote{Data on county pairs are from \citet{dube+etal:10}.} There are 1,172 such pairs for which we have complete data.

Data on unemployment benefit durations in each state is based on trigger reports provided by the Department of Labor. These reports contain detailed information for each of the states regarding the eligibility and adoption of the two unemployment insurance programs over our primary sample period: Extended Benefits program (EB) and Emergency Unemployment Compensation (EUC08).\footnote{See http://ows.doleta.gov/unemploy/trigger/ for trigger reports on the EB program and \\ http://ows.doleta.gov/unemploy/euc\_trigger/ for reports on the EUC08 program.}

The EB program allows for 13 or 20 weeks of extra benefits in states with elevated unemployment rates. The EB program is a joint state and federal program.  The federal government pays for half of the cost, and determines a set of ``triggers" related to the insured and total unemployment state rates that the states can adopt to qualify for extended benefits.   At the onset of the recession, many states chose to opt out of the program or only adopt high triggers. The American Recovery and Reinvestment Act of February 2009 turned this into a federally funded program. Following this, many states joined the program and several states adopted lower triggers to qualify for the program.\footnote{Outside of the extensions induced by the Great Recession, the EB program was triggered on in Louisiana in the aftermath of hurricane Katrina. It provided 13 extra weeks of benefits to those whose regular 26 weeks of benefits ran out between October 30, 2005 and February 26, 2006. Excluding this extension form our analysis leaves our conclusions unaffected.}


The EUC08 program enacted in June 2008, on the other hand, has been a federal program since its onset. The program started by allowing for an extra 13 weeks of benefits to all states and was gradually expanded to have 4 tiers, providing potentially 53 weeks of federally financed additional benefits. The availability of each tier is dependent on state unemployment rates. The trigger reports contain the specifics of when each state was eligible and activated the EB program and different tiers of the EUC08 program. We have constructed the data through December 2012.

Prior to the end of 2012, the duration of benefits in a given state varied over time but the level of benefits remained constant. In 2013, however, just before the  expiration of EB and EUC08 programs in December of that year, the system went through considerable upheaval due to the sequestration of the federal budget. The sequester mandated a 10.7\% reduction on spending on EB and EUC08 benefits. However, the implementation of these cuts varied widely across states with some implementing across the board cuts, others implementing much larger cuts for new entrants into the programs (and in some cases for individuals starting new EUC tiers), yet others reducing the number of weeks of benefits available under various EUC tiers but leaving the benefit levels unchanged. While these 2013 changes provide a source of significant variation that might help identify the labor market effect of unemployment benefit levels, we do not attempt to exploit it in this paper given our focus on the effects of benefit duration.

There is a substantial heterogeneity in the actual unemployment benefit durations across time and across the U.S. states. Appendix Figure \ref{fig:Benefit_Maps} presents some snapshots that illustrate the extent of this variation. Among 1,172 border county pairs used in our analysis,  1,132 have different benefits for at least one quarter. The median county pair has different benefit durations for 14 quarters during 2008-2012. The difference in available benefit duration within a county-pair ranges from 0 to 18 quarters.

Some of the data series used in the analysis are available at a monthly frequency while others are quarterly. Therefore, we aggregate all monthly data to obtain quarterly frequency.  Logs are taken after aggregation. When constructing the quasi-differences, we use the separation rate directly from JOLTS data.\footnote{In Section \ref{sec:alt_sep}, we consider the sensitivity of our estimates to constructing the quasi-differences using county measures of separations from QWI.}  For the quarterly discount factor we set $\beta = 0.99$ typically used in the studies of the aggregate economy, such as ours, and corresponding to a $4\%$ annual discount rate. We verify, however, that results are not sensitive to setting $\beta = 0.9975$ or $\beta = 0.98$, corresponding to a 1\% or 8\% annual discount rate, respectively.

%%%%%%%%%%%%%%%%%%%%%%%%%%%%%
%---------------------------%
%---------------------------%
%%%%%%%%%%%%%%%%%%%%%%%%%%%%%
\section{UI Benefit Extensions and Unemployment}
\label{sec:benefits_and_unemployment}

\begin{table}[t]
  \centering
   \caption{Unemployment Benefit Extensions and Unemployment}\label{tab:Benefits_on_unemp}
\footnotesize
\input{Benefits_on_unemp.tex}
\parbox{16.5cm}{\noindent Note - $p$-values (in parentheses) calculated via bootstrap. Bold indicates $p<0.01$.\\
Column (1) - Baseline sample,\\
Column (2) - Baseline sample controlling for State GDP per worker,\\
Column (3) - Scrambled border county pairs sample,\\
Column (4) - Scrambled border county pairs sample controlling for State GDP per worker,\\
Column (5) - Sample of border counties with  $<15\%$ share of state's employment,\\
Column (6) - Sample of border counties with similar industrial composition,\\
Column (7) - Sample of border counties with population centers $<30$ miles apart,\\
Column (8) - Sample of border counties within the same Core Based Statistical Areas,\\
Column (9) - Baseline sample with perfect foresight measure of available benefits,\\
Column (10) - Baseline sample with controls for all other state-level policies, \\
Column (11) - Sample of border counties with unemployment correlation $>0.5$ from 1996-2000.
}
\vspace{-0.0cm}
\end{table}




%%%%%%%%%%%%%%%%%%%%%%%%%%%%%%%%%%%%%%%%%%%%%%%%%%%%%%%%%%%%%%%%%%%%%%%%%%%%%%%%%%
%--------------------------------------------------------------------------------%
%%%%%%%%%%%%%%%%%%%%%%%%%%%%%%%%%%%%%%%%%%%%%%%%%%%%%%%%%%%%%%%%%%%%%%%%%%%%%%%%%%
\subsection{Baseline Empirical Results}
\label{sec:baseline_results}
Column (1) of Table \ref{tab:Benefits_on_unemp} contains the results of the estimation of the effect of unemployment benefit duration on unemployment using the baseline specification in Eq. (\ref{eq:baseline_spec}). We find that changes in unemployment benefits have large and statistically significant short-run effect on unemployment: a $1\%$ rise in benefit duration for only one quarter increases unemployment rate by $0.053$ log points.

Eq. (\ref{eq:partial_sum}) helps us extrapolate these effects to assess various policy scenarios of interest. For example, assuming perfect foresight of future benefits, our estimate of $0.053$ implies that if standard durations of unemployment benefits (26 weeks) had prevailed following the Great Recession, the unemployment rate in 2010 and 2011 would have been $3.02$ and $2.15$ percentage points lower, respectively. Although the duration of benefits was fairly similar in 2010 and 2011, the implied effect on unemployment was substantially larger in 2010. Assuming perfect foresight, this is because firms creating jobs in 2011 expected benefit duration to decline in the following year (lowering the expected wage bill and making it easier to cover the costs of vacancy creation). In contrast, firms considering creating jobs in 2010 were confronted with an additional year, 2011, of high benefit durations, which discouraged vacancy creation. This example illustrates an important general property of the effects of unemployment benefits on unemployment. To the extent that employers anticipate future changes in benefits, unemployment evolves smoothly over time and fully responds to future changes in benefits before those changes actually occur. Thus, unless expectations are known or controlled for, the relationship between contemporaneous changes in benefits and unemployment is uninformative of the true labor market impact of unemployment benefit policies.

When comparing the magnitude of the effect we find to the experience in the data, it is also important to keep in mind that it is based on the difference across pairs of border counties. Thus, the effects of various other shocks or policies that affect these counties symmetrically are differenced out. For example, the 2\% reduction to an employee's share of Social Security payroll taxes implemented in all states in 2011 and 2012 might have had a substantial negative impact on unemployment, counteracting some of the effects of unemployment benefit extensions.

Another counterfactual policy experiment of interest involves a permanent unanticipated increase in benefit durations. Using the average quarterly separation rate of $10\%$ in JOLTS data, our estimate implies a sizable impact of permanently ($n=\infty$) increasing benefits from $\omega_1=26$ to $\omega_2=99$: the long-run average unemployment rate would rise from $5\%$ to $9.6\%$. Clearly, a permanent increase in benefit durations would have the strongest negative impact on labor demand as firms would expect to bargain with workers entitled to high benefits at all future dates. It differs from the experience during the Great Recession because the enacted benefit extensions were not designed to be permanent and benefit durations varied over time.\footnote{Different choices for the discount factor affect both the estimated coefficient on benefits and the mapping between those coefficients and permanent effects. The effect is small, however. For example, re-estimating the baseline specification in Eq. (\ref{eq:baseline_spec}) under fairly extreme assumptions of $\beta=0.9975$ or $\beta=0.98$, corresponding to annual discount rates of 1\% and 8\% percent, and computing the permanent effect $\hat{\alpha}/(1-\beta(1-s))$ in Eq. (\ref{eq:partial_sum}) using these values, yields 0.484 and 0.491, respectively. These effects are very close to the 0.488 estimate obtained under the standard assumption of $\beta=0.99$, corresponding to an annual discount rate of 4\%.}


%%%%%%%%%%%%%%%%%%%%%%%%%%%%%
%---------------------------%
%%%%%%%%%%%%%%%%%%%%%%%%%%%%%
\subsubsection{Under the Hood of the Quasi-Difference}
\label{sec:decomposing_qd}

\begin{figure}%[!Ht]%
\centering
\subfigure[Current and Lagged Benefits and Unemployment]{\includegraphics[trim = 0.1in 0.1in 0.1in 0.1in, clip = true, width=0.48\columnwidth]{DiffWks.pdf}
\label{fig:binscatter_du}}
\subfigure[Current Benefits and QD Unemployment]{\includegraphics[trim = 0.1in 0.1in 0.1in 0.1in, clip = true, width=0.48\columnwidth]{QDWks.pdf}
\label{fig:binscatter_qdu}}
\vspace{-0.2cm}
\caption{(Quasi-Differenced) Unemployment Differences between Border Counties versus \newline Differences in Benefit Duration.}
\vspace{-0.3cm}
\end{figure}

Our methodology calls for estimating the effect of benefit duration on the quasi-differenced unemployment, i.e., on $\tilde{u}_t := \log(u_t) - \beta (1-s_t) \log(u_{t+1})$. It is instructive to consider the effects of benefit duration on the two components of the quasi-difference, i.e., on $\log(u_t)$ and $\log(u_{t+1})$. To this end, we regress the difference of these variables between border counties on the difference in the log benefit duration between these counties using the same interactive effects model as in the estimation of the effect on quasi-differenced unemployment. We obtain coefficients of $0.418$ ($p$-value 0.000) and  $0.406$ ($p$-value 0.000), respectively. Thus, higher benefit duration is associated with higher current and higher future unemployment, so that both components of the quasi-difference increase in benefit generosity. Note that the effect of benefits on $u_t$ or $u_{t+1}$ is substantially larger than the effect on $\tilde{u}_t$. The reason is clear. The regression of $\tilde{u}_t$ on benefits measures the effect of higher benefits for one period only. The regression of $u_t$ or $u_{t+1}$ on benefits captures, in addition, the effect of the expected duration of the high benefit policy.

To help visualize these results, in Figure \ref{fig:binscatter_du} we provide a binned scatter plot of the difference between border counties in the current or next quarter's unemployment on the difference in current benefits. In Figure \ref{fig:binscatter_qdu} we provide a similar plot for the difference between border counties in quasi-differenced unemployment. These figures suggest an unambiguous increase in current, future, and quasi-differenced unemployment with benefit duration.

%%%%%%%%%%%%%%%%%%%%%%%%%%%%%
%---------------------------%
%%%%%%%%%%%%%%%%%%%%%%%%%%%%%
\subsection{The Impact of Expected Future Policy Changes on Current Unemployment}
\label{sec:forward_spec}

Our key methodological contribution is to develop an estimator that controls for the effects of future policies on current decisions, resulting in a specification that includes only current period benefits in the regression. We now extend this methodology to assess more directly the effects of future benefit durations on current unemployment. Recall that the value of a filled job in period $t$ is given by $J_t = \pi_t+ \sum_{l=1}^{k-1} (\prod_{m=1}^l \beta (1-s_{t+m-1}))  \pi_{t+l}+   (\prod_{m=1}^k \beta (1-s_{t+m-1})) J_{t+k}$, which yields the $k$-period ahead quasi-difference value of a job $J$,
\begin{equation}
J_t - (\prod_{m=1}^k \beta (1-s_{t+m-1})) J_{t+k} = \pi_t + \sum_{l=1}^{k-1} (\prod_{m=1}^l \beta (1-s_{t+m-1}))  \pi_{t+l}.
\end{equation}
Using previous derivations that have established that 
$J -\beta (1-s) J'$, $u$, $\theta$ and $\pi$ are proportional, yields to the $k$-period ahead quasi-differenced unemployment rate, which equals the discounted sum of future one-period quasi-differenced unemployment rates:
\begin{equation}
\tilde{u}^k_{t} := \log(u_t) - (\prod_{m=1}^k \beta (1-s_{t+m-1}))  \log(u_{t+k}) = \tilde{u}_t+\sum_{l=1}^{k-1} (\prod_{m=1}^l \beta (1-s_{t+m-1}))  \tilde{u}_{t+l}
\end{equation}




This $k$-period ahead quasi-difference equals the sum of one-period quasi-differences and thus depends not only on current benefits $b_t$ but on the full future sequence of benefits until period $t+k-1$. This implies that the effects of current and future benefit policies can be directly estimated through the regression
\begin{equation}\label{eq:basic_spec_forward}
\Delta \tilde{u}^k_{p,t} = \sum_{m=1}^k \alpha_m \Delta b_{p,t+m-1} +\lambda_{p}^{k\, '} F^{k}_{t}+ \nu^{k}_{p,t},
\end{equation}
where $\nu^{k}_{p,t}$ includes expectation errors for future benefit levels in periods $t+1$ through period $t+k-1$ because firms are assumed to have rational expectations but not perfect foresight of future policies.

Our benchmark specification is a special case for $k=1$. The expanded specification in (\ref{eq:basic_spec_forward}) allows us to investigate in the data the expectational channel of policy as it includes the future benefits directly in the regression. Subject to the available panel duration, this specification can be estimated for an arbitrary $k$, that is it allows to assess the impact of benefits at  an arbitrary future date on current unemployment. Only benefits after period $t+k$ are not included in the regression but instead are controlled for by using the $k$-period ahead quasi-difference. 

Similar to the benchmark, the extended specification also allows to compute the effect of increasing benefit duration from $\omega_1$ to $\omega_2$ weeks for $n = k$ time periods (using $s_t=s$) as
\beq
\sum_{m=1}^k \alpha_m \times\left(\log(\omega_2)-\log(\omega_1)\right), \label{eq:partial_sum_forward}
\eeq
and for a permanent change in benefit duration from $\omega_1$ to $\omega_2$ weeks as
\beq
\frac{\sum_{m=1}^k \alpha_m }{1- (\beta(1-s))^k}\times\left(\log(\omega_2)-\log(\omega_1)\right). \label{eq:partial_sum_forward_perm}
\eeq

Note that the theory implies that the effect of a permanent increase in benefits calculated in (\ref{eq:partial_sum_forward_perm}) using the $\hat{\alpha}_m$ estimates from (\ref{eq:basic_spec_forward}) should be the same for all $k \geq 1$. The results of performing this experiment in the data are presented in Table \ref{tab:Forward_Spec}. We find that the estimated effects of benefit extensions are indeed quite stable across specifications which vary the number of leads in benefits. Furthermore, we test the null hypothesis that the permanent effect from our baseline estimate ($k=1$) is equal to each estimate from $k=2,\cdots,8$. The $p$-values for all seven tests are $>0.05$, implying that we cannot reject the null hypothesis that all of the coefficients are equal to the baseline at the 95\% confidence level. This result is noteworthy for two reasons. First, it confirms the appropriateness of our benchmark specification that includes only current benefits. This indicates that a one-period-ahead quasi-difference indeed fully captures the expectations of future policies. Second, it directly shows that current unemployment responds to expected changes in benefits in future periods in a way consistent with the basic theory underlying our empirical methodology.\footnote{It also corroborates and extends our finding in Section \ref{sec:decomposing_qd} that the data reject the hypothesis that quasi-differenced unemployment increases because benefits decrease future unemployment. If this were the case, then, since the future period $t+k$ unemployment rate enters the $k$-period ahead quasi-difference discounted by $\prod_{m=1}^k \beta (1-s_{t+m-1})$, the effect of a permanent increase in benefits would decrease the higher is $k$. Our findings in this section again reject this hypothesis.
}



\begin{table}[t]
\centering
\caption{Estimated Permanent Effects Using $k$-Period ahead Quasi-Difference}\label{tab:Forward_Spec}
\input{Forward_Spec.tex}
\parbox{16.5cm}{\noindent \footnotesize Note - Column (3) contains estimated effect of a permanent increase in benefit duration from 26 to 99 weeks calculated via (\ref{eq:partial_sum_forward_perm}) using $\hat{\alpha}_m$ estimates from specification in (\ref{eq:basic_spec_forward}) for various values of $k$ in Column (1). Column (4) contains the implied unemployment rate assuming a 5\% equilibrium unemployment when 26 weeks of benefits are available. For example, the entry in Column (4) for $k=1$ is obtained from $\log(0.05)+0.65=\log(9.6)$. The calculation for other values of $k$ is analogous.}
\vspace{-0.2cm}
\end{table}

%%%%%%%%%%%%%%%%%%%%%%%%%%%%%%%%%%%
%---------------------------------%
%%%%%%%%%%%%%%%%%%%%%%%%%%%%%%%%%%%
\subsection{Placebo Test for the Quasi-Difference Estimator}
\label{sec:placebo_test}
In the previous section we verified the performance of the quasi-difference estimator in the data and established the importance of expectations of future policy changes by analyzing the results from the $k$-period ahead quasi-difference estimator. We now perform another direct test that verifies the empirical performance of the estimator (we will also verify its performance in the model generated data below).
Specifically, we apply the estimator to the data from a time period when there were no benefit extensions with an artificially created placebo measure of weeks of benefits available based on a hypothetical trigger of benefit extensions.

Accordingly, we consider data from 1996-2000 when no extended benefits were available in the US.\footnote{Except for a brief extension of benefit duration in New Jersey studied by \citet{card+levine:2000} and discussed in Appendix \ref{sec:literature}. Eliminating 1996 New Jersey data from the placebo sample has no impact on the findings.}
In practice, a state triggers on a benefit extension in a given month if the three month average of the state seasonally adjusted unemployment rate exceeds a pre-determined threshold. Consequently, we obtain data on the monthly seasonally adjusted unemployment rate at the state level and specify our placebo extension of 13 weeks in any month when the preceding three month average of the state seasonally adjusted unemployment rate exceeded 6\%. We then aggregate from monthly to quarterly data, take logs and difference across border county pairs replicating exactly the procedures in our analysis of the actual UI extensions. Finally, we estimate the main specification of the paper using the interactive effects estimator and compute standard errors via block-bootstrap.

Applied to these data, our quasi-differenced estimator correctly recovers a negligible and statistically insignificant coefficient of $0.008$ (p-value of $0.35$) measuring the impact of placebo benefits on unemployment.\footnote{It is also interesting to note that extended benefits where available to those unemployed in 2001 through the extensions triggered in early 2002. The Temporary Extended Unemployment Compensation act passed in March 2002 provided extended benefits to anyone who became unemployed after March 2001. In addition, several states triggered benefit extensions in the first half of 2002 via the Extended Benefits (EB) program, which also applied to people who became unemployed in 2001. To the extent that these extensions were anticipated (e.g. because of the collapse in asset prices in the technology sector, the September 11 terrorist attacks, etc.), 2001 data is not appropriate for a placebo experiment. Nevertheless, implementing the experiment on 1996-2001 data, we find a small but now marginally statistically significant positive estimate of $0.015$ on unemployment. This increase in the coefficient is expected given our theory and reinforces our other findings on the role of expectations. Its specific magnitude is not readily interpretable, however, because we do not know the correlation between the benefit durations assigned through an artificial placebo trigger with the true anticipated durations.}




%%%%%%%%%%%%%%%%%%%%%%%%%%%%%
%---------------------------%
%%%%%%%%%%%%%%%%%%%%%%%%%%%%%
\subsection{Testing for Endogeneity}
\label{sec:endogeneity_test}

In this section we develop and implement a test to detect the presence of the endogeneity problem. We begin, however, by outlining the origin of the problem informally using an intuitive example that imposes stronger conditions than those actually required for identification.\footnote{In Appendix \ref{app:endogeneity}, we discuss endogenity testing further and provide additional examples and empirical results.}

Imagine a border county pair consisting of county $a$ belonging to state $A$ and county $b$ belonging to state $B$. State $A$ also has some geographic area $\mathcal{A}$ that excludes county $a$. We now consider two cases.

\vspace{0.3cm} \noindent {\bf Case 1. Continuous economic conditions at the state border.}\\
Suppose there is a large shock affecting the economy of $\mathcal{A}$. The economic effects of this shock might spread geographically to reach county $a$. However, there is no particular reason for these effects to stop upon reaching the state border. Thus, they will continue spreading and would affect county $b$ similarly to their effect on county $a$. If this is the case, there is no endogeneity problem in our baseline specification (\ref{eq:baseline_spec}) as the difference in unemployment between counties $a$ and $b$ is due solely to the difference in benefit policies, perhaps triggered by the developments in $\mathcal{A}$. With geographically continuous economic fundamentals, shocks directly to counties $a$ and $b$ also do not create an endogeneity problem even if either one or both counties are large enough to trigger a change in policies in the corresponding states.


\vspace{0.3cm} \noindent {\bf Case 2. Discontinuous economic conditions at the state border.}\\ The endogeneity problem can arise only if shocks to e.g., productivity, stop when reaching a state border. In this case, a shock to $\mathcal{A}$ may affect, say, productivity in county $a$ and trigger a change in unemployment benefit policy in state $A$. In contrast, this shock stops when reaching the state border so that neither $b$'s productivity nor $B$'s benefit policy is affected. In this case, the difference in unemployment between counties $a$ and $b$ is driven by both the difference in productivities and the difference in benefits, with the latter at least partially induced by the difference in productivities. In this case, the estimate of the effect of benefits would be biased if the difference in state productivities is not controlled for.\footnote{Note that an endogeneity problem would not arise even in Case 2 if there are discontinuous idiosyncratic shocks to counties $a$ or $b$ as long as these shocks do not affect the state average conditions and do not trigger changes in benefit policy at the state level. This is not a very strong restriction as the median border county has only one half of one percent of its state's employment. In addition, in Column (5) of Table \ref{tab:Benefits_on_unemp} we redo the analysis where we drop counties that have greater than 15\% share of state employment and find our estimates unchanged.}

\vspace{0.4cm} We now turn to a more formal exposition. The identifying assumption of our empirical strategy is that the error term $\nu_{p,t}$ in estimation Eq. (\ref{eq:baseline_spec}) is uncorrelated with benefits $\Delta b_{p,t}$. Unemployment at the county level is driven by benefits $b$, the time varying factors $F$ and county-specific factors such as productivity or demand which are unobserved and are part of the term $\nu_{p,t}$. The assumption that $\nu_{p,t}$ is not correlated with benefits then means that the differences in productivity, demand, etc. across border counties are not correlated with the benefits across the same counties. Since benefits are a function of state level variables, a sufficient condition for this assumption to be valid is that the difference in county level productivity, demand, etc. is uncorrelated with the corresponding differences at the state level,
\begin{equation} \label{assu:identify}
Corr (\nu_{p,t}, \Delta Z_{p,t}) = 0,
\end{equation}
where $Z_{t}$ is state level measure of productivity or demand and $\Delta Z_{p,t}$ is the difference in this measure across states. 

To test whether this condition is satisfied in the data, we can decompose the term $\nu_{p,t}$ into a part that depends on the state, $\Delta Z_{p,t}$, and another part that depends on various county-specific factors only, $\tilde{\nu}_{p,t}$,
\beq
\nu_{p,t} = \chi \Delta Z_{p,t} + \tilde{\nu}_{p,t},
\eeq
so that we rewrite the empirical specification as
\begin{equation} \label{eq:endogenous}
\Delta \tilde{u}_{p,t} =  \alpha \Delta b_{p,t} +\lambda_{p}'F_{t}+ \chi \Delta Z_{p,t} + \tilde{\nu}_{p,t}
\end{equation}
and test whether coefficient $\chi=0$.

The economics behind this test is clear. Unemployment benefit extensions are determined at the state level and thus depend on a state's economic conditions such as state-level productivity or demand $Z_t$. Thus, a negative state-level shock to $Z_t$ can cause unemployment to increase in all the counties in the state and simultaneously lead to an extension of benefits. If state-level shocks do not affect border counties similarly, i.e., $\chi \not= 0$, the estimated coefficient $\alpha$ would be biased in our baseline specification in Eq. (\ref{eq:baseline_spec}).  The presence of this bias would be revealed, however, by implementing specification (\ref{eq:endogenous}). If the bias is present, we would expect the coefficient $\chi$ on $\Delta Z_{p,t}$ to be statistically different from zero and the coefficient $\alpha$ on benefit duration to change in magnitude and perhaps lose its statistical significance.


To implement this test in the data, we use two measures of state-level conditions $Z_t$ -- state productivity (discussed in the main text) and state-level unemployment instrumented with Bartik shocks (discussed in Appendix \ref{app:endogeneity_test_unemp}). We measure state productivity as real gross state product per worker.  We obtain data on state real GDP at a quarterly frequency from the Regional Economic Accounts at the Bureau of Economic Analysis.\footnote{https://www.bea.gov/newsreleases/regional/gdp\_state/2015/xls/qgsp0915\_real.xlsx} We then divide quarterly state GDP by quarterly state employment. The results are provided in Column (2) of Table \ref{tab:Benefits_on_unemp}. Note that including the difference in state productivity has almost no effect on the estimate of the effect of benefit duration on unemployment. Including Bartik-instrumented state unemployment also has almost no effect on the estimate of the effect of benefit duration on unemployment as reported in Appendix \ref{app:endogeneity_test_unemp}. These results provide clear evidence that our findings are not driven by a mechanical relationship between economic conditions at the state level and the duration of unemployment benefits.\footnote{Columns (1) and (2) of the Appendix Table \ref{tab:Benefits_on_unemp_OLS} lead to the same conclusion based on the results of this test implemented in the additive fixed effects specification instead of the interactive effects one used here.}

The fact that the endogeneity tests indicate that economic fundamentals evolve smoothly across state borders implies that the size of the border counties relative to their state is not a relevant consideration for our analysis. Even if one of the counties is large enough so that a shock to that county triggers a policy change at the state level, both counties in the pair are affected similarly by the shock so that the difference between them depends only on the difference in benefit policies. We can also verify this implication directly. To do so, we restrict the sample to border county pairs such that each county in the pair accounts for no more than 15\% of employment in the state that county belongs to. The results of re-estimating the benchmark specification on this sample are reported in Column (5) of Table  \ref{tab:Benefits_on_unemp}. They confirm that the estimate of the effect of benefit extensions is virtually unaffected.



%%%%%%%%%%%%%%%%%%%%%%%%%%%%%%%%%%
%--------------------------------%
%%%%%%%%%%%%%%%%%%%%%%%%%%%%%%%%%%
\subsubsection{Scrambled Border County Pairs}

In the previous section, we tested for endogeneity and found that including the difference in state productivities or instrumented state unemployments has a negligible effect on the estimated effect of benefit extensions, $\alpha$, and that the effects of these state-level variables, $\chi$, are not statistically different from zero. An unanswered question, however, is whether these tests have power. To verify that they do, we generate 200 scrambled samples, where in each sample we randomly reassign our baseline set of county pairs. That is, instead of pairing neighboring counties from different states, pairs are formed by randomly matching counties from the original set of border counties. This mechanically introduces a discontinuity in economic conditions across the constructed ``border'' county pairs. Thus, a shock to productivity in state 1 of a scrambled border pair will affect productivity in county 1, but not in county 2. If that shock also moves unemployment in state 1, it would be correlated with the difference in policies between state $1$ and $2$. This invalidates our identification assumption (\ref{eq:bc_identify}).

Consequently, estimating our benchmark specification (\ref{eq:baseline_spec}) on a scrambled border county sample would yield a biased coefficient of interest $\alpha$ because $\nu_{p,t}$ is correlated with $\Delta b_{p,t}$ since both are correlated with $\Delta Z_{p,t}$. We estimate our baseline specification on the 200 scrambled samples and report the mean coefficient and $p$-value across scrambles in Column (3) of Table \ref{tab:Benefits_on_unemp} and show that the estimate of $\alpha$ is indeed substantially upward biased in the samples of randomly paired counties.

Next, we add the difference in state-level productivities to these regressions as in specification (\ref{eq:endogenous}). We expect to find a negative $\chi$ because of the endogeneity problem induced by the random pairing of counties. Adding state-level productivity, however, alleviates the endogeneity problem and diminishes the bias in estimating $\alpha$. The bias is not expected to fully disappear when we add state-level productivity since we do not control for other state variables, such as state demand, which correlate with $\nu_{p,t}$ leading to a bias. Results in Column (4) of Table \ref{tab:Benefits_on_unemp} confirm this logic. Importantly, state-level productivity is highly significant across samples of scrambled county pairs, illustrating the power of the test.

Similarly, adding the difference in state unemployment instrumented with Bartik shocks to the regression in (\ref{eq:endogenous}) estimated on the set of scrambled  ``border'' county samples reveals a highly significant coefficient of $-1.871$ ($p$-value of $0.005$) on this variable and leads to a substantial change of the coefficient on benefits from $0.106$ ($p$-value of $0$) to $-0.026$ ($p$-value of $0.49$). This is consistent with this estimation being biased (as expected) but reveals the power of the test.





%%%%%%%%%%%%%%%%%%%%%%%%%%%%%%%%%%%%%%%%%%%%%%%%%%%%%%%%%%%%%%%%%%%%%%%%%%%%%%%%%%
%--------------------------------------------------------------------------------%
%%%%%%%%%%%%%%%%%%%%%%%%%%%%%%%%%%%%%%%%%%%%%%%%%%%%%%%%%%%%%%%%%%%%%%%%%%%%%%%%%%


%%%%%%%%%%%%%%%%%%%%%%%%%%%%
%--------------------------%
%%%%%%%%%%%%%%%%%%%%%%%%%%%%
\subsection{Sensitivity of Estimates}


%%%%%%%%%%%%%%%%%%%%%%%%%%%
%-------------------------%
%%%%%%%%%%%%%%%%%%%%%%%%%%%
\subsubsection{Border Counties with Similar Industrial Composition}
As pointed out by \citet{holmes:98}, the density of manufacturing industry employment varies systematically across counties within border pairs that belong to states with different right-to-work legislation. Manufacturing industries and thus states with a large manufacturing sector have more cyclical unemployment. They may also have a more cyclical unemployment benefit policy, potentially giving rise to the endogeneity problem. If this cyclical heterogeneity across states is sufficiently empirically important, however, our interactive effects estimator picks it up through assigning a higher loading on the cyclical aggregate factor for more cyclical states.

As an additional and more general check, we now investigate whether differences in industrial composition affect our results. To this aim, we repeat the benchmark analysis on a subset of border counties with similar industrial composition.  If the industrial composition affected our results, we would expect a different result in the subsample than in the full sample. We obtain data on county employment by industry from the Bureau of Economic Analysis, Regional Economic Information System.\footnote{http://www.bea.gov/regional/}  Using sample average industry employment shares within each county, we construct the $l^2$-distance between border counties within each pair. The results, presented in Column (6) of Table \ref{tab:Benefits_on_unemp}, are based on the sample of 50\% of county pairs with the most similar industrial composition out of all border county pairs. The effect of unemployment benefit extensions on unemployment on this subsample is similar to the one found in our full sample.

As a further, more agnostic check, we restrict our sample of counties to those with high co-movement in unemployment rates before the Great Recession. In particular, we focus on the same period as our placebo sample (since benefit extensions introduce a policy discontinuity that would alter the correlation between county unemployment rates) and select only county pairs where the correlation between county unemployment rates is greater than 0.5 for the five years 1996 to 2000. We then re-run our estimator on this subsample of $698$ border pairs. The results, presented in Column (11) of Table \ref{tab:Benefits_on_unemp} show that the effect of unemployment benefit extensions on unemployment is nearly identical to our baseline sample.


%%%%%%%%%%%%%%%%%%%%%%%%%%%%%
%---------------------------%
%%%%%%%%%%%%%%%%%%%%%%%%%%%%%
\subsubsection{Degree of Economic Integration between Border Counties}

The degree of economic integration varies across county border pairs. This is relevant for the following reason. If two border counties have a fully integrated labor market with perfect mobility of workers, the residence and employment decisions are separated. In other words, the decision in which of the two counties to (look for) work is independent of the decision in which of the counties to live. Thus, in response to a change in benefits, say, in one of the states, residents of both counties adopt the same strategy of which county to work in. As unemployment is measured by the place of residence, it will be the same in both counties. Thus, our estimate of the effect of unemployment benefit extensions on unemployment would be severely biased toward zero.

Here we explore whether the potential bias is large by restricting attention to a subsets of border counties with different degrees of geographic proximity and economic integration. In Column (7) of Table \ref{tab:Benefits_on_unemp} we consider a subset of counties with population centers that are at most 30 miles apart (which is close to the median distance between population centers in our sample of border county pairs).\footnote{We are grateful to Bob Hall for sharing the geolocation data on county population centers with us.} A more extensive analysis of the effects of distance on the estimates is provided in Appendix \ref{app:benefits_and_unemployment_distance}. We find that the estimated effects of benefit extensions increase slightly with distance, further confirmation of our finding that any bias from commuting is small. In Column (8) of Table \ref{tab:Benefits_on_unemp} we further restrict the sample to include only counties with most integrated labor markets. To do so, we repeat the analysis on a restricted sample of border counties that belong to the same Core Based Statistical Areas (CBSAs). CBSAs represent a geographic entity associated with at least one core of 10,000 or more population, plus adjacent counties that have a high degree of social and economic integration with the core (see \citet{omb:2010} for detailed criteria). The results of both experiments imply similar effect of unemployment benefit extensions on unemployment to the one found in our full sample. Finally, in Appendix \ref{sec:mobility} we present evidence that workers do not change the location of employment in response to changes in benefits and that labor markets in border counties are well approximated as closed economies.


%%%%%%%%%%%%%%%%%%%%%%%%%%%%
%--------------------------%
%%%%%%%%%%%%%%%%%%%%%%%%%%%%
\subsubsection{Alternative Benefit Duration Measure}
\label{sec:perfect_foresight}

Our baseline measure of weeks of benefits available corresponds to the number of weeks a newly unemployed worker can expect to receive if current policies and aggregate conditions remained in force for the duration of the unemployment spell. An alternative, albeit extreme, assumption is that individuals have perfect foresight of the future path of benefits.

To construct the perfect foresight measure of available benefits, for a worker who becomes unemployed in a given week, we compute the realized maximum number of weeks available to him during the course of his unemployment spell (this takes into account extensions that are enacted after the spell begins).

The following example illustrates the construction of the two measures of benefit duration. Consider October 2009 in California. At the time, up to 26 regular weeks were available, and, in addition 20 weeks in Tier 1 and 13 weeks in Tier 2 of EUC08, and 20 weeks in EB.  Thus, under our baseline specification the measure of weeks available would be 26+20+13+20=79 weeks. In November of 2009, the weeks available were expanded up to 99 total (two additional tiers were added) and the program continued to be extended at those benefit levels through September of 2012.  So the perfect foresight measure would assign 99 weeks available to a worker that became unemployed in 2009.

The results based on the perfect foresight measure of available benefit duration are reported in Columns (9) of Table \ref{tab:Benefits_on_unemp}. Similar to the results based on the baseline measure of benefit availability, they continue to imply a quantitatively large impact of unemployment benefit duration on unemployment.


%%%%%%%%%%%%%%%%%%%%%%%%%%%%%%
%----------------------------%
%%%%%%%%%%%%%%%%%%%%%%%%%%%%%%
\subsubsection{Controlling for Other State-Level Policies}

Unemployment insurance policies are not the only discontinuous policy at state borders; all state policies are. Thus, an important question is whether our analysis isolates the effects of unemployment benefit policies or picks up the effects of changes in some other state policies correlated over the sample period with unemployment and unemployment benefit extensions. While this is a crucial element of the analysis, due to space constraints, a detailed analysis of the effects of individual state-level policies is provided in Appendix \ref{app:benefits_and_unemployment_state_policies}. Specifically, we collect data on and control for numerous tax, transfer, and regulatory policies at the state and county levels. We also control for the effects of stimulus spending and variations in foreclosure laws. We find that no other state policy changes were sufficiently correlated with unemployment benefit extensions to affect our estimates. In Column (10) of Table \ref{tab:Benefits_on_unemp}, we report the estimates of the benchmark specification that simultaneously controls for all other policies documented in Appendix \ref{app:benefits_and_unemployment_state_policies}. The estimated impact of unemployment benefit extensions is barely affected.

%%%%%%%%%%%%%%%%%%%%%%%%%%%%%%
%----------------------------%
%%%%%%%%%%%%%%%%%%%%%%%%%%%%%%
\subsubsection{Alternative Separation Rate Measure}
\label{sec:alt_sep}
Our baseline construction of quasi-differences uses aggregate data from JOLTS, the gold-standard measure of separations in labor economics. However, the derivation of our quasi-difference estimator in Section \ref{sec:methodology} clarifies that county separation rates may differ for exogenous and endogenous reasons. County-level separation data is available from QWI. We chose not to use these data on separations as our benchmark, as they appear noisy and often implausible (e.g., separation rates in excess of 80\%). Moreover, QWI separations are much higher than separation rates in the CPS, JOLTS, or the universe of W-2 data \citep[see, e.g., Figure E2 in][]{hyatt2020business}. Nevertheless, to assuage any concerns about this choice, we report a complete set of results based on QWI separation rates in Appendix Tables \ref{tab:Benefits_on_unemp_beg} and \ref{tab:macroeffects_beg}. All of the results are virtually identical to using the aggregate JOLTS separation rates. This is consistent with the noise in QWI being orthogonal to benefits and with our findings reported below in Section \ref{sec:vacancy} where we detect virtually no effect of benefits on unemployment operating through the endogenous separation margin.

%%%%%%%%%%%%%%%%%%%%%%%%%%%%%%%%%%%%%%
%------------------------------------%
%------------------------------------%
%%%%%%%%%%%%%%%%%%%%%%%%%%%%%%%%%%%%%%
\section{The Role of Equilibrium Macro Effects}
\label{sec:macro_effects}
In equilibrium labor market search models, the response of unemployment to changes in policies are primarily driven by employers' vacancy creation decisions. Consider, for example, an increase in unemployment benefit duration. Having access to longer spells of benefits improves the outside option of workers and leads to an increase in the equilibrium wage. This lowers the accounting profits of firms and reduces vacancy posting to restore the equilibrium relationship between the cost of firm entry and the expected profits. Lower vacancy creation leads to a decline in labor market tightness, defined as the ratio of vacancies to unemployment. This lowers the job-finding rate of workers and results in an increase in unemployment. %\textbf{KM: We've talked about this channel many times, do we need this paragraph here?}

In this section, we present evidence on the empirical relevance of these macro effects. In particular, we document the effect of unemployment benefit extensions on vacancy creation, employment, and wages in the data. We also compare the magnitude of these empirical findings to those in a calibrated equilibrium search model (with the detailed model analysis relegated to Appendix \ref{sec:simulation}).


%%%%%%%%%%%%%%%%%%%%%%%%%%%%%%%%%%%%%
%-----------------------------------%
%%%%%%%%%%%%%%%%%%%%%%%%%%%%%%%%%%%%%
\subsection{Unemployment Benefit Extensions and Vacancy Creation}
\label{sec:vacancy}
We begin by considering the effect of unemployment benefit extensions on vacancy posting by employers and on labor market tightness using HWOL data and the basic specification in Eq. (\ref{eq:baseline_spec}). The results are reported in Columns (1) and (2) of Table \ref{tab:Benefits_on_JobCreation}. We find that changes in unemployment benefits have a large and statistically significant short-run effect on vacancy creation: a $1\%$ rise in benefit duration for only one quarter lowers the number of vacancies by $0.052$ log points and labor market tightness by $0.101$ log points.\footnote{In theory there should be an exact mapping between our results in Tables \ref{tab:Benefits_on_unemp} and \ref{tab:Benefits_on_JobCreation} for (log) unemployment, vacancies and tightness, since $\log(\theta) = \log(V) - \log(U)$. The very minor discrepancy in our estimates arises because (1) we have unemployment data for more county pairs than we do for vacancy data so that the estimates are obtained on slightly different samples, and (2) benefit duration is not the only regressor as we use an interactive-effects estimator \citep{bai:09} which adds an estimated number of factors to the regression. The estimation procedure is unlikely to select exactly the same aggregate factors for the three data series (even simply including time and county-pair fixed effects would not guarantee an exact mapping of the two tables' results). The coefficients do add up exactly, of course, if we use an OLS on the same samples and with benefits as the only regressor.}


\begin{table}[t]
  \centering
   \caption{Unemployment Benefit Extensions and Job Creation} \label{tab:Benefits_on_JobCreation}
\small
\input{Benefits_on_JobCreation.tex}
\vspace{-0.2cm}
\end{table}



In the standard equilibrium search model, the matching function implies a tight relationship between changes in unemployment, vacancies, and tightness. As we have obtained independent estimates of the effects of benefit extensions on these variables, it is of interest whether their magnitudes are mutually consistent. The following calculation establishes that this is indeed the case.

Assuming that the matching function is of the commonly used Cobb-Douglas type,
\begin{equation*}
M(u,v) = \mu v^{1-\gamma} u^{\gamma},
\end{equation*}
allows us to relate the change in tightness to the change in unemployment.
Since the job-finding rate is given by
\begin{equation*}
f = \mu \theta^{1-\gamma},
\end{equation*}
the implied change in $f$ induced by a change in benefits equals $- (1-\gamma) \times 0.101$. Since the elasticity of the steady-state unemployment rate $u$ with respect to $f$ equals $1-u$, the implied change in $u$ (due to the change in tightness induced by the change in benefits) equals
\begin{equation*}
(1-u) (1-\gamma) \times 0.101.
\end{equation*}



There are many independent estimates exploiting different sources of variation of the parameter $\gamma$ in the literature. The consensus range of the estimates is between $0.4$ and $0.5$, \citet{brugemann:08}, \citet{petrongolo_pissarides:01}. Using $\gamma = 0.45$ at the midpoint of this range and assuming $u = 0.05$, the implied change equals $0.0528$. This value is very close to the actual change in unemployment reported in Table \ref{tab:Benefits_on_unemp}.

Note that this result also points toward only a small effect of benefit extensions on search effort and job acceptance decisions by the unemployed. Changes in these decisions by the unemployed would translate into changes in the parameter $\mu$ of the matching function (i.e., a decline in the job-finding rate for a given vacancy-unemployment ratio). We find very limited evidence of such changes being induced by unemployment benefit extensions, given the conventional estimates of the matching function elasticity.

The quantitative importance of the equilibrium impact of benefit extensions on job creation is further highlighted by our analysis of administrative unemployment claims data in Appendix~\ref{app:data_quality_claims}. Specifically, we find that the estimated effect of unemployment benefit extensions on the job-finding rate of unemployment benefit claimants is the same as the estimated effect of the overall job-finding rate measured above using county-level job vacancy data. This is relevant because our baseline analysis measures the effects of benefit extensions on all of the unemployed, including those who are ineligible to receive benefits. That measure, therefore, is a combination of the macro effect (which affects all unemployed) and the differential micro effect on search behavior of unemployed who are either eligible or ineligible to receive benefits. Using claims data, however, we exclusively focus on the unemployed who are eligible. If the micro effect was quantitatively important, the estimated coefficient on the claims sample should be significantly different from the one on the sample of all unemployed since even at the depths of the recession, the fraction of unemployed receiving claims did not exceed half of the unemployed. Our finding of similar effects of benefit extension on all unemployed in baseline data on all unemployed and on benefit recipients in claims data once again attributes only a very small role to the micro elasticity.

The results in this section also allows us to separate the two effects of benefit extensions on unemployment -- the one operating through vacancy posting decisions by employers and the other operating through the effect of extensions on the unemployment entry rate through endogenous job separations. In Appendix \ref{app:end_sep} we show that our benchmark estimate $\hat{\alpha}=0.053$ captures both effects. The calculation above indicates that for a conventional matching function elasticity the effects of benefits on unemployment operating through the vacancy posting margin is equal to $0.0528$. This leaves very little room for the effect of benefits on unemployment operating through the endogenous separations margin.

Finally, these findings allows us to address evidence provided by \citet{ahn+hamilton:18} that the evolving composition of the unemployment pool may affect the relationship between labor market tightness and unemployment in the time-series. These changes are not an important concern for our analysis because our estimate aggregates many ``local'' estimates, i.e., it is based on differences in benefit durations across border counties that span the entire range of benefit duration levels experienced over the Great Recession and many time periods. Thus, what we effectively estimate, is the effect for the average composition of the unemployment pool. At this level, we find that when we work directly with labor market tightness the implied results for unemployment are the same as when we work directly with unemployment.



%%%%%%%%%%%%%%%%%%%%%%%%%%%%%%
%----------------------------%
%%%%%%%%%%%%%%%%%%%%%%%%%%%%%%
\subsection{Unemployment Benefit Extensions and Employment}
\label{sec:UI_and_Emp}
In Column (3) of Table \ref{tab:Benefits_on_JobCreation} we report the effect of unemployment benefit extensions on employment using QCEW data. We find a large negative effect implying that a rise in unemployment associated with an extension of unemployment benefits is similar in magnitude to the decline in employment. This finding challenges the wisdom of relying on unemployment benefit extensions as a policy to stimulate aggregate demand. The large decline in employment associated with such policies is likely to substantially dampen any potential stimulative effects. Note that the effect of unemployment benefit extensions on employment that we find is quantitatively similar to the one measured in \citet{hagedorn+manovskii+mitman:15} who used a different source of variation for identification and relied on a different empirical methodology.


A hypothesis often mentioned in the literature is that the rise in unemployment in response to unemployment benefit extensions might be driven by measurement issues. In particular, workers who collect benefits claim to be actively searching for a job in response to surveys used to determine the unemployment rate, while in reality they are not. In other words, had benefits not been extended, these workers would have reported themselves as being out of the labor force. The decline in the vacancy rates and employment documented here provides evidence against the quantitative importance of this hypothesis.
In fact, using Eq. (\ref{eq:partial_sum}) and assuming perfect foresight of future benefits, in Section \ref{sec:baseline_results} we found that benefit extensions increased unemployment by 2.15 percentage points in 2011. The same calculation implies that employment in 2011 was lower by 2 percentage points due to benefit extensions. Thus, benefit extensions led to a decline in employment and an increase in unemployment of similar magnitude. Alternatively, we can also compute the effect  of permanently extending benefits to 99 weeks on employment:
\[
\frac{-0.0038}{1-\beta(1-s)}\times\left(\log(99)-\log(26)\right)=-0.046.
\]
Thus, the long-run average employment rate would decrease from 95\% to 90.4\%.  This 4.6 percentage point decrease is of the same magnitude as the 4.6 percentage point increase in the unemployment rate found in the corresponding experiment in Section \ref{sec:baseline_results}.


Note that our estimate of the effect of unemployment benefit extensions on employment is based on the difference across border counties. We then use the resulting coefficient to predict the effect of a nation-wide extension. A potential concern with such a procedure is that when some states extend benefits more than others, economic activity and, thus, employment may reallocate to states with lower benefits. This reallocation is picked up by our estimates but would be absent if the policy was changed nation-wide. We find no empirical justification for such a concern. In particular, we apply our empirical methodology to measure the change in employment shares of sectors producing output that is plausibly non-tradable across states, such as retail or food services. If the change in employment is driven to an important degree by reallocation, we would not expect benefit extensions to have a large effect on these sectors. Instead, we find  that a 1\% rise in benefit duration for one quarter leads to a very small decline of retail and food services employment shares by $0.003$ and $0.004$ log points, respectively. Both effects are statistically insignificant at any conventional level and the effects remain economically very small at the boundaries of the respective 95\% confidence intervals. In addition, \citet{hagedorn+handbury+manovskii:15} use the Nielsen Consumer Panel Data to measure the responsiveness of cross state border shopping to changes in unemployment benefit generosity. Their results indicate that this effect is negligible.




%%%%%%%%%%%%%%%%%%%%%%%%%%%%%%%
%-----------------------------%
%-----------------------------%
%%%%%%%%%%%%%%%%%%%%%%%%%%%%%%%
\subsection{Unemployment Benefit Extensions and Wages}
\label{sec:wages}

We have established that extensions of unemployment benefits lead to a decline in job creation by employers. In a standard equilibrium search model such a response is induced by the fact that longer expected benefit eligibility improves the outside option of workers and leads to an increase in the equilibrium wage.\footnote{\citet{pissarides:98} compares equilibrium wage effects across a wide range of models.} We now assess whether this equilibrium effect is consistent with the data.

Consider the wage of a worker $i$ in county $a$ in pair $p$ which depends on county productivity $z^a$, county market tightness $\theta^a$, benefits $b^a$ and idiosyncratic component $\phi^i$:
\beq \label{eq:wage_raw}
\log (w_t^i) = \beta_0 + \beta_z \log(z_t^a) + \beta_{\theta} \log(\theta_t^a) + \beta_b \log(b_t^a) + \log(\phi_t^i) + \eta_t^i,
\eeq
where $\eta$ is a measurement error.
Theory predicts that the equilibrium wage, conditional on county productivity, demand, etc, increases when UI becomes more generous. It is important to emphasize that we are referring to the response of the equilibrium wage, which is also negatively affected  by a drop in market tightness caused by a negative response of job creation to the policy change. The fact that the equilibrium wage combines the positive direct effect of benefit extensions and the negative effect induced by the equilibrium response of job creation, makes the identification of the net equilibrium effect on wages more demanding on the data.

Double differencing Eq. (\ref{eq:wage_raw}) across time and counties $a,b$ within a border pair yields:
\beq \label{eq:wage_thetab}
\bar{\Delta} \log(w^a) - \bar{\Delta} \log(w^b)
=  \beta_{\theta} (\bar{\Delta} \log(\theta^a) - \bar{\Delta}\log(\theta^b))
+ \beta_b (\bar{\Delta}\log(b^a) - \bar{\Delta}\log(b^b)) + \vartheta_t,
\eeq
where $\bar{\Delta}$ is the difference across time, and $\vartheta_t$ collects all error terms and stochastic components unrelated to policy. We then regress this double difference of wages on the double difference in benefits:
\beq \label{eq:wageequation_estimate}
\bar{\Delta}\log(w^a) - \bar{\Delta}\log(w^b)
=\tilde{\beta}_b (\bar{\Delta}\log(b^a) - \bar{\Delta}\log(b^b)) + \tilde{\vartheta}_t.
\eeq
In practice, to reduce the measurement error in wages, we average over two quarters, so that, e.g. $\bar{\Delta}\log(w^a_t) = (\log (w_{t+1}^a)+\log (w_{t+2}^a))/2 - (\log (w_{t}^a)+\log (w_{t-1}^a))/2$. This choice, however, does not matter for the sign or significance of the coefficient of interest.  Implementing the above regression of the double difference of the average wages of all workers on the double difference in benefits, we obtain
\beq
\tilde{\beta}_b = \beta_b + \beta_{\theta} \beta_{\theta,b} + \beta_{\phi} \beta_{\phi,b}. \notag
\eeq
As such, the estimated aggregate wage response combines the direct effect of benefits on wages, $\beta_b$, the indirect effect of benefits on market tightness, $\beta_{\theta} \beta_{\theta,b}$, and the indirect effect of benefits on the match quality distribution, $\beta_{\phi} \beta_{\phi,b}$.

Another important concern in studying the dynamics of wages is selection. First, the wages may depend on the quality of job matches the distribution of which may depend on labor market conditions \citep{gertler+trigari:09, hagedorn+manovskii:13}. We show in Appendix \ref{app:wages} how one can overcome the effect of such selection by considering wages of job stayers, whose match qualities do not change over time. Another concern is that workers in ongoing employment relationships may be partially insured against aggregate economic conditions through long-term contracts. To address this concern, in Appendix \ref{app:wages} we also consider the wages of newly hired workers who renegotiate their contracts and whose wages might be more reflective of what job creators might expect to pay when filling their vacancies. 

A dataset that that allows us to construct wages for the universe of job stayers, new hires, and all employees at the county level is QWI. We find that all three measures of wages statistically significantly increase in response to an increase in benefits. Note that this provides strong evidence for the general equilibrium effects. Indeed, if higher unemployment was not caused by unemployment benefit extensions, one would expect wages to be \emph{lower} in counties with higher unemployment.

In addition to the results based on raw wages, we also report the results based on adding the state unemployment insurance (``suta'') payroll tax to wages. The latter represents a more comprehensive measure of labor costs to employers. While unemployment benefit extensions were federally financed, states adjust suta tax rates to fund the regular benefits (typically the first 26 weeks). To the extent that federally financed benefit extensions increase the number of unemployed, this may increase the number of recipients of the regular benefits and lead to an increase in state UI taxes needed to fund these benefits. This leads to a further increase in the labor costs. Our results indicate that this is indeed the case, although the effect is quantitatively small.


The results of implementing the regression in Eq. (\ref{eq:wageequation_estimate}) are reported in Columns (1) and (2) of Table \ref{tab:Benefits_on_Wages}, for raw wages and wages adjusted for the UI payroll tax, respectively. We find that average wages statistically significantly increase in response to an increase in benefits. Note that this provides strong evidence for the general equilibrium effects. Indeed, if higher unemployment was not caused by unemployment benefit extensions, one would expect wages to be \emph{lower} in counties with higher unemployment. If wages are set by long term contracts, the coefficient $\tilde{\beta}_b$ estimated on all wages is biased downward if there is mean-reversion of state economic conditions. Despite these offsetting effects, the positive equilibrium response of wages to benefit extensions dominates. Further, in Appendix \ref{app:wages}, we find that both the wages of stayers and new hires also statistically significantly increase in response to an increase in benefits.



\begin{table}[t]
  \centering
   \caption{Unemployment Benefit Extensions and Wages} \label{tab:Benefits_on_Wages}
\small
\input{Benefits_on_Wages.tex}
\vspace{-0.2cm}
\end{table}



\paragraph{Discussion of wage effects}

The fact that benefit extensions pass through to wages is unsurprising given recent research on the extent of private and public insurance available to most unemployed. Early work by \citet{gruber1997consumption}, and subsequent work by \citet{krueger2016macroeconomics} and \cite{herkenhoff2019impact} have shown that the unemployed have very few liquid assets and can replace less than 3\% of lost income using financial assets. Moreover, \citet{braxton2024can} document very low substitutability between private and public insurance, implying that benefit extensions should have large effects on outside options.


%%%%%%%%%%%%%%%%%%%%%%%%%%%%%%%%%%%%%%%%
%--------------------------------------%
%--------------------------------------%
%%%%%%%%%%%%%%%%%%%%%%%%%%%%%%%%%%%%%%%%
\section{Conclusion}
We developed an empirical methodology to measure the effect of unemployment benefit extensions on unemployment that includes the impact of benefit extensions on job creation neglected in the existing empirical literature. In particular, we exploited the discontinuity of unemployment insurance policies at state borders to identify their impact. Our estimator controls for the effect of expectations of future changes in benefits and has a simple economic interpretation. It is also robust to the heterogeneous effects of aggregate shocks on local labor markets.  Moreover, our estimator allows us to utilize
the entire panel data on benefit duration across space and time without having to identify any unexpected policy changes.

We found that unemployment benefit extensions have a large effect on total unemployment and can potentially account for a significant share of the persistently high unemployment following the Great Recession. The analysis of administrative claims data reveals that unemployment benefit extensions have a significant negative impact on labor demand. We found further support for this conclusion by documenting direct evidence of a significant negative impact of unemployment benefit extensions on vacancy creation and employment through their effect on wages. These results indicate the importance of taking the impact of benefit extension on job creation into account for a more complete evaluation of this policy tool. They also illustrate the need for the new methodology required to measure the response of forward-looking job creators that we develop.




One motivation for increasing unemployment benefit durations during the Great Recession, in addition to helping unemployed workers smooth their consumption, was to increase employment through its stimulative effect on local demand. We cannot do full justice to evaluating this effect because our methodology is based on differencing between border counties. Nevertheless, to the extent that the unemployed spend a significant fraction of their income in their home counties (in the form of, e.g., rent payments or purchases of services), the corresponding part of the stimulative effect is fully captured by our analysis. Yet, we find that border counties with longer benefit durations have much higher unemployment despite the potential beneficial effects of spending. Future research on the potential stimulative demand effects of unemployment benefit extensions should account for the magnitude of the macro effects documented in this paper. Specifically, the stimulative impact of higher spending by the unemployed may be partially or even fully offset by the negative impact of benefit extensions on employment, which operates through reduced vacancy creation.


%%%%%%%%%%%%%%%%%%%%%%%%%%%%%%%%%%%%%%%%%
%---------------------------------------%
%---------------------------------------%
%%%%%%%%%%%%%%%%%%%%%%%%%%%%%%%%%%%%%%%%%
%\clearpage
% \setstretch{1.1}
\singlespacing
%\setstretch{0.9}
% \setlength\bibitemsep{0pt}
%\bibliographystyle{aea}
% \small
\bibliography{articles}

\pagebreak
%%%%%%%%%%%%%%%%%%%%%%%%%%%%%%%%%%%%%%%%%%%%%%%%%%%%%%%%%%%%%%%%%%%%%%%%%%%%%%%%%%%%

\begin{appendix}

\onehalfspacing
% \normalsize
%\renewcommand{\thepage}{A-\arabic{page}}
\setcounter{page}{1}

\renewcommand{\thetable}{A-\arabic{table}}
\setcounter{table}{0}

\renewcommand{\thefigure}{A-\arabic{figure}}
\setcounter{figure}{0}

\renewcommand{\theequation}{A\arabic{equation}}
\setcounter{equation}{0}

\renewcommand{\thesection}{\Roman{section}}
\setcounter{section}{0}


\begin{center}
{\bf APPENDICES FOR ONLINE PUBLICATION}
%{\bf  APPENDICES}
\end{center}



%%%%%%%%%%%%%%%%%%%%%%%%%%%%%%%%%%%%%%%%%%%%%%%%%%%%%%%%%%%%%%%%%%%%%%%%%%%%%%%%%%
%--------------------------------------------------------------------------------%
%%%%%%%%%%%%%%%%%%%%%%%%%%%%%%%%%%%%%%%%%%%%%%%%%%%%%%%%%%%%%%%%%%%%%%%%%%%%%%%%%%
\section{Comparison to Existing Empirical Estimates}
\label{sec:literature}
The existing literature measuring the impact of unemployment benefit extensions on unemployment has not yet attempted to measure the joint contribution of micro and macro effects. Nevertheless, the research design underlying some strands of the literature suggests that their estimates might be interpreted as suggestive of the magnitude of the total effect. We now review this literature.





One key strand in the relevant literature is based on the seminal contributions by \citet{moffitt:85}, \citet{katz+meyer:90}, and \citet{meyer:90}.\footnote{\citet{krueger+meyer:02} provide a survey of other important contributions to this strand of the literature.} These authors used administrative data on unemployment benefit recipients and exploited the cross-state variation in unemployment benefit extensions to measure the effect of the extensions on the hazard rate of leaving compensated unemployment. The effects found in this literature are sizable, implying that a one week increase in potential benefit duration increases the average duration of the unemployment spells of UI recipients by 0.1 to 0.2 weeks.\footnote{Many influential studies find even larger effects, e.g. \citet{ham+rea:87}.} These studies were based on a sample of unemployed workers who collect benefits. To assess the effect of benefit duration on overall unemployment one also needs to know the impact on those unemployed who do not collect benefits. This was studied by \citet{hagedorn+manovskii+mitman:15_micro}, who show that the job-finding rate of ineligible workers responds as much as that of the eligible ones to the specific benefit extensions studied in these seminal contributions as well as during the Great Recession.





Consider the implications of these estimates. During the Great Recession unemployment benefits were extended 73 weeks from 26 to 99 weeks. Thus, these estimates imply an increase in unemployment duration between  7.3 ($=73\times 0.1$) and 14.6 ($=73\times 0.2$) weeks, i.e. the duration approximately doubles. But a doubling of duration implies that the exit rate from unemployment falls by a factor of two. This would then imply roughly a doubling of the unemployment rate, as can be seen from, e.g., the basic steady state relationship that balances flows in and out of unemployment, $u=s/(s+f)$. This is a considerably larger effect than the one we find.





However, the interpretation of these large consensus estimates is not clear. The literature was exploiting the variation of benefit duration across US states during the recessions of the early 1980s to measure the effects of benefit duration on the probability to find a job. As the regressions were estimated on state-level data, to obtain an unbiased estimate of this effect, one needs to control for all other aggregate factors affecting the state-level job-finding rate. This can be accomplished by including a state-time dummy in the regression. The effect of benefits would then be identified from cross sectional variation among unemployed in time until their benefits run out. While this approach would identify the coefficient properly, the interpretation is important. The coefficient would only reflect the micro effect, that is the effect of benefits on individuals' incentives to search. The macro effect of benefits on labor demand would be instead subsumed by the state-time dummy alongside all other aggregate shocks.





Virtually all existing studies, however, include the state unemployment rate instead of a state-time dummy in the regressions of state-level job-finding rates on benefit duration. While this cannot control for all aggregate factors, the underlying assumption must be that benefits and unemployment are orthogonal to all other aggregate factors. Note that even if this assumption were correct, the coefficient on benefits would be a convolution of the micro effect and some part of unknown magnitude and sign of the macro effect. However, this assumption is clearly erroneous because state unemployment and/or benefits are endogenous. This is easy to see. Consider some aggregate shock that affects the job-finding rate in period $t-1$. By definition, this shock also affects unemployment (at the end of) period $t-1$. If the shock is persistent (and most aggregate shocks are), this shock also affects the job-finding rate in period $t$. Thus, the right-hand-side variable unemployment is correlated with the error of the regression in which the job-finding rate is the dependent variable. A biased regressor contaminates the whole regression. Similarly, past shocks to the labor market affect current benefits, rendering the estimated coefficient on benefits biased as well. Thus, the interpretation of coefficients in a regression of the state-level job-finding rate on benefits and state unemployment is very unclear, especially given that state-level benefits and unemployment are highly correlated.


Thus, at least in theory, the methodology in this literature can yield an estimate of the micro effect at the cost of ignoring the effect of benefit extensions on labor demand (by differencing it out or capturing it together with all other shocks in the state-time dummy). What the literature lacks is an identification strategy that accounts for the effect of unemployment benefit extensions on labor demand by forward looking job creators -- the macro effect. The development of the necessary methodology is a contribution of this paper.


Another strand in the literature that might be suggestive of the magnitude of the total effect of unemployment benefit extensions is based on a quasi-experimental research design. In a seminal study following this research strategy, \citet{card+levine:2000} consider the experience of New Jersey that awarded 13 extra weeks of benefits to those whose regular 26 weeks of benefits expired between June and November of 1996. This unemployment benefit extension was a ``side-effect'' of a political process and was not driven by underlying economic conditions, minimizing the need to control for other aggregate shocks and the endogeneity of policy. The authors estimate that this policy lead to a 16.6\% decline in the exit rate from unemployment, which implies a slightly smaller effect of policy than what is implied by the literature studying the effects of 1980s extensions discussed above.\footnote{Specifically, \citet{card+levine:2000} find that a one week increase in potential benefit duration increases the average duration of  unemployment spells of UI recipients by 0.08 weeks.} The authors suggest that this might be driven by the fact that their experiment took place during a relatively prosperous period while earlier studies considered the effects of benefit extensions in recessions. What might be also important, however, is that neither of the two literatures has controlled for the effects of expectations regarding future policy changes. It is plausible that job creators where correctly expecting the 1980's benefit extensions to be more persistent, explaining a quantitatively larger response. Taking this into account, the magnitude of the effect found in \citet{card+levine:2000} is much larger than what our estimates would imply for such a small and transitory extension.


The fact that our estimates imply significantly smaller effects of unemployment benefit extensions relative to most earlier studies might also be consistent with a secular decline in the responsiveness of the labor market to this policy.\footnote{\citet{katz:10} proposes that this might be due to a declining reliance by firms and industries on using temporary layoffs with recall dates tailored to benefit durations.} Indeed, we repeated the analysis in this paper using the data on benefit extensions during the earlier 1991 recession (using the 1990-1996 sample) and found a coefficient on benefit duration of $0.072$ with a $p$-value of $0$. Repeating the analysis using the data on benefit extensions during the 2001 recession (the 1996-2005 sample) yields a coefficient on benefit duration of $0.062$ with a $p$-value of $0$. While these estimates still suggest smaller effects of benefit extensions than found in the earlier literature using the 1980s data, it is larger than our estimate of $0.053$ using the later data from the Great Recession period.

A recent quasi-experimental study by \citet{johnston+mas:15} assessed the impact of the sudden and unanticipated cut in potential benefit duration by 16 weeks in Missouri in April 2011. The cut applied only to new claimants while those who claimed benefits prior to the reform were grandfathered into the old potential benefit duration schedule. Comparing the hazards out of unemployment for these two samples the authors find a large micro effect, the magnitude of which is quite surprising given the estimates in the existing literature. Their experiment is, however, consistent with the presence of a very large macro effect on job creation which is evident in Figure \ref{fig:MO_vacPuD}. Specifically, in Panel \ref{fig:MO_All_Borders} we plot the seasonally adjusted monthly difference in the log vacancy-unemployment ratio between border counties in Missouri and their paired counties from all bordering states during a year before and a year after the reform. We observe a sharp and discontinuous rise in the vacancy-unemployment ratio in Missouri at the time of the reform. As predicted by the standard equilibrium search model, the jump was caused by a sharp increase in job vacancy creation. Note that at least on impact, the jump in vacancy creation cannot be in response to higher search effort of new claimants eligible for fewer weeks of benefits following the reform simply because they represent a tiny fraction of all unemployed (only about a third of all unemployed were claiming benefits in Missouri at the time of the reform and it took a long time for the claimants under the new rules to account for a meaningful share of all claimants).


\begin{figure}[t]%[!Ht]%
\centering
\subfigure[All Missouri Borders.]
    {
        \scalebox{0.85}{\includegraphics[trim = 0.15in 0.1in 0.15in 0.1in, clip = true, width=0.52\columnwidth]{motight_pooled.pdf}}
        \label{fig:MO_All_Borders}
    }
%\hspace{-0.15cm}
%\vspace{-0.05cm}
\subfigure[Each Missouri Border.]
    {
        \scalebox{0.55}{\includegraphics{motight_all.pdf}}
        \label{fig:MO_Each_Border}
    }
\caption{The Macro Effect of the Cut in Potential Benefit Duration in Missouri in April 2011.}\label{fig:MO_vacPuD}
\vspace{0.2cm}
\footnotesize
\parbox{16.5cm}{\noindent Note - Panel \ref{fig:MO_All_Borders} plots the difference in the log vacancy-unemployment ratio between border counties in Missouri and adjacent border counties across all Missouri state borders. Panel \ref{fig:MO_Each_Border} plots the same statistic separately for each Missouri border. The dashed vertical line indicates the date when potential benefit duration was cut by 16 weeks in Missouri.}
\vspace{-0.3cm}
\end{figure}


In Panel \ref{fig:MO_All_Borders} we plot a similar figure separately for each Missouri state border and find that the sharp rise in the vacancy-unemployment ratio following the reform is apparent relative to every state that Missouri borders. To improve exposition, we normalize the log vacancy-unemployment ratio along each state border to be zero in the quarter before the reform. We provide separate unnormalized plots for each border in Figure \ref{fig:Tightness_MO_Borders} and in Table \ref{tab:MO_V_U} we decompose the jump in labor market tightness in Missouri relative to each bordering state into the contribution of the change in vacancies and in unemployment. \citet{mitman2022micro} extend this analysis and provide additional direct evidence for the macro effect on job creation following the benefit cut in Missouri using a variety of identification strategies and find increases in unemployment benefit duration significantly decrease labor market tightness. 


While the implied macro effect appears very large, interpreting its magnitude is difficult for two reasons. First, this simple experiment does not control for job creators' expectations. Indeed, this experiment is likely to yield a lower bound on the true macro effect if job creators in Missouri assigned, prior to the reform, some probability to a future cut in benefit duration but were surprised by the specific timing of when the reform was implemented. Similarly, there is a downward bias if, upon observing the policy change in Missouri, job creators in neighboring states assigned a higher probability to similar benefit duration cuts taking place in their states. Second, it is difficult to draw econometrically sound conclusions robust to the influence of sampling uncertainty from essentially one data point provided by this reform. Consequently, in Section \ref{sec:macro_effects} in the main text we apply the methodology we developed to control for the effects of expectations in measuring the macro effects and expand the scope of the analysis to include all changes in benefit duration in all states in the contiguous US following the Great Recession. This analysis yields an estimate of the macro effect that is one third smaller than what the estimate based on the Missouri reform, taken at face value, suggests.



\vspace{0.5cm}\begin{figure}[H]
\makebox[\textwidth]{\parbox{1.5\textwidth}{
\centering
	\captionsetup{justification=centering}
     \caption{The Dynamics of Vacancy-Unemployment Ratio along each Missouri State Border \\ Surrounding the Cut in Potential Benefit Duration in Missouri in April 2011.}\label{fig:Tightness_MO_Borders}
\subfigure[MO-AR Border.]
    {
        \scalebox{0.33}{\includegraphics{motight_border11.pdf}}
        \label{fig:motight_border11}
    }
\hspace{-0.15cm}
%\vspace{-0.05cm}
\subfigure[MO-IL Border.]
    {

        \scalebox{0.33}{\includegraphics{motight_border46.pdf}}
        \label{fig:motight_border46}

    }
\hspace{-0.05cm}
%\vspace{-0.25cm}
\subfigure[MO-IA Border.]
    {
        \scalebox{0.33}{\includegraphics{motight_border52.pdf}}
        \label{fig:motight_border52}
    }
\hspace{-0.05cm}
%\vspace{-0.25cm}
\subfigure[MO-KS Border.]
    {
        \scalebox{0.33}{\includegraphics{motight_border56.pdf}}
        \label{fig:motight_border56}
    }

\vspace{-0.05cm}
  \subfigure[MO-KY Border.]
    {
        \scalebox{0.33}{\includegraphics{motight_border59.pdf}}
        \label{fig:motight_border59}
    }
		%\vspace{-0.25cm}
   \subfigure[MO-NE Border.]
    {
        \scalebox{0.33}{\includegraphics{motight_border81.pdf}}
        \label{fig:motight_border81}
    }
%\vspace{-0.25cm}
  \subfigure[MO-OK Border.]
    {
        \scalebox{0.33}{\includegraphics{motight_border82.pdf}}
        \label{fig:motight_border82}
    }
%\vspace{-0.25cm}
   \subfigure[MO-TN Border.]
    {
        \scalebox{0.33}{\includegraphics{motight_border83.pdf}}
        \label{fig:motight_border83}
    }

\vspace{-0.1cm}
%\end{center}
\vspace{0.2cm}
\footnotesize
\parbox{16.5cm}{\noindent Note - Each panel plots the difference in the log vacancy-unemployment ratio between border counties in Missouri and adjacent border counties separately for each Missouri state border. The dashed vertical line indicates the date when potential benefit duration was cut by 16 weeks in Missouri.}
\vspace{0.0cm}
}}
\end{figure}
%\end{changemargin}


\begin{table}[H]
  \centering
   \caption{Changes in Vacancies, Unemployment, and their Ratio along Missouri Borders } \label{tab:MO_V_U}
\small
\input{MO_V_U.tex}
\vspace{0.2cm}
\parbox{16.5cm}{\noindent Note - The table summarizes the jump in the log vacancy-unemployment ratio following the potential benefit duration cut by 16 weeks in Missouri in April 2011 and decomposes it into the rise in vacancies and a decline in unemployment. Each row measures the difference between border counties in Missouri and adjacent border counties separately for each Missouri state border.}
\vspace{-0.3cm}
\end{table}

\FloatBarrier

\clearpage

%%%%%%%%%%%%%%%%%%%%%%%%%%%%%%%%%%%%
%----------------------------------%
%----------------------------------%
%%%%%%%%%%%%%%%%%%%%%%%%%%%%%%%%%%%%
\clearpage
\section{Endogenous Separations}
\label{app:end_sep}

In this appendix, we explain that our methodology identifies the combined effect of unemployment benefit extensions that operates through the job creation margin and through the effect on endogenous job separations.

Recall that in Section \ref{sec:methodology} we started from the firm vacancy posting condition and derived equation (\ref{eq:jc}) that market tightness  
$\log(\theta_t)$ is determined by  $J_t$, the value of the job for the firm:
\begin{equation}
\log(\theta_t)=\tilde{\kappa} \log(J_t) + \nu^a_t,
\end{equation}
for an approximation error $\nu^a_t$ and constant $\tilde{\kappa}$.
In a simple search model in which unemployment is exclusively driven by market tightness $\theta_t$, the ratio of vacancies to unemployment, \citep{pissarides:00}, in particular with exogenous separations, the log unemployment rate, $\log(u_t)$, equals (omitting approximation errors):
\beq\label{eq:fec}
\log(u_t) =  \lambda_x \log(\theta_t) = \lambda_x \frac{\lambda_J}{\tilde{\kappa}} \log(J_t),
\eeq
for a parameter $\lambda_x$. In words, these equations indicate that to the extent that firms' profits and job values are reduced by higher unemployment benefits, vacancy creation is driven down and this leads to higher unemployment ($\lambda_x$ is negative).

If job separations $s_t$ respond to benefits $b_t$, the impact of benefits on vacancy creation remains the same but now there is an additional channel through which benefits affect unemployment. We show next that this only affects the interpretation of the parameter $\lambda_x$ but leaves our measurement of the total effect of benefits unaffected.

An endogenous separation is a mutually consensual decision between the firm and the worker to end the employer-employee relationship. It is efficient for the worker and the firm to separate if the surplus of the relationship is negative. 
The firm's share of the surplus from a specific match is given by the sum of $J$ and some idiosyncratic component specific to the match drawn from a continuous distribution.
We then make the standard assumption that there is a one-to-one mapping between the firm's surplus and the surplus of the match, which is the case, for example, if both the worker and the firm benefit from a higher match surplus.
Note that without idiosyncratic components we would not have a meaningful model of endogenous separations since either no-one or everyone would separate.



The magnitude of endogenous separations is then negatively related to $J_t$ so that  we obtain a continuous relationship between $s_t$and $J_t$,
\begin{equation}\label{eq:sep}
\log(s_t)=\tilde{\kappa}_s \log(J_t).
\end{equation}
With endogenous separations, the unemployment rate changes now for two reasons, a change in $\theta$ and/or a change in $s$, using Eq. (\ref{eq:sep}):
 \beq\label{eq:decomp}
\log(u_t) = \lambda_x \log(\theta_t) + \lambda_s \log(s_t) = (\lambda_x \frac{\lambda_J}{\tilde{\kappa}} + \lambda_s \tilde{\kappa}_s ) \log(J_t),
\eeq
establishing that the unemployment rate can be written as function of $J$ as claimed in (\ref{eq:uJ}), which captures that vacancies and endogenous separations drive the unemployment rate. 
The unemployment rate thus equals
 \beq
\log(u_t) =  \lambda_x \log(\theta_t) +  \lambda_s \tilde{\kappa}_s  \log(J_t).
\eeq
Then using $\log(J_t) = \log(\theta_t) / \tilde{\kappa}$ from Eq. (\ref{eq:fec}) yields the unemployment rate as a function of market tightness,
\beq
\log(u_t) &  = &  \lambda_x \log(\theta_t) +  \frac{\lambda_s \tilde{\kappa}_s}{\tilde{\kappa}} \log(\theta_t) \\
	       & =&  \underbrace{(\lambda_x  + \lambda_s \frac{\tilde{\kappa}_s}{\tilde{\kappa}} )}_{=: \tilde{\lambda}_x }      \log(\theta_t).
\eeq
The remaining derivations in the draft go through unchanged except that $\lambda_x$ is replaced throughout by $\tilde{\lambda}_x$. Specifically, we obtain the quasi-difference
\beq
\tilde{u}_t := \log(u_t) - \beta (1-s_t) \log(u_{t+1}) = \tilde{\kappa} \tilde{\lambda}_x (1-\beta (1-s)) \log(\pi_{t}) + \tilde{\lambda}_x \log(\eta_t)
\eeq
and estimate
\begin{equation}
\Delta \tilde{u}_{p,t} = \tilde{\alpha} \Delta b_{p,t}+\Delta\epsilon_{p,t}.
\end{equation}
Thus, introducing endogenous separations does not change our baseline estimation equation which remains well identified. It only changes the structural interpretation of the estimated coefficient $\tilde{\alpha}$.


%%%%%%%%%%%%%%%%%%%%%%%%%%%%%%%%%%%%%%%%
%--------------------------------------%
%--------------------------------------%
%%%%%%%%%%%%%%%%%%%%%%%%%%%%%%%%%%%%%%%%
\clearpage
\section{Iterative Effects Estimator}
\label{app:implementation}

%%%%%%%%%%%%%%%%%%%%%%%%%%%%%%%%%%%%%%%%%%
%----------------------------------------%
%%%%%%%%%%%%%%%%%%%%%%%%%%%%%%%%%%%%%%%%%%
\subsection{Formal Identifying Assumptions}
\label{app:moon}
In this Appendix, we state the formal assumptions for identification using the interactive effects estimator in our context. We can re-write our specification \ref{eq:baseline_spec} in matrix form:

\begin{equation*}
\tilde{U} =  \alpha^0 b +\lambda^0 (f^0)'+ \nu, 
\end{equation*}
where $\tilde{U}, b, \nu \in\mathbb{R}^{N\times T}$ are $N\times T$ matrices (where $N$ is the number of pairs), $\lambda^0 \in\mathbb{R}^{N\times R^0}$ is an $N\times R^0$ matrix, $f^0 \in\mathbb{R}^{T\times R^0}$ is a $T\times R^0$ matrix, and $\alpha^0$ is a scalar. The superscript zero indicates the true value of the parameters. Let $B \in\mathbb{R}^{NT\times 1}=vec(b)$ be the $NT$ vectorization of  $b$.

Following \cite{moon2015linear}, sufficient conditions for identification of $\alpha^0$, $\lambda^0 (f^0)'$, and $R^0$, are that there exists a nonnegative integer $R$ such that:
\begin{enumerate}
\item The second moments of $b_{p,t}$ and $\nu_{p,t}$ exist for all $p,t$.
\item $\mathbb{E}(\nu_{p,t})=0$ and $\mathbb{E}(b_{p,t}\nu_{p,t})=0$ for all $p,t$.
\item $\mathbb{E}\left(B'\left(M_F\otimes M_{\lambda^0} \right)B\right)>0$, for all $F\in\mathbb{R}^{T\times \tilde{R}}$
\item $\tilde{R}\geq R^0:=\textrm{rank}(\lambda^0 (f^0)')$ 
\end{enumerate}
where $M$ is the standard annihilator matrix, e.g. $M_F=\mathbb{I}-F(F'F)^{-1}F'$.\footnote{While \cite{moon2015linear} and \cite{bai:09} study the identification of the same model, the mathematics of their proofs are different.  \cite{moon2015linear} allow for regressors to be predetermined, which is not allowed by \cite{bai:09}. On the other hand, \cite{moon2015linear} assume that errors are iid normal, while \cite{bai:09} allows them to be weakly dependent. \cite{moon2015linear} conjecture in their paper that iid assumption is not needed for their results but the math needed for the more general proofs has not been developed yet. They provide Monte Carlo simulations which suggest that the result also holds for nonnormal and correlated errors. The results of our own Monte Carlo simulations below in this Appendix support the conjecture in \cite{moon2015linear}. }

The first two assumptions are standard and equivalent to the assumptions needed for identification in OLS. Condition (1) requires finite second moments of the regressors and errors. Condition (2) requires that the error term is mean zero and the error term is uncorrelated with benefits. We provide explicit tests in the draft that assumption (2) holds in our setting by looking at the endogeneity test in the data and the $k$-period quasi-difference. Assumption (3) is a noncollinearity condition (analogous to the standard no-multicollinearity condition in OLS), which says that benefits (our regressor) have significant variation across pairs $p$ and over time $t$ after projecting out all of the variation that can be explained by the true factor loadings $\lambda^0$ and by arbitrary factors $F\in\mathbb{R}^{T\times \tilde{R}}$. The condition is analogous to OLS where one subset of regressors cannot be expressed as a linear combination of the others (in which case multiplying by the annihilator matrix would leave no variation). This condition is a generalization of the within-variation assumption for an additive fixed-effect regression with time-invariant pair fixed effects, which would read $\mathbb{E}\left(b'\left(M_{1_T}\otimes \mathbf{1}_N \right)b\right)>0.$ The conventional fixed effect assumption rules out time-invariant regressors. Similarly, assumption (3) rules out more general low-rank regressors. To check that we satisfy this assumption, we calculate the rank of $b$ in our data. We find that the $\textrm{rank}(b)=20$, whereas we estimate 2 factors in the data, much lower than the rank of 20 of $b$, implying that we do not suffer from a problem of low-rank regressors. In calculating the number of factors, we entertain up to $\tilde{R}=7$, which is still a significantly lower rank than 20 from $b$. Finally, assumption (4) says there must be a finite number $\tilde{R}$, which is greater than the true number of factors $R^0$. We estimate $R^0=2$ in the data. Thus, we conclude that all of the assumptions needed for identification are satisfied here.

%%%%%%%%%%%%%%%%%%%%%%%%%%%%%%%%%%%%%%%%%%%%%%%%%
%-----------------------------------------------%
%%%%%%%%%%%%%%%%%%%%%%%%%%%%%%%%%%%%%%%%%%%%%%%%%
\subsection{Implementation}
The following is a brief description of the algorithm implementing our iterative two-stage estimator.
\begin{enumerate}
\item Start with a guess for $\alpha$, say $\alpha_{1}$.
\item At each iteration $\xi$, do the following:

\begin{enumerate}
\item given $\alpha_{\xi}$, for each $p$, construct $\upsilon_{p,t}=\Delta x_{p,t}-\beta\left(1-s_t \right)\Delta x_{p,t+1}-\alpha_{j}\Delta b_{p,t}.$
Then, $\upsilon_{p,t}=\lambda_{p}'F_{t}$ is a pure factor model and can be estimated consistently using principal components.\footnote{The exposition of the estimator assumes that there are no missing observations. We use the generalized procedure described in \citet{bai:09} and allow for missing observations.}
\item Given the estimates for $\lambda_p$ and $F_t$, estimate Eq. (\ref{eq:baseline_spec}) via OLS and update the guess to obtain $\alpha_{\xi+1}$.
\end{enumerate}
\item Repeat 2 until $\alpha_{\xi}$ converges.%\footnote{We have conducted a number of Monte Carlo simulations with sample sizes similar to our sample. The estimator described here is found to converge to the true parameter. Results are available upon request.}
\end{enumerate}


%%%%%%%%%%%%%%%%%%%%%%%%%%%%%%%%%%%%%%%%%%%%%%%%%%%%%%%%%%%%%%%%%%%%%%%%%%%%%%%%%%
%--------------------------------------------------------------------------------%
%%%%%%%%%%%%%%%%%%%%%%%%%%%%%%%%%%%%%%%%%%%%%%%%%%%%%%%%%%%%%%%%%%%%%%%%%%%%%%%%%%
\subsection{Monte Carlo Evaluation}

%%%%%%%%%%%%%%%%%%%%%%%%%%%%%%%%%%%%%%%%%%%%%%%%%%%%%%%%%%%%%%%%%%%%%%%%%%%%%%%%%%
%--------------------------------------------------------------------------------%
%%%%%%%%%%%%%%%%%%%%%%%%%%%%%%%%%%%%%%%%%%%%%%%%%%%%%%%%%%%%%%%%%%%%%%%%%%%%%%%%%%
\subsubsection{Convergence}
In practice we start the algorithm with the initial guess $\alpha_{1}=0$. The estimator converges to the true parameter irrespective of the initial guess, however. A simple Monte Carlo study illustrates. We take 1000 random draws from the standard normal distribution, and using each draw as a starting point for the algorithm, we re-estimate the effect of unemployment benefit extensions on unemployment. Figure \ref{fig:initial-guess} shows a scatterplot of the estimates
against the initial guesses. It indicates that for each initial guess the factor model yields a coefficient virtually identical to the benchmark. The 5th percentile out of 1000 estimated coefficients is 0.05317857 and the 95th percentile is 0.05317863.

\begin{figure}[t]
\centering
\caption{Irrelevance of the initial guess for the factor model estimation}
\label{fig:initial-guess}
\includegraphics[scale=0.7]{initial_guess.pdf}
\end{figure}

%%%%%%%%%%%%%%%%%%%%%%%%%%%%%%%%%%%%%%%%%%%%%%%%%%%%%%%%%%%%%%%%%%%%%%%%%%%%%%%%%%
%--------------------------------------------------------------------------------%
%%%%%%%%%%%%%%%%%%%%%%%%%%%%%%%%%%%%%%%%%%%%%%%%%%%%%%%%%%%%%%%%%%%%%%%%%%%%%%%%%%
\subsubsection{Correlation Structure of Errors.}\label{app:MC_corr}

We now describe a set of Monte Carlo experiments which assess the performance of our estimator in realistic scenarios that replicate serial and spatial dependence structure of the data. To replicate our dataset, we simulate data for 1172 pairs, each for 32 periods. We assume that the pairs are located along 108 border segments, with the same mapping between pairs and border segments as in the data. The true data generating process is given below:
\[
y_{pt}=\alpha \Delta b_{p,t}+\lambda_{p}'F_{t}+\varepsilon_{p,t},
\]
where $F_{t}$ is the vector of factors and $\lambda_{p}$ is the factor loading for pair $p$. We assume, following our benchmark estimates, that the data are generated by two factors so that $\lambda_{p}$ and $F_{t}$ are $2\times1$. We take $\Delta b_{p,t}$ from the data and the estimated values of $\alpha$, $\lambda_{p}$ and $F_{t}$ from the benchmark estimation. We specify the error term $\varepsilon_{pt}$ as the sum of two uncorrelated components -- a border-specific component $\varepsilon_{pt}^{b}$ and an idiosyncratic component $\varepsilon_{pt}^{i}$, with weight $\rho^{s}$ on the common component (to capture the degree of spatial correlation):
\[
\varepsilon_{pt}=\sqrt{\rho^{s}}\varepsilon_{pt}^{b}+\sqrt{(1-\rho^{s})}\varepsilon_{pt}^{i}.
\]

All county pairs along a state border $b$ share $\varepsilon_{pt}^{b}$, allowing
us to capture the spatial correlation in residuals. We model $\varepsilon_{pt}^{b}$
as iid across different border segments with a Gaussian innovation $\zeta_{pt}^{b}$ and $\varepsilon_{pt}^{i}$ as iid across border county pairs with a Gaussian innovation $\zeta_{pt}^{i}$. Both  $\varepsilon_{pt}^{b}$ and $\varepsilon_{pt}^{i}$ are potentially autocorrelated as determined by $\rho^{t}$:
\begin{eqnarray*}
\varepsilon_{pt}^{b}&=\rho^{t}\varepsilon_{pt-1}^{b}+\zeta_{pt}^{b},\\
\varepsilon_{pt}^{i}&=\rho^{t}\varepsilon_{pt-1}^{i}+\zeta_{pt}^{i}.
\end{eqnarray*}

The results of all Monte Carlo experiments in this Section are reported in Table \ref{tab:Monte-Carlo-Results}. Each experiment runs a Monte Carlo with 2000 repetitions. For each value of $\rho^{t}$, we adjust the standard deviation of innovations $\zeta_{pt}^{i}$ and $\zeta_{pt}^{b}$ so that the unconditional variances of $\varepsilon_{pt}^{i}$ and $\varepsilon_{pt}^{b}$ are constant across experiments. The overall unconditional time series variance of $\epsilon_{pt}$ across pairs, constant across the experiments, is given by the baseline data estimate of $0.0008641$, the estimate for the serial correlation is $\rho^{t}=0.08$ (s.e. 0.02) and the estimate for the spatial correlation is $\rho^{s}=0.56$ (s.e. 0.005).

We verify the properties of our estimator under realistic correlation structures in the data, namely centered around the baseline estimates of serial and spatial correlation and three standard errors in each direction around that point. The results indicate that the estimator recovers well the effect of interest given the moderate serial and spatial correlation implied by the data.


\begin{table}[t]
\caption{Monte Carlo Results}
\label{tab:Monte-Carlo-Results}
\centering{}%
\input{Monte_Carlo_Results}
\parbox{13.0cm}{\noindent \footnotesize Note - The table collects the results of all Monte Carlo experiments described in Appendix \ref{app:MC_corr}.}
\end{table}





%%%%%%%%%%%%%%%%%%%%%%%%%%%%%%%%%%%%%%%%%%%%%%%%%%%%%%%%%%%%%%%%%%%%%%%%%%%%%%%%%%
\clearpage
\section{Additive Fixed Effects Specification}
\label{app:OLS}
In this Appendix, we consider a simpler and commonly used additive fixed effects model. We first test it against the interactive effects model and find that the additive fixed effects model is rejected. Nevertheless, we replicate every result in Table \ref{tab:Benefits_on_unemp} in the main text that was based on the interactive effects model using the additive fixed effects model. The resulting estimates do not affect any of the conclusions in the paper.


\subsection{Testing Factor Model against Additive Fixed Effects}

Section 9 of \citet{bai:09} provides two tests to determine which of the two models - the additive fixed effects model (Model A) or the interactive effects model (Model B) - is a better description of the data. The first method is the Hausman test, and the second is based on estimating the number of factors.

\underline{Hausman test}
We apply the test between additive and interactive effects models developed in Section 9 of \citet{bai:09}.  The null hypothesis is that the additive effects model is the true one:

\begin{equation}
    \Delta\tilde{u}_{p,t}=\alpha \Delta\log(b_{p,t})+\eta_p+\xi_t+\epsilon_{p,t},
    \label{eq:additive}
\end{equation}
and the alternative hypothesis is the interactive effects model:
\begin{equation}
    \Delta\tilde{u}_{p,t}=\alpha \Delta\log(b_{p,t})+\lambda'_p F_t+\epsilon_{p,t}.
    \label{eq:factor}
\end{equation}
Note that the additive effects model is nested in the interactive effects model with $\lambda_p=(\eta_p,1)$ and $F_t=(1,\xi_t)$.

As explained by \citet{bai:09}, the interactive-effects estimator for $\alpha$ is consistent under both models (\ref{eq:additive}) and (\ref{eq:factor}), but is less efficient than the least squares dummy-variable estimator for model (\ref{eq:additive}), as the latter imposes restrictions on factors and factor loadings. The fixed effect estimator is inconsistent under (\ref{eq:factor}), and the principle of the Hausman test is applicable.

The $J$-statistic computed from the Hausman test should converge in distribution to a $\chi^2$. However, the result of our implementation of the Hausman test is a \emph{negative} $J$-statistic. The standard error is smaller for the interactive-effects estimator (0.0039) than the additive-effects one (0.0061). Following \citet{bai:09} and using his notation as on page 1257, this means that
$$
var(\hat{\beta}_{IE}-\hat{\beta}_{FE}) = var(\hat{\beta}_{IE}) - var(\hat{\beta}_{FE}) < 0,
$$

As explained in \citet{schreiber:08}, the $J$-statistic can be negative asymptotically. Therefore, in large samples (our case here with large $N$ and $T$), a negative test statistic is only compatible with the alternative model and should be interpreted accordingly. Thus, we reject the additive effects model. Given that the additive effects model is rejected, results from \citet{bai:09} and \citet{gobillon+magnac:13} show that the estimator is inconsistent and rejected by the data, which is why we kept the interactive effects model as the benchmark.


Despite rejecting the additive effects model with the first test, we move to the second test.

\underline{Number of factors test}
The starting point is that \citet{bai:09} shows that the number of factors can be consistently estimated. He then uses this result to determine whether an additive model or interactive model is a better description of the data.

The first step of the test demeans the data within border county pairs so that the time-invariant county-pair fixed effects drop out. Then the interactive effects model is estimated on these demeaned data and the optimal number of factors is estimated. If the additive effects model is correct, the number of factors should be zero. If the number of factors is one or larger, the additive effects model is not appropriate and the interactive effects model is preferred.

We implement this test and find two factors in the demeaned data, rejecting the additive-effects model.

Both testing approaches offered in \citet{bai:09} reject the additive effects model in favor of interactive effects. We, therefore, report the results of the interactive effects model, which is consistent, in the main text.

\subsection{Results based on Additive Fixed Effects Specification}
Although the above analysis implies that estimates based on the additive fixed effects specification are biased and inconsistent (since the model is rejected), for completeness we report a full set of results based on this specification in Table \ref{tab:Benefits_on_unemp_OLS}. These estimates do not present any substantive contradiction to any of the results based on the factor model reported in the main text.


\vspace{0.5cm}\begin{table}[ht]
  \centering
   \caption{Unemployment Benefit Extensions and Unemployment: Results Based on Additive Effects Specification}\label{tab:Benefits_on_unemp_OLS}
\small
\input{Benefits_on_unemp_OLS.tex}
\parbox{16.5cm}{\noindent Note - Clustered standard errors in parentheses. Bold indicates $p<0.01$.\\
Column (1) - Baseline sample,\\
Column (2) - Baseline sample controlling for State GDP per worker,\\
Column (3) - Scrambled border county pairs sample,\\
Column (4) - Scrambled border county pairs sample controlling for State GDP per worker,\\
Column (5) - Sample of border counties with  $<15\%$ share of state's employment,\\
Column (6) - Sample of border counties with similar industrial composition,\\
Column (7) - Sample of border counties with population centers $<30$ miles apart,\\
Column (8) - Sample of border counties within the same Core Based Statistical Areas,\\
Column (9) - Baseline sample with perfect foresight measure of available benefits,\\
Column (10) - Baseline sample with controls for all other state-level policies.
}
\vspace{-0.0cm}
\end{table}



\subsection{Additional Controls}

The unobserved aggregate factors are treated as latent variables which are inferred from the data using the model. The county-specific loadings on each latent factor are also unobserved and are estimated at the same time. In fact one of the advantages of this approach is to reduce the dimensionality of the problem where the unobserved factors may represent a combined effect of numerous individual aggregate variables that one cannot include individually into the regression.

Given the nature of the Great Recession, one might be concerned that factors are standing in for observable variables like county-level debt, income and house prices. It seems difficult to estimate the effects of those variables consistently. The problem is that these variables at the county level are clearly endogenous. Thus, including them into the specification leads to a potential bias of all coefficients, including those on unemployment benefits. In Table \ref{tab:Endog_Vars}, we nevertheless explore the consequences of including these variables.

In Column 1 we add these variables to our baseline specification. We still select two factors as the optimal number of factors. The estimated effect of unemployment benefits on unemployment rises marginally relative to our benchmark. In Column 2 we consider an OLS regression with county-pair fixed effects (as we did in the section above). This has a minor effect on the coefficient of interest. Finally, in Column 3 we consider an instrumental variable regression where we attempt to account for the endogeneity of county-level variables. We use instruments employed in the literature. One is the well known Saiz instrument \citep{saiz:10} that identifies the elasticity of investment to house prices primarily using geographic features of areas (e.g. mountains, bodies of water) that restrict the supply of available land. The other instrument is the Wharton Residential Land Use Regulatory Index \citep{gyourko+saiz+summers:08}, which measures the stringency of local land regulation, which affects the supply of land that can be developed. Unfortunately, both of instruments are only available at the MSA level (which only covers a subset of county pairs, leading do a significant reduction in the number of observations). We assume that the instrument takes the same value for all counties within the MSA. Moreover, it is constant over time, which precludes us from including fixed effects in the IV regression in Column 3. Nevertheless, the substantive results are relatively little changed, although the coefficient is less precisely estimated, in part because the sample size declines significantly. 

\begin{table}[H]
  \centering
	\captionsetup{justification=centering}
   \caption{Unemployment Benefit Extensions and Unemployment, \\ Controlling for County-Level House Prices, Income, and Debt} \label{tab:Endog_Vars}
\small
    \begin{tabular}{lccccc} \hline
    VARIABLES & Factor Model &&  OLS &&  IV-GLS \\ \cline{2-2} \cline{4-4} \cline{6-6}
    				  & (1) && (2) && (3) \\ \hline
Weeks of  Benefits & \bf{0.059} && \bf{0.044} && \bf{0.046} \\
									 & (0.000)	  && (0.000)	 	 &&  (0.001)\\

N. factors   & 2 && -- && --   \\
Observations & 23,332 && 23,318 && 4,216  \\
R-squared    & 0.479 && 0.073 && -0.017  \\ \hline
\multicolumn{6}{c}{\footnotesize Note - $p$-values (in parentheses). Bold indicates $p<0.01$.} \\
%\multicolumn{5}{l}{\footnotesize } \\
\end{tabular}
\vspace{-0.2cm}
\end{table}


%%%%%%%%%%%%%%%%%%%%%%%%%%%%%%%%%%%%%%%%%%%%%%%%%%%%%%%%%%%%%%%%%%%%%%%%%%%%%%%%%%
%--------------------------------------------------------------------------------%
%%%%%%%%%%%%%%%%%%%%%%%%%%%%%%%%%%%%%%%%%%%%%%%%%%%%%%%%%%%%%%%%%%%%%%%%%%%%%%%%%%
\section{Validation using Model-Generated Data}
\label{sec:simulation}

In this Section we evaluate the performance of our empirical method on data generated by a calibrated equilibrium search model. The model is an extension of \citet{mortensen+pissarides:94} to allow for unemployment benefit expiration.

To address the border county design, the model features a nested state-county structure. In particular, there is a stochastic process for state's productivity.

 The unemployment benefit policy depends on the endogenous unemployment level in the state economy.  The county economy takes the endogenously induced joint stochastic process for state unemployment, productivity and benefits as exogenous. The assumption is that counties are ``small'' relative to the state of which they are apart. We explain that this assumption is consistent with the empirical evidence presented above and quantify the implications of alternative assumptions below.

Preferences, technology and frictions are the same across the state and county economies.

\vspace{-0.4cm}\paragraph{\textit{Agents.}} In any given period, a worker can be either employed (matched with a firm) or unemployed. Risk-neutral workers maximize expected lifetime utility
\begin{equation*}
U=\E_0 \sum\limits_{t=0}^{\infty} \beta^t c_t,
\end{equation*}
where $\E_0$ is the period-0 expectation operator, $\beta \in \left(0,1\right)$ is the discount factor, $c_t$ denotes consumption in period $t$.  An unemployed worker produces $h$, which stands for the combined value of leisure and home production.  In addition, unemployed workers may be eligible for benefits $b$.  Unemployed workers who are eligible for benefits lose eligibility stochastically at rate $e_t(\cdot)$, which depends on the state unemployment rate as specified below.

Firms are risk-neutral and maximize profits. Workers and firms have the same discount factor $\beta$. A firm can be either matched to a worker or vacant. A firm posting a vacancy incurs a flow cost $k$.

\vspace{-0.4cm}\paragraph{\textit{Matching.}}
The number of new matches in period $t$ is given by $M(u_t,v_t)$, where $u_t$ is the number of unemployed in period $t$, and $v_t$ is the number of vacancies.  The matching function is assumed to be constant returns to scale, and strictly increasing and strictly concave in both arguments.  We define $\theta_t=v_t/u_t$ as the market tightness in period $t$. We then define the job-finding probability as $f(\theta_t)=M(u_t,v_t)/u_t=M(1,\theta_t)$ and the probability of filling a vacancy as $q(\theta_t)=M(u_t,v_t)/v_t=M(1/\theta_t,1)$. By the assumptions on $M$ made above, the function $f(\theta_t)$ is increasing in $\theta_t$ and $q(\theta_t)$ is decreasing in $\theta_t$. Existing matches are destroyed with exogenous job separation probability $\delta$.

\vspace{-0.4cm}\paragraph{\textit{Production.}}  A matched worker-firm pair produces output $z_t$, which follows a first order Markov process.  Firms pay workers a wage $w_t$, determined through Nash bargaining with workers' bargaining power $\xi$. Thus, the period profit of a matched firm is given by $\pi_t = z_t - w_t$.


\subsection{State Economy}

In the state economy the benefit expiration policy depends on the state unemployment rate, $e_t(\ust)$.  We assume ineligible workers regain eligibility as soon as they are matched with a firm.  The relevant state variables for the state economy are thus the exogenous state productivity $\zst$ and the endogenous unemployment rate $\ust$.  Let $\ost=(\zst,\ust)$. The state law of motion for employment is therefore:
\begin{eqnarray}
L^S_{t+1}(\ost)=(1-\delta)L^S_{t}+f(\tst)\left(1-L^S_t\right) \label{LOML}
\end{eqnarray}
and $\ust=1-L^S_t$.

\vspace{-0.3cm}\paragraph{\textit{Value Functions.}}
The flow value for a firm employing a worker is
\begin{equation}
J^S_t(\ost) = \zst-w^S_t+\beta\left(1-\delta\right)\E J_{t+1}(\Omega^S_{t+1}) \label{TJ}
\end{equation}
and the flow value of a vacant firm is:
\begin{equation}
V^S_t(\ost)=-k+\beta q\left(\tst\right)\E J_{t+1}(\Omega^S_{t+1}),
\label{TFE}
\end{equation}
where $k$ is the flow cost of maintaining a vacancy.  The surplus for a firm employing a worker is thus $J^S_t-V^S_t$. The value functions for workers can be written as:
\begin{eqnarray}
W^S_t(\ost) &=& w^S_t +\beta\left(1-\delta\right)\E W^S_{t+1} + \beta\delta\left(1-e_t(\ost)\right) \E
U_{t+1}^{S,E}(\Omega^S_{t+1}) \notag \\
&& +\beta\delta e_t(\ost) \E U_{t+1}^{S,I}(\Omega^S_{t+1}),  \label{TW} \\ [2mm]
U^{S,E}_t(\ost) &=& h+b +\beta f\left(\tst\right)\E W^S_{t+1}(\Omega^S_{t+1})  +\beta\left(1-f\left(\tst\right)\right)\left(1-e_t(\ost)\right)\E U_{t+1}^{S,E}(\Omega^S_{t+1})\notag \\
&&   +\beta\left(1-f\left(\tst\right)\right)e_t(\ost) \E U_{t+1}^{S,I}(\Omega^S_{t+1}),  \label{TUE} \\ [2mm]
U^{S,I}_t(\ost) &=& h+\beta f\left(\tst\right)\E W_{t+1}(\Omega^S_{t+1}) + \beta\left(1-f\left(\tst\right)\right)\E U_{t+1}^{S,I}(\Omega^S_{t+1}), \label{TUI}
\end{eqnarray}
where $W^S_t$ s the value of a job for a worker, $U^{S,E}_t$ is the value of unemployment for an agent eligible for benefits and
$U^{S,I}$ is the value of unemployment for a non-eligible agent. Define the surplus of being employed as $\Delta_t^{S,E}=W^S_t-U_t^{S,E}$. Also define the surplus for an unemployed worker of being eligible: $\Phi^S_t=U_t^{S,E}-U_t^{S,I}$. The laws of motion for these quantities are:
\begin{eqnarray}
\Delta_t^{S,E}(\ost)&=&w^S_t-h-b+\beta\left(1-\delta-f\left(\tst\right)\right)\E \Delta_{t+1}^{S,E}(\Omega^S_{t+1})  \notag \\
&& +\beta\left(1-\delta-f\left(\tst\right)\right)e_t(\ost)\E\Phi^S_{t+1}(\Omega^S_{t+1}), \label{TDelta} \\ [2mm]
%\notag \\
\Phi^S_t(\ost)&=&b+\beta\left(1-f\left(\tst\right)\right)\left(1-e_t(\ost)\right)\Phi^S_{t+1}(\Omega^S_{t+1}). \label{TPhi}
\end{eqnarray}
The wage is chosen to maximize:
\beq
\left(\Delta_t^{S,E}\left(\ost\right)\right)^\xi\left(J^S_t\left(\ost\right)-V^S_t\left(\ost\right)\right)^{1-\xi}.
\label{wageS}
\eeq

\vspace{-0.3cm}\paragraph{\textit{State Equilibrium Definition.}}
Given a policy $\left(b,e_t\left(\cdot\right)\right)$ and an initial condition $\Omega^S_0$ an equilibrium is a sequence of $\ost$-measurable functions for wages $w_t$, market tightness $\tst$, employment $L^S_t$, and value functions
$\left\{W^S_t, U_t^{S,E}, U_t^{S,I}, J^S_t, V^S_t, \Delta^S_t\right\}$ such that:
\begin{enumerate}
\vspace{-0.2cm}\item The value functions satisfy the worker and firm Bellman Eqs (\ref{TJ}), (\ref{TFE}), (\ref{TW}), (\ref{TUE}),  (\ref{TUI});
\vspace{-0.2cm}\item Free entry: The value $V^S_t$ of a vacant firm is zero for all $\ost$;
\vspace{-0.2cm}\item Nash bargaining: The wage satisfies equation (\ref{wageS});
\vspace{-0.2cm}\item Law of motion for employment: The employment process satisfies (\ref{LOML}).
\end{enumerate}

\subsection{County Economy}

The county is assumed to be small with respect to the state of which it is a member.  That is, the county unemployment rate is not assumed to affect the state unemployment rate and the county productivity process is orthogonal to the state one.  The benefit expiration policy for the county, however, depends on the state unemployment rate.  Thus, in addition to exogenous county productivity, $z^C$, the state productivity and the state unemployment rate will be state variables (since they are jointly sufficient to forecast benefit policy).  Thus, denote the vector of states for the county $\oct=\left(\zct;\zst,\ust\right)$.  

The law of motion for county employment is:
\begin{eqnarray}
L^C_{t+1}(\oct)=(1-\delta)L^C_{t}+f(\tct)\left(1-L^C_t\right). \label{app_eq:LOML}
\end{eqnarray}
and $\uct=1-L^C_t$.

\vspace{-0.3cm}\paragraph{\textit{Value Functions.}}
The flow value for a firm employing a worker is
\begin{equation}
J^C_t(\oct) = \zct-w^C_t+\beta\left(1-\delta\right)\E J_{t+1}(\Omega^C_{t+1}) \label{TJC},
\end{equation}
and the flow value of a vacant firm is:
\begin{equation}
V^C_t(\oct)=-k+\beta q\left(\tct\right)\E J^C_{t+1}(\Omega^C_{t+1}).
\label{TFEC}
\end{equation}
The surplus for a firm employing a worker is thus $J^C_t-V^C_t$.

The value functions for workers can be written as:
\begin{eqnarray}
W^C_t(\oct) &=& w^C_t +\beta\left(1-\delta\right)\E W^C_{t+1} + \beta\delta\left(1-e_t(\oct)\right) \E U_{t+1}^{C,E}(\Omega^C_{t+1}) \notag \\
&& +\beta\delta e_t(\oct) \E U_{t+1}^{C,I}(\Omega^C_{t+1}),  \label{TWC} \\[2mm]
%\notag \\
U^{C,E}_t(\oct) &=& h+b +\beta f\left(\tct\right)\E W^C_{t+1}(\Omega^C_{t+1})  +\beta\left(1-f\left(\tct\right)\right)\left(1-e_t(\oct)\right)\E U_{t+1}^{C,E}(\Omega^C_{t+1})\notag \\
&&   +\beta\left(1-f\left(\tct\right)\right)e_t(\oct) \E U_{t+1}^{C,I}(\Omega^C_{t+1}),  \label{TUEC} \\[2mm]
%\notag \\
U^{C,I}_t(\oct) &=& h+\beta f\left(\tct\right)\E W_{t+1}(\Omega^C_{t+1}) + \beta\left(1-f\left(\tct\right)\right)\E U_{t+1}^{C,I}(\Omega^C_{t+1}). \label{TUIC}
\end{eqnarray}
Define the surplus of being employed as $\Delta_t^{C,E}=W^C_t-U_t^{C,E}$. Also define the surplus for an unemployed worker of being eligible: $\Phi^C_t=U_t^{C,E}-U_t^{C,I}$. The laws of motion for these quantities are:
\begin{eqnarray}
\Delta_t^{C,E}(\oct)&=&w^C_t-h-b+\beta\left(1-\delta-f\left(\tct\right)\right)\E \Delta_{t+1}^{C,E}(\Omega^C_{t+1})  \notag \\
&& +\beta\left(1-\delta-f\left(\tct\right)\right)e_t(\oct)\E\Phi^C_{t+1}(\Omega^C_{t+1}),
 \label{TDeltaC}\\[2mm]
%\notag \\
\Phi^C_t(\oct)&=&b+\beta\left(1-f\left(\tct\right)\right)\left(1-e_t(\oct)\right)\Phi^C_{t+1}(\Omega^C_{t+1}). \label{TPhiC}
\end{eqnarray}

The wage is chosen to maximize:
\beq
\left(\Delta_t^{C,E}\left(\ost\right)\right)^\xi\left(J^C_t\left(\ost\right)-V^C_t\left(\ost\right)\right)^{1-\xi}.
\label{wageC}
\eeq

\vspace{-0.4cm}\paragraph{\textit{County Equilibrium Definition.}}
Taking as given an initial condition $\Omega^C_0$, benefit expiration policy, and the joint stochastic process for state productivity and unemployment, we define an equilibrium given policy:

\begin{definition}
Given a policy $\left(b,e_t\left(\cdot\right)\right)$ and an initial condition $\Omega^C_0$ an equilibrium is a sequence of $\oct$-measurable functions for wages $w_t$, market tightness $\tct$, employment $L^C_t$, and value functions
\[\left\{W^C_t, U_t^{C,E}, U_t^{C,I}, J^C_t, V^C_t, \Delta^C_t\right\}\] such that:
\begin{enumerate}
\vspace{-0.2cm}\item The value functions satisfy the worker and firm Bellman equations (\ref{TJC}), (\ref{TFEC}), (\ref{TWC}), (\ref{TUEC}),  (\ref{TUIC});
\vspace{-0.2cm}\item Free entry: The value $V^C_t$ of a vacant firm is zero for all $\oct$;
\vspace{-0.2cm}\item Nash bargaining: The wage satisfies equation (\ref{wageC});
\vspace{-0.2cm}\item Law of motion for employment: The employment process satisfies (\ref{app_eq:LOML});
\vspace{-0.2cm}\item The joint process for $\left(z^S_t,u^S_t\right)$ is consistent with the state equilibrium.
\end{enumerate}
\end{definition}

\subsection{Calibration}

\begin{table}[t]
\centering
\caption{Internally Calibrated Parameters}
\small
\input{Model_Calibration.tex}
\label{calib}
\vspace{-0.2cm}
\end{table}


The calibration strategy we employ is to require the state economy to be consistent with key labor market statistics and to match the effect of unemployment benefit extensions on unemployment estimated in Section \ref{sec:baseline_results}. The model period is taken to be one week.  We match the average labor market tightness, the average job-finding rate, and the regression coefficient of quasi-differenced unemployment on benefit duration.  The calibrated parameters are summarized in  Table \ref{calib}.  In order to be consistent with the existing EB program, in the calibration we set benefit expiration policy at 26 weeks when state unemployment is less than 6.5\%, 39 weeks when unemployment is between 6.5\% and 8\% and 46 weeks when greater than 8\%. The level of benefits $b$ that eligible workers receive remains fixed for any unemployment rate at 0.4. The remainder of the parameters are calibrated externally, using the same values and parametric forms for the matching function as \citet{hagedorn+manovskii:08}.

\subsection{Quantitative Evaluation}

The goal of the simulation exercise is to generate synthetic data at the county level comparable to the actual data.  We simulate two states and one county in each of them. The two states and the two counties each have the same process for productivity.  The counties, consistent with our border county assumption, have the same \textit{realized sequence of shocks}.  The two states, however, have \text{different realized sequences of productivity shocks}.  Consequently, the realized endogenous sequences of state unemployment will be different.  Thus, the two counties will have a different time series of unemployment benefits.

We simulate the two states and the two counties for 100 years and throw out the first 15 years of data as ``burn-in.''  We then estimate the same regression (with quasi-differenced unemployment on the left-hand side) as we do on the data from the Great Recession. Recall that our calibration strategy ensures that coefficient on the difference in benefits in this regression is the same in the data and in the simulations of the model. Then, we calculate the effect of a permanent 10-week increase in benefits on unemployment, vacancies and tightness.  We then compare these true permanent effects from the model to the calculated permanent effects from the data. The results and relevant comparisons are displayed in Table \ref{valid}.  The model generated data confirms the empirical validity of our specification, as our model, calibrated to generate the same regression coefficient on unemployment benefit duration from the data delivers right permanent effects on unemployment, vacancies and tightness.




\begin{table}[t]
\centering
\small
  \captionsetup{justification=centering}
  \caption{Estimated Permanent Effect of a 10 Week Benefit Extension \\ from Regressions Coefficients in Model Generated Data}\label{valid}
\input{Model_Validation.tex}
\vspace{-0.2cm}
\end{table}


\subsection{The ``Small'' County Assumption}

We have performed the quantitative analysis using a simple model in which counties are ``small'' relative to the state they belong to. This is without loss of generality if economic fundamentals evolve smoothly across state borders, as we have found to be the case in the data in Section \ref{sec:endogeneity_test}. In this case the size of the border counties relative to their state's and the correlation between the state and county productivities are not a relevant consideration for our analysis and can be safely ignored when simulating the model. In other words, considering a more sophisticated model with some ``large'' counties and correlation between county and state productivity will not impact the findings.



We can use the model, however, to assess the potential consequences of the violation of our identifying assumption that state-level variables that induce state-level policies, such as unemployment benefit extensions, are differenced out between border counties. To this end, we now quantify the size of the potential bias that would be induced by the extreme alternative assumption that shocks to county productivity affect only their state's productivity and have no impact on the neighboring states. In our data each border county, conditional on accounting for less than 15\% of state employment (there are $1,113$ such counties in our sample), accounts for on average $2\%$ of state employment. In contrast, each border county, conditional on accounting for more than 15\% of state employment (there are $59$ such counties in our sample), accounts for on average $35\%$ of state employment. We now simulate the model using these shares as a measure of correlation between each county and its state's productivity (and assume that shocks to county productivity are orthogonal to productivity of the neighboring state). In other words, the correlation between county and its state productivity is set to $0.02$ for $1,113$ counties and to $0.35$ for $59$ counties. The results of performing this experiment indicate that the estimated coefficients are biased but the bias is quite small relative to the true effect. Specifically, estimating the effects of benefit extensions on unemployment in the model that features this correlation returns the estimate of 0.0569 while the true coefficient is 0.0532.


%%%%%%%%%%%%%%%%%%%%%%%%%%%%%%%%%%%%%
%-----------------------------------%
%-----------------------------------%
%%%%%%%%%%%%%%%%%%%%%%%%%%%%%%%%%%%%%
\section{Change in Location of Employment in Response to \\ Changes in Benefits}
\label{sec:mobility}

 A potential concern arises from the observation that households may live in different states than where they work. This would bias our estimates if the households systematically change their job search behavior in response to changes in unemployment benefits. For example, if households search in states with less generous benefits to take advantage of a higher job-finding rate, our estimate of the effect of benefit extensions on unemployment would be biased \emph{downwards}, since those households would face a higher job-finding rate, which would translate into a lower unemployment rate in that county. In this section, we use two different methods to show that our analysis is not affected by such a bias. First, we construct an upper bound on the effect of cross-border mobility using direct empirical evidence of job search behavior. Second, we develop an imputation procedure that allows to estimate the effects of unemployment benefit extensions while fully accounting for mobility. Both approaches confirm that search behavior does not vary systematically with changes in benefits, validating our use of a simple and transparent specification that ignores mobility decisions.
This finding is not very surprising. Residents of the border counties face a trade-off between receiving higher wages with lower job-finding probability in a county belonging to the state with higher benefit eligibility and receiving lower wages with higher job-finding probability in the state with lower benefits (recall that benefits depend on the state of employment, and not on the state of residence). Moreover, the difference in the available duration of benefits across border counties is typically small and may not justify larger commuting expenses.

We begin by using data from the LEHD Origin-Destination Employment Statistics (LODES) dataset to construct counts of how many people live in a county but work in the county across the state border and the change in the number of such commuters over time.\footnote{The data are available at https://lehd.ces.census.gov/data/lodes/LODES7/. We use the origin-destination data ``OD,'' which provides counts of commuters at the Census Block level. Data tabulates these flows by job type. We use data constructed for primary jobs, which is the highest paying job for an individual worker for the year. We
aggregate these flows to county level and obtain the number of commuters
for all county pairs in the data. Data are annual, cover the period
2005--2012, and are available for 50 states and the District of Columbia, with the exceptions of D.C. over 2005-2009 period and Massachusetts over 2005-2010 period.} In order to bound the potential effect that commuters may have on our estimates, we make an extreme assumption that instead of being employed in the border county, those individuals would have been unemployed in their home county.\footnote{To correct for time aggregation due to the fact that our unemployment count data is quarterly, but the commuter data is annual, we adjust by half the change in commuters across all quarters in the year, which assumes that on average each commuting change would have taken place halfway through the year.} Accordingly, we construct a counterfactual unemployment rate given by adding to the count of unemployed in the home county the change in the number of commuters and then divide that total by the labor force in the county.\footnote{We do not need to add the change in commuters to the labor force in the denominator as it is measured based on the place of residence, so that commuters working in the border county are already counted as part of the labor force in the home county.} As in the benchmark analysis, we then take logs, quasi-difference this counterfactual unemployment rate, and then difference across the border counties. We use the same factors as in the benchmark analysis but re-estimate the factors loadings and coefficient on the difference in the quasi-differenced log benefit duration. The resulting estimate of the effect of benefit duration on unemployment reported in Table \ref{tab:Mobility_LODES} is virtually unchanged from our benchmark. Thus, even under the extreme assumption that commuters would have been unemployed in their home counties, we see no evidence of a bias in our estimates induced by worker mobility across state borders.


\begin{table}[t]
\centering
\small
	\captionsetup{justification=centering}
   \caption{Unemployment Benefit Extensions and \\ Mobility Across State Borders. LODES data.} \label{tab:Mobility_LODES}
\input{Mobility_LODES.tex}
\vspace{-0.3cm}
\end{table}

Next, we develop an imputation procedure that fully accounts for mobility in our estimates. Because integrated labor markets generally contain multiple neighboring counties, instead of focusing on the county pair as the unit of analysis for search behavior we aggregate all counties on both sides of a border segment and perform the imputation on that ``border segment'' pair. To impute what fraction of workers search in the state where they live, consider the following model. We consider the local economy to consist of a pair of state border segments $A$, $B$.  The segments are populated by labor forces of size $n_t^A$ and $n_t^B$ (taken as the sum of all the border county labor forces in each state on the resp. side of the border) and populations $\pat$ and $\pbt$.

In any given period, a worker can be either employed (matched with a firm), unemployed or not in the labor force. In period $t$, firms in state $A$ post vacancies in state $A$, $\vat$.  An unemployed worker in state $A$ searches either in state $A$ or in state $B$.  We assume that a fraction $\zeta$ of non-labor force participants (observed in the LAUS data) enter the labor force and search for jobs.  The number of new matches in state A in period $t$ equals
\begin{equation*}
M\left(\uat, \vat\right),
\end{equation*}
where $\uat$ is the measure of individuals in period $t$ searching in state $A$. The number of matches is the same for state $B$ \textit{mutatis mutandis}. $M(\cdot,\cdot)$
exhibits constant returns to scale and is strictly increasing and  strictly concave in both arguments.

We define
\begin{equation*}
\tat = \frac{\vat}{\uat}
\end{equation*}
to be the market tightness in state $A$ in period $t$.  We define the job-finding and vacancy-filling probabilities as in Appendix \ref{sec:simulation}.






We can thus write for the number of unemployed searching in state A and B, resp.:
\begin{align} \label{eq:utilde}
\uat&=(\uatt+\zeta(\pat-\nat))\xat+(1-\xbt)(\ubtt+\zeta(\pbt-\nbt)),\\
\ubt&=(\ubtt+\zeta(\pbt-\nbt))\xbt+(1-\xat)(\uatt+\zeta(\pat-\nat)), \label{eq:utilde2}
\end{align}
where $\tilde{u}^{i}_{t}$ is the observed number of unemployed who live in state $i$, $x^{i}_{t}$ is the fraction of the unemployed in state $i$ that searches in state $i$, and $\delta_t$ is the separation probability into unemployment, calculated from the Current Population Survey (CPS) following \citet{shimer:05b}.  We follow \citet{hall:13} and set $\zeta$ to $\nicefrac{5}{27}$ to match the ratio of the job-finding rates of non-participants to the unemployed in the CPS.




We can measure the probabilities for an unemployed worker from states A and B to find a job, $\phi_t^A$ and $\phi_t^B$, in the data:
\begin{align} \label{eq:JFR_Data}
\phi_t^A = \frac{\uatt-\uattp+\delta_t \left(\nat-\uatt\right)}{\uatt} ,\\
\phi_t^B = \frac{\ubtt-\ubttp + \delta_t \left(\nbt-\ubtt\right) }{\ubtt},
\end{align}
as all right-hand variables are observed in the data.
We can relate the measurable $\phi_t^A$ and $\phi_t^B$ to the unobservable variables $\xat$, $\xbt$, $f(\tat)$, $f(\tbt)$:

\begin{align} \label{eq:JFR}
\phi_t^A = \xat f\left(\tat\right)-\left(1-\xat\right) f\left(\tbt\right),\\
\phi_t^B = \xbt f\left(\tbt\right)-\left(1-\xbt\right) f\left(\tat\right).\label{eq:JFR2}
\end{align}
The four equations $(\ref{eq:utilde})$, $(\ref{eq:utilde2})$, $(\ref{eq:JFR})$ and $(\ref{eq:JFR2})$ have $4$ unknowns, $\xat, \xbt, f(\tat), f(\tbt)$.\footnote{We do not directly observe $\xat$, and thus we don't observe $\uat$ and $\tat$, nor the matching function.} These equations are not linearly independent and thus do not allow us to recover these 4 unknowns. Instead they give us a set of solutions $S$.

In order to proceed to identify $\xat,\xbt$, we assume that the matching function is Cobb-Douglas, $\mu u^\gamma v^{1-\gamma}$.   Note, however, that we do not necessarily see the true level of vacancies.  However, if we assume that we see the same \textit{fraction}, $\psi$, of total vacancies for both counties in a pair, we can still estimate the effective matching function given our observed vacancies.  If we observe $\tilde{v}=\psi v$, then the total number of matches is $\tilde{\mu} u^\gamma \tilde{v}^{1-\gamma}$, where $\tilde{\mu}=\psi^{\gamma-1}\mu$.  Thus, we propose to identify $\tilde{\mu}$ and $\gamma$ in addition to the $x$'s.

We allow $\tilde{\mu}$ to change over time, to capture any possible time trends in the adoption of online vacancies.  The algorithm consists of selecting $\gamma$, $\{\tilde{\mu}_t,\xat,\xbt\}_{t=1}^T$ to minimize the error in the equations (\ref{eq:JFR}), (\ref{eq:JFR2}) and $\frac{q(\tat)}{q(\tbt)}=\left(\frac{\tbt}{\tat}\right)^{\gamma}$:
\begin{eqnarray}
\phi_t^A &=& \xat \tilde{\mu}_t\left(\frac{\vatt}{\uat}\right)^{1-\gamma}-\left(1-\xat\right)\tilde{\mu}_t\left(\frac{\vbtt}{\ubt}\right)^{1-\gamma} ,\\
\phi_t^B &=& \xbt \tilde{\mu}_t\left(\frac{\vbtt}{\ubt}\right)^{1-\gamma}-\left(1-\xbt\right) \tilde{\mu}_t\left(\frac{\vatt}{\uat}\right)^{1-\gamma},\\
\frac{1 - \frac{v^A_{t+1}-v^{A,new}_{t+1}}{v^A_t}}{1 - \frac{v^B_{t+1}-v^{B,new}_{t+1}}{v^B_t}}&=&\left(\frac{\frac{\vbtt}{\ubt}}{\frac{\vatt}{\uat}}\right)^{\gamma},
\end{eqnarray}
where we observe all left hand side variables for all $t$ border segment by border segment.\footnote{The probability to fill a vacancy $q_t = 1 - \frac{v_{t+1}-v^{new}_{t+1}}{v_t}$, where $v_t$ is the stock of vacancies at $t$ and $v^{new}_{t}$ are newly posted vacancies at $t$, so that $v_{t+1}-v^{new}_{t+1}$ are not filled vacancies from period $t$. Both $v_t$ and $v^{new}_t$ are observable in the data. Note, we are assuming here that all vacancies that drop out of the sample are filled. However, our measurement is not affected if some unfilled vacancies also drop out of the sample as long as the fraction of unfilled vacancies that drop out is the same across the two sides of the border segment (since we take the ratio of the imputed vacancy-filling rates).}


\begin{table}[t]
\centering
\small
\caption{Effect of UI Benefits on Imputed Labor Market Variables} \label{tab:Imputed_Results}
\begin{tabular}{lccc} \hline
VARIABLES & Out of State Search & Imputed Tightness  & Imputed Job-finding\\ \hline \noalign{\vskip 1mm}
Weeks of  Benefits& 0.0002 &  \bf{-0.1154}  & \bf{-0.0524}\\
 & (0.510) &  (0.000) & (0.000)\\ \noalign{\vskip 1mm}
Factors & 2 & 2 & 2\\
Observations & 29,492 & 29,492 & 29,492\\
 R-squared & 0.066 & 0.2816 & 0.2996\\ \hline
\multicolumn{4}{l}{\footnotesize Note - $p$-values (in parentheses) calculated via bootstrap. Bold indicates $p<0.01$.} \\
%\multicolumn{4}{l}{Bold font indicates $p<0.01$.} \\
\end{tabular}
\vspace{-0.3cm}
\end{table}

We measure the effect of benefits on search behavior by examining the difference between the imputed fraction of workers searching away from their home states $(1-\xat)-(1-\xbt)$.  Further, we construct imputed tightness by dividing county level vacancies by the imputed measure of unemployed workers searching in that county ($\nicefrac{\vat}{\uat}$), corrected for the search behavior along that border segment (we impose the same $x$'s for all counties within a state for each border segment).  Then, the job-finding rate is constructed using the imputed tightness and the estimated parameters of the matching function.    Table \ref{tab:Imputed_Results} Column (1) shows, using the difference-in-difference estimator with interactive effects, that there is only a very small and statistically insignificant response of search behavior, to changes in benefits, so that mobility does not bias our estimates.   Further, the effect on imputed tightness, which now fully accounts for changes in mobility in response to changes in benefits, is not statistically significantly different from the baseline estimate. %The effect of extending benefits to 82.5 weeks for approximately 16 quarters (the average during the Great Recession) on the quarterly job-finding rate would predict a drop from $77.6\%$ to $48.6\%$.


%%%%%%%%%%%%%%%%%%%%%%%%%%%%%
%---------------------------%
%%%%%%%%%%%%%%%%%%%%%%%%%%%%%

%\clearpage
\section{LAUS Data Quality}
\label{app:data_quality}
A potential concern with any empirical investigation is the quality of the underlying data. In the case of the analysis in this paper, the concern is with the construction of the county-level unemployment data by LAUS. In particular, if a component of county-level unemployment data was somehow imputed using state-level unemployment, this could give rise to the endogeneity problem discussed in Section \ref{sec:endogeneity_test}, in which both benefit duration and measured county unemployment would be driven by underlying economic conditions in the state, such as productivity or demand in the state $Z_{p,t}$. Fortunately, endogeneity tests revealed that this is not the case. In fact, we find that the difference in measured unemployment between pairs of border counties is not correlated with the difference in state-level productivity or the difference in instrumented state unemployment (more precisely, $Corr (\nu_{p,t}, \Delta Z_{p,t}) = 0$). This reveals that county-level unemployment estimated by the BLS does not reflect state-level variables to an important degree since (otherwise, this correlation would not be zero), reflecting the negative correlation between state-level productivity or demand and state unemployment.



In this Appendix, we describe the data construction by LAUS and implement modifications to their procedures to provide additional checks of the appropriateness of the LAUS data for the analysis in this paper.

\subsection{Removing Additivity Factors}
\label{app:data_quality_additivity}
The primary labor force survey used to measure unemployment in the US is the Current Population Survey (CPS). Unfortunately, this survey is not representative at the county level. The objective of the LAUS program is then to estimate county level employment and unemployment in a way that would match as closely as possible the estimate that would have been obtained if a representative labor force survey was conducted in each county. To do so, the LAUS program draws on a variety of data sources. In addition to the CPS, the LAUS relies on large-scale surveys and quarterly censuses of payroll employment as well as the data from the universe of jobs covered by the UI laws and data on the universe of UI claims. Using these data LAUS estimates using the observed relationships in the aggregate (not state-level) data the number of unemployed workers who left a UI covered job and remain unemployed at a particular point in time (have not exited the labor force). In addition, it uses the age distribution of the population in a county from the Census Bureau to predict the number of new entrants (or re-entrants) into the labor market who are not yet covered by the UI system. Taken together, this represents the estimate of county unemployment. The estimation of county-level number of new entrants is mainly based on aggregate relationships but it does use the five-year average state-level estimates of the number of new entrants. Although the use of only the long-run average of state-level variables in this step minimizes the concern that it may induce the endogeneity problem, we formally verify that it does not in Appendix \ref{app:data_quality_claims}.
Prior to doing so, we assess the last step of the procedure in which current state-level variables may indeed enter the county-level estimates. In this step, the LAUS multiplies the unemployment estimate in each county of the state by the same ``additivity factor'' necessary to ensure that the sum of unemployment estimates across all counties in a state adds up to the total estimate of state unemployment.

Introducing this additivity adjustment may cause a bias in our estimates if the relationship between LAUS estimate and unmeasured unemployment (the additivity error) varies across the state. The endogeneity tests performed above indicate the lack of such a bias. In contrast, if the proportionality assumption is approximately correct, then not applying the additivity factors may lead to biased estimates when comparing border counties. Fortunately, we can directly assess the consequences of using the additivity adjustment on our estimates by undoing this step of the LAUS procedure.\footnote{We are very grateful to the BLS for releasing the additivity factors underlying LAUS estimates to us.} Estimating the baseline specification using the data with LAUS additivity factors removed yields the coefficient of $0.054$ with $p-$value of $0.000$ on weeks of benefits. Thus, a direct comparison of the estimates on the data before and after the additivity adjustment is performed reveals that they are very similar.




\subsection{Using Administrative Unemployment Claims Data}
\label{app:data_quality_claims}
To enable an additional independent verification of the quality of LAUS data, the BLS has agreed to provide us with 2005-2012 administrative data from the unemployment insurance system. The data include all continuing unemployment claims in the regular state unemployment insurance benefit program (i.e., during the first 26 weeks of an unemployment spell) by county.\footnote{Importantly, unemployment benefit recipients must exhaust all state benefits before being eligible for federal extensions under either EUC or EB programs.} In addition, the data include the number of final payments, i.e. the number of unemployed who collected their final unemployment check under the regular state system (the $26^{\mbox{\footnotesize th}}$ week). The final payment data are available to us only at the Labor Market Area (LMA) level as defined by the BLS.\footnote{The LMA directory is available at http://www.bls.gov/lau/lmadir.pdf.} A Labor Market Area may include a single or multiple counties in one or more states. As a consequence, we restrict the sample to single county LMAs. This leaves us with a sample of 282 border county pairs and 9,024 observations. Importantly, these continuing claims and final payments data are pure counts, and are not subject to any imputation or other adjustments (in particular, these data contain neither estimation of the number of new entrants nor any adjustment for additivity - the two steps in LAUS estimation procedure that might in theory give rise to endogeneity problem).



These data allow us to measure the job-finding rate of unemployment insurance claimants in their first 26 weeks of compensated unemployment as follows. The counts of continued claims are provided monthly.  The claims are a count of the total number of continued claimants in the reference week of that month (the week including the 12th, as with the CPS).  The final payments data are at a weekly frequency and represent the count of claimants who receive their $26^{\mbox{\footnotesize th}}$ benefit check, i.e., the final payment from the regular state UI system.  The continued claims data, therefore include the number of final payments from the reference week of that month.  We can write the number of continued claimants in month $t$ as:
\begin{eqnarray}
{u}_t^{CC} = u^{FP}_{t,2} + \sum_{k=0}^4 \frac{\bar{u}_{t+k}^{FP}}{\prod_{j=0}^{k}(1-f_{t+j})}, \label{eq:claims}
\end{eqnarray}
where $u^{CC}_t$ is the number of continuing claims at month $t$, $f_t$ is the probability to find a job in month $t$,\footnote{We assume that during the first 26 weeks that an individual is receiving benefits the only reason why the claimant stops receiving benefits is because of a transition to employment. Alternatively, the claimant could stop collecting benefits without finding a job. We do not expect this to be quantitatively important.}

$u^{FP}_{t,\tau}$ is the number of those who received the final payment in week $\tau$ of month $t$, and $\bar{u}_{t}^{FP}$ is the total number of those who received their final payments after the reference week in month $t$ but not later than the reference week in month $t+1$.  For example, in a month with five weeks, the measure would include final payments in weeks 3-5 of that month, and then the final payments in the first two weeks of the subsequent month, i.e.
\begin{equation*}
\bar{u}_{t}^{FP}=\sum_{\tau=3}^{5}u^{FP}_{t,\tau} + \sum_{\tau=1}^{2}u^{FP}_{t+1,\tau}.
\end{equation*}
Rearranging Eq. (\ref{eq:claims}) we can express the job-finding rate as
\begin{eqnarray}
f_t= 1-\frac{\bar{u}^{FP}_{t}+\sum_{k=1}^4 \frac{\bar{u}_{t+k}^{FP}}{\prod_{j=1}^{k}(1-f_{t+j})}}{{u}_t^{CC}-u^{FP}_{t,2}},
\end{eqnarray}
yielding a system of equations with four more unknowns than equations. Setting terminal conditions for the final four measures of the job-finding rate we can solve backwards and calculate all of the preceding job-finding rates. We use the national job-finding rate from the CPS for the first four months of 2013 as our terminal conditions.\footnote{Alternatively, we could back out the job-finding rate from the LAUS unemployment data.  Doing so does not affect our results.}  We then perform our analysis through 2012:Q3, dropping the observations with the terminal conditions.


Implementing our estimation procedure using the baseline specification in Eq. (\ref{eq:baseline_spec}) to measure the effect of benefit extensions on this job-finding rate, we find that the coefficient $\alpha_f=-0.0486$, with a p-value of $0.03$. The coefficient on benefit duration in the regression for the job-finding rate multiplied by one minus the unemployment rate should approximately equal the coefficient in the regression for unemployment, i.e. $\alpha_f(1-u)\approx -\alpha_u$.  Using the value for the average unemployment rate of $7.01\%$ over the period, we get $\alpha_u=0.0452$. Using instead LAUS county unemployment data to re-estimate our baseline specification on unemployment on this sample we find a nearly identical coefficient on benefit duration of $0.0475$, with a p-value of $0$. Thus, the effect of unemployment benefit extensions on unemployment as measured with administrative claims data is quantitatively consistent with the effect measured with LAUS county-level unemployment data.\footnote{Note that our results imply that the ratio of county LAUS unemployment over county claimants does not respond to changes in benefits as it would if county unemployment was imputed using state-level unemployment.}

In addition to confirming the appropriateness of using LAUS data for our baseline empirical strategy, this result also highlights the quantitative importance of the equilibrium impact of benefit extensions on job creation. Our baseline analysis using LAUS data measures the effects of benefit extensions on all of the unemployed, including those who are ineligible to receive benefits. That measure therefore is a combination of the macro effect (which affects all unemployed) and the differential micro effect on search behavior of unemployed who are either eligible or ineligible to receive benefits. Using claims data, however, we exclusively focus on the unemployed who are eligible. If the micro effect was quantitatively important, the estimated coefficient on the claims sample should be significantly different from the one on the LAUS sample, since even at the depths of the recession the fraction of unemployed receiving claims did not exceed half of the unemployed. Our finding of similar effects of benefit extension on all unemployed in LAUS data and on benefit recipients in claims data suggests only a small role of micro elasticity, confirming existing estimates based on the data from the CPS.\footnote{This conclusion is further confirmed in Section \ref{sec:macro_effects} of the main text, where we estimate the effect of unemployment benefit extensions on the overall job-finding rate measured using county-level job vacancy data and find it to be the same as the effect on the job-finding rate of benefit claimants documented in this Appendix.}



%%%%%%%%%%%%%%%%%%%%%%%%%%%%%
%---------------------------%
%---------------------------%
%%%%%%%%%%%%%%%%%%%%%%%%%%%%%
\clearpage
\section{Additional Discussion of Endogeneity}
\label{app:endogeneity}

%%%%%%%%%%%%%%%%%%%%%%%%%%%%%
%---------------------------%
%%%%%%%%%%%%%%%%%%%%%%%%%%%%%
\subsection{Some Informal Examples}
\label{app:endogeneity_examples}

The key identification problem that our methodology was designed to overcome is the positive feedback from state unemployment to benefit extensions. This is guaranteed in our setting if shocks that affect state unemployment (and consequently benefits) impact symmetrically the two counties on the two sides of the state border. In this case, when we difference the variables between these two counties, the shock is differenced out so the only remaining difference between counties is the difference in benefits (that was induced by that shock).

To make things concrete, consider the following stylized (and not necessarily factually correct) examples. The discovery of hydraulic fracking was a big positive shock to the economy of Pittsburgh, PA during the Great Recession. This shock raised productivity and lowered unemployment in Pennsylvania relative to New Jersey. The border counties of Philadelphia, PA and Camden, NJ benefited equally from the availability of now cheaper natural gas. Faced with the same negative aggregate shock during the Great Recession, due to the feedback effect from unemployment to benefits, New Jersey raised unemployment benefits while Pennsylvania did not. Now consider the difference between Philadelphia and Camden Counties. They are affected by the same aggregate shock and have the same positive effect from the access to natural gas from Pittsburgh. These effects cancel out so that the difference in unemployment between them comes only from the fact that one has higher benefits than the other. Similarly, consider the effects of Hurricane Sandy that was a much bigger shock for New Jersey than Pennsylvania. The fact that this makes it more likely that benefits would be changed in NJ is not relevant because the storm had a similar impact on Philadelphia and Camden. Note that for this argument it is quite irrelevant whether Philadelphia is large or small relative to Camden or relative to Pennsylvania.

Alternatively, imagine a more idiosyncratic shock to counties. Suppose the Interstate Highway I-95 that is a major route connecting Philadelphia and Camden to other states is temporarily closed down for mechanical reasons. This will be a major shock to both counties. However, this shock is more likely to trigger benefit extensions in Pennsylvania since Philadelphia is much larger and important to Pennsylvania's economy than Camden is relative to New Jersey. Does this present a challenge to our identification strategy? No. The difference between the Philadelphia and Camden counties still arises only because of benefits and not because of underlying economic conditions.

The identification tests developed in the main body of the paper formally verify that shocks to economic conditions of the two states indeed have a symmetric impact on the two border counties, except through their effect on benefits.

%%%%%%%%%%%%%%%%%%%%%%%%%%%%%%
%----------------------------%
%----------------------------%
%%%%%%%%%%%%%%%%%%%%%%%%%%%%%%
%\clearpage
\subsection{Endogeneity test using Bartik shocks and state unemployment}
\label{app:endogeneity_test_unemp}

In this appendix we discuss our second implementation of the endogeneity test from Section \ref{sec:endogeneity_test} using state-level unemployment instrumented with Bartik shocks. Specifically, we estimate the following specification:
\begin{equation} \label{eq:endogenousFE_Bartik}
\Delta \tilde{x}_{p,t} =  \alpha \Delta b_{p,t} +\lambda_{p} + \chi \Delta_{t,p} u_{t,s}  + \tilde{\nu}_{p,t},
\end{equation}
where $\Delta_{t,p} u_{t,s} $ is the difference of the change in unemployment across the states the counties in the pair belong to instrumented with Bartik shocks. The results of implementing this test in the data yield a statistically insignificant coefficient of $0.043$ ($p$-value $0.90$) on the instrumented difference in state unemployment and including it in the regression affects neither the size nor the high statistical significance of the coefficient on benefits, once again consistent with the benchmark estimate being unbiased.

We instrument state unemployment with Bartik shocks because the endogeneity test must be based on an exogenous variable reflecting state-level economic conditions. As discussed in the related literature section, the empirical literature typically relies on the state-level unemployment rate for this purpose. This raises the question whether the difference in non-instrumented state-level unemployment rates can be used directly in place of $\Delta Z_{p,t}$ when testing for endogeneity. The answer is no, and the logic is very simple. State unemployment is endogenous to benefits. When benefits are raised in a state, unemployment increases in every county of the state. Moreover, unemployment is a slow moving state (in a mathematical sense) variable in response to benefits. Thus, while benefits change abruptly, county and state unemployment co-move in response to this change. As a consequence, if one regresses county unemployment on benefits and state unemployment, benefits are not relevant, as all their impact is summarized by state unemployment. In other words, state unemployment is the only variable that has a predictive power for county unemployment in such a regression while benefits - that drive both the county and state unemployment - are completely wiped out. We illustrate this point using data simulated from the model calibrated in Appendix \ref{sec:simulation}. The model imposes exogeneity because the county and state-level shocks are orthogonal, yet the coefficient of benefits is erroneously estimated to be zero when non-instrumented state unemployment is included in the regression. Note that the same logic applies to the evaluation of any state-level policy that affects unemployment throughout the state. The effect of any such policy change would be entirely (but erroneously) attributed to state unemployment if it is included in the regression and is not instrumented.

We simulate data from our calibrated model where we impose exogeneity - i.e. we assume the productivity processes at the county and state level are independent.  In Figure \ref{fig:endogeneity_test_unemp} we plot the time series for state unemployment, county unemployment and weeks of benefits available.  Notice that both state and county-level unemployment are smooth-moving variables, whereas the weeks of benefits jump when a benefit extension is triggered.  The correlation between state and county unemployment is significantly higher than between county unemployment and benefits, and controlling for state unemployment completely takes out the effect of benefits.  However, it is important to note that the only channel through which the state economy affects the county economy is through the benefit policy (because in this example the productivity processes are orthogonal). Thus, controlling for state unemployment, which is endogenous to benefits, is not a valid test for exogeneity.

\vspace{0.5cm}\begin{figure}[ht]
\centering
\includegraphics[width=0.8\columnwidth]{StateUCrop.pdf}
\vspace{-0.3cm}\caption{County and State Unemployment: Model.}
\label{fig:endogeneity_test_unemp}
\vspace{-0.3cm}
\end{figure}



%%%%%%%%%%%%%%%%%%%%%%%%%%%%%
%---------------------------%
%---------------------------%
%%%%%%%%%%%%%%%%%%%%%%%%%%%%%
\subsection{Alternative ``Endogeneity Test''}
\label{app:Hall_endogeneity}

\citet{hall:13} proposed trying to detect the presence of contamination on county-level LAUS unemployment data with state level variable by studying the correlation between the county unemployment rate measured in LAUS and state unemployment rate (relative to the correlation with the adjacent county's unemployment rate). To implement this suggestion, as proposed by \citet{hall:13}, we use LAUS data and regress county unemployment on state unemployment and adjacent county unemployment in 2007. The estimated coefficient on the state is 0.951 with a standard error of 0.015 and the coefficient on the adjacent out-of-state county is 0.316 with a standard error of 0.008.

To test whether the large estimate on the state-level variable implies the presence of contamination we perform the same analysis with administrative unemployment insurance claims data (that are not subject to any imputation). Specifically, we perform the same regression with continuing claims divided by population. The estimated coefficient on the state is 1.121 with a standard error of 0.045 and the coefficient on the adjacent out-of-state county is 0.415 with a standard error of 0.024. Thus, the results of these experiments do not provide support for the hypothesis that LAUS county-level unemployment contains an imputed state-level component.

The fact that state unemployment is an important predictor of county unemployment is also not very surprising and does not indicate shortcomings of our methodology. To see this, consider a purposefully simplified example that describes an ideal scenario with no endogeneity problem but where state unemployment is nevertheless a significant predictor. Suppose the shocks to the state, including the border county in that state and its adjacent out-of-state county have a common component. Each county is also subject to independent idiosyncratic shocks and measurement error in local unemployment. When we difference between the two border counties, the common component of the shock is differenced out implying no endogeneity concern. But consider now the regression of border county unemployment on unemployment in the state and adjacent out-of-state county. As a state consists of many counties, idiosyncratic shocks and measurement error are largely eliminated. Thus, the regression includes two regressors that have the same informational content for the dependent variable, but one of the regressors (state unemployment) is measured with little error while the other (out-of-state county unemployment) with lots of error. Clearly, the regressor measured with little error will be a significant predictor. Thus, such a regression cannot be interpreted as a test of endogeneity. It only provides evidence of the importance of common components and the distribution of errors at various levels of aggregation.

In addition, states differ with respect to many policies, e.g., taxes, regulations, UI and other benefit policies that apply to all locations within a state. Thus, it is expected that there is a difference in unemployment between border counties even when benefit durations are the same in both of them. Our estimation is based on a panel with large variation over time and these differences are accounted for by controlling for pair fixed effects (the factor model) and explicitly controlling for changes in these policies.

Further, since we know that county-level QCEW employment data is also not subject to any imputation using state-level variables, we now repeat the same experiment using 2007 LAUS unemployment and QCEW employment data.  In order to make the regressions comparable, we run them in logs of the level of unemployment and employment, as opposed to rates (because the calculation of employment rate would involve a potentially imputed county labor force variable which is only available in LAUS data the quality of which we aim to ascertain). The results of the regressions are displayed in the Table \ref{app_tab:A-1}. Both for county employment and unemployment we find large, positive significant coefficients on both the other county and on the state-level variables, and that the coefficient on the state is larger than that of the adjacent county.\footnote{To understand the relationship between these estimates in levels of unemployment to the estimates in rates reported above, note that since the level of a state variable is about two orders of magnitude larger than the county, in logs this would translate into roughly a factor of two difference on the state variable, which is what we observe.} Both pairs of coefficient estimates have overlapping 95\% confidence intervals. This experiment once again reinforces the conclusion that the large coefficient on the state variables is not evidence of an imputation problem. Instead, it is likely picking up a relationship between county and state variables driven by real economic factors and the fact that county-level variables are noisier than the state-level ones.

\begin{table}
  \centering
  \caption{County Unemployment and Employment as a Function of Pair County and State}\label{app_tab:A-1}
\input{LAUS_Imputation.tex}
\vspace{-0.2cm}
\end{table}

While the results of these test support the conclusion that LAUS county-level unemployment data are not inappropriately imputed using state-level variables, we find the analysis of endogeneity in the main text of the paper to provide a more direct and powerful verification of the lack of an endogeneity problem in our analysis. Moreover the fact that the results using county-level LAUS unemployment data are quantitatively consistent with the results based on independent datasets on vacancies, employment, and unemployment insurance claims strongly points to the same conclusion. This is comforting not only from the point of view of our research design but also because LAUS estimates are the basis for determining local spending under many State and Federal programs.




%%%%%%%%%%%%%%%%%%%%%%%%%%%%%
%---------------------------%
%---------------------------%
%%%%%%%%%%%%%%%%%%%%%%%%%%%%%
\clearpage
\section{Unemployment Benefit Extensions and Unemployment: Effect of Distance}
\label{app:benefits_and_unemployment_distance}
In this section we provide a more detailed analysis of the distance between border counties on the estimates. The median distance between population centers of border counties is close to 30 miles. The number of county pairs falls off sharply as we restrict to the distance to be either below 20 or above 50 miles. Thus, we re-estimate the benchmark specification on subsets of county pairs with the distance of more or less than 20, 30, 40, or 50 miles. The results are reported in Table \ref{tab:Effect_of_Distance}. As discussed in the main text, the closer are the border counties to each other, the larger is the downward bias on the estimated effect of UI extensions on unemployment because it tends to become easier to live in one county and work in the other (and unemployment measure is residence-based). This effect is indeed born out by the data as we see the estimated effects of unemployment benefit extensions on unemployment generally increasing in the distance between counties. The bias is quite small relative to the size of the estimate, however.

\vspace{0.5cm}\begin{table}[ht]
  \centering
   \caption{Unemployment Benefit Extensions and Unemployment: Effect of Distance} \label{tab:Effect_of_Distance}
\small
\input{Effect_of_Distance.tex}
\vspace{-0.2cm}
\end{table}

%%%%%%%%%%%%%%%%%%%%%%%%%%%%%%%%
%------------------------------%
%------------------------------%
%%%%%%%%%%%%%%%%%%%%%%%%%%%%%%%%
%\clearpage
\section{Alternative Separation Rate Measure}
\label{app:alt_sep}

In this Appendix, we discuss how we construct our alternative separation measure using QWI data. We construct pair-level separation rates by taking the ratio of separations during the quarter to total jobs at the beginning of the quarter (summing across both counties in the pair). Given that the QWI data appear to systematically overestimate separations relative to all other benchmark labor-market data sources, we rescale the QWI county separation rates to have the same mean as the separation rate from JOLTS. We then construct our quasi-differenced variances using this rescaled separation rate. Note that this preserves the covariance between benefits and county separations, allowing the quasi-differenced variables to respond to county-level movements in the separation rates. We report the results of our benchmark unemployment experiments using the county-level separation rates  in Table \ref{tab:Benefits_on_unemp_beg} and of the macro effects experiments in Table \ref{tab:macroeffects_beg}.

\begin{table}[H]
  \centering
  
   \caption{Unemployment Benefit Extensions and Unemployment: QWI Separations}\label{tab:Benefits_on_unemp_beg}
   
\small
\input{Benefits_on_unemp_beg.tex}
\parbox{16.5cm}{\noindent Note - $p$-values (in parentheses) calculated via bootstrap. Bold indicates $p<0.01$.\\
Column (1) - Baseline sample,\\
Column (2) - Baseline sample controlling for State GDP per worker,\\
Column (3) - Scrambled border county pairs sample,\\
Column (4) - Scrambled border county pairs sample controlling for State GDP per worker,\\
Column (5) - Sample of border counties with  $<15\%$ share of state's employment,\\
Column (6) - Sample of border counties with similar industrial composition,\\
Column (7) - Sample of border counties with population centers $<30$ miles apart,\\
Column (8) - Sample of border counties within the same Core Based Statistical Areas,\\
Column (9) - Baseline sample with perfect foresight measure of available benefits,\\
Column (10) - Baseline sample with controls for all other state-level policies.
}
\vspace{-0.0cm}
\end{table}



\begin{table}[H]
  \centering
   \caption{Unemployment Benefit Extensions and Job Creation: QWI Separations}\label{tab:macroeffects_beg}
   
\small
\input{macroeffects_beg.tex}
\vspace{-0.2cm}
\end{table}

\FloatBarrier

%%%%%%%%%%%%%%%%%%%%%%%%%%%%%
%---------------------------%
%---------------------------%
%%%%%%%%%%%%%%%%%%%%%%%%%%%%%
%\clearpage
\section{Unemployment Benefit Extensions and Unemployment: Controlling for Other State-Level Policies}
\label{app:benefits_and_unemployment_state_policies}

In this Appendix we control for tax and transfer policies that might be correlated with unemployment and unemployment benefit extensions at the county or state levels.


%%%%%%%%%%%%%%%%%%%%%%%%%%%%%
%---------------------------%
%%%%%%%%%%%%%%%%%%%%%%%%%%%%%
\subsection{Controlling for the Expansion of Food-Stamps Programs}


\citet{mulligan:12} has argued that in addition to unemployment benefit extensions, the Department of Agriculture's food-stamp program, now known as the Supplemental Nutrition Assistance Program, or SNAP, was also expanded considerably following the Great Recession. It is possible that the expansion of this program at the state level was correlated with unemployment benefit extensions so that the results reported above combine the effects of these programs. We now isolate their impacts.

Food-stamps were originally designed as a means-tested program for the poor. During the Great Recession the Federal government has allowed states to adopt  broad eligibility criteria that effectively eliminated the asset test and states received waivers from work requirements for the participants in the program. As a result, the participation in the program increased dramatically so that by 2010 half of non-elderly households with an unemployed head or spouse were receiving food stamps, with large variation across states.

To asses the extent to which the effects of unemployment benefit extensions documented above are affected by the expansion of eligibility for food-stamps program, we obtained USDA's SNAP Policy Database which documents policy choices of each state at monthly frequency.\footnote{http://www.ers.usda.gov/data-products/snap-policy-database.aspx} We construct a dummy variable equal to one during all periods when states use broad-based categorical eligibility to eliminate the asset test. The variable is zero otherwise. We include this variable in our baseline regression and report the results in Column (2) of Table \ref{tab:Benefits_on_unemp_SNAP_Mortgage}. The results confirm the argument in \citet{mulligan:12} that the expansion of food-stamps eligibility represents a marginal tax on working and thus leads to an increase in unemployment. It is, however, only weakly correlated with unemployment benefit extensions and thus does not significantly affect our estimate of their impact.



\begin{table}[t]
\centering
%Unemployment Benefit Extensions and Unemployment: \\ \hspace{2cm}
   \caption{Controlling for State SNAP and Foreclosure Policies }\label{tab:Benefits_on_unemp_SNAP_Mortgage}
\small
\input{Benefits_on_unemp_SNAP_Mortgage.tex}
\vspace{-0.2cm}
\end{table}


%%%%%%%%%%%%%%%%%%%%%%%%%%%%%%%%%%%%%%%%%%%%%%%%%%%%%%%%%%%%%%%%%%%%%%%%%%%%%%%%%%
%--------------------------------------------------------------------------------%
%%%%%%%%%%%%%%%%%%%%%%%%%%%%%%%%%%%%%%%%%%%%%%%%%%%%%%%%%%%%%%%%%%%%%%%%%%%%%%%%%%
\subsection{Controlling for Variation in State Foreclosure Policies}

The Great Recession has began with a sharp but heterogeneous across states decline in house prices. The government has responded by introducing various mortgage modification programs with the objective of helping underwater mortgagors. Various of these programs were either asset-tested or designed to write down mortgage principle to ensure that housing costs did not exceed a certain proportion of household income. In a series of papers, \citet{mulligan:08, mulligan:09, mulligan:10} has noted that this represents an implicit subsidy to unemployed workers. Moreover, \citet{herkenhoff+ohanian:13} have argued that the duration of the foreclosure process has been extended considerably following the Great Recession and that unemployed mortgagors use their ability to skip payments without being foreclosed upon as an implicit loan subsidy negatively affecting their job search and acceptance decisions.


\citet{cordell_etal:13} use proprietary data to measure the heterogeneity in foreclosure delay following the Great Recession across states. They find that in judicial states, in which state law requires a court action to foreclose, the delay is much larger than in statutory foreclosure states that do not require judicial intervention. Our use of the interactive effects estimator was specifically motivated by the concerns that aggregate shocks, such as shocks to house prices, may have heterogeneous impacts across border-county pairs depending, in part, on their state foreclosure law.  To verify the performance of the estimator, we define a dummy variable taking the value of one for border counties belonging to states with judicial foreclosure laws and zero otherwise. We then include in the benchmark specification the difference of the value of this dummy between border counties $i$ and $j$ in pair $p$.

The results reported in Column (4) of Table \ref{tab:Benefits_on_unemp_SNAP_Mortgage} indicate that this variable (the difference of the two dummies) is not statistically significant and does not affect the estimate of the effect of unemployment benefit extensions. This finding does not imply that foreclosure delay was not an important determinant of unemployment. It only means that our interactive effects estimator accounted for some of this aspect of heterogeneity across states and it did not impart a bias on our estimate of the effect of unemployment benefit extensions.



%%%%%%%%%%%%%%%%%%%%%%%%%%%%%%%%%%%%%%%%%%%%%%%%%%%%%%%%%%%%%%%%%%%%%%%%%%%%%%%%%%
%--------------------------------------------------------------------------------%
%%%%%%%%%%%%%%%%%%%%%%%%%%%%%%%%%%%%%%%%%%%%%%%%%%%%%%%%%%%%%%%%%%%%%%%%%%%%%%%%%%
\subsection{Controlling for the Effect of Stimulus Spending}


\begin{table}[t]
\centering
%Unemployment Benefit Extensions and Unemployment: \\
   \caption{Controlling for State Tax and Spending Policies}\label{tab:Benefits_on_unemp_stimulus_taxes}
   \footnotesize
\input{Benefits_on_unemp_stimulus_taxes.tex}
\parbox{16.5cm}{\noindent Note - $p$-values (in parentheses) calculated via bootstrap. Bold font indicates $p<0.05$.}
%\parbox{16.5cm}{\noindent }
\vspace{-0.2cm}
\end{table}


In the specification of Column (2) of Table \ref{tab:Benefits_on_unemp_stimulus_taxes} we control for the effects of stimulus spending. We use data on actual county level spending arising from the American Recovery and Reinvestment Act (ARRA) - commonly referred to as the ``stimulus package.'' We obtain an accounting of all stimulus spending at the zip code level under the ARRA.\footnote{www.recovery.gov.}  We then match counties to zip codes.  We run our specification both in levels and by dividing the spending by the population in the county, obtained from the Census. We find that that controlling for ARRA spending does not affect our estimate of the effect of unemployment benefit extensions.\footnote{The coefficient on spending however has to be interpreted with caution. It is conceivable, in contrast to unemployment benefits which depend on economic conditions at the state level, that spending at the county level depends
on the economic conditions at the county level. In this case the coefficient on spending will be biased.}


%%%%%%%%%%%%%%%%%%%%%%%%%%%%%%%%%%%%%%%%%%%%%%%%%%%%%%%%%%%%%%%%%%%%%%%%%%%%%%%%%%
%--------------------------------------------------------------------------------%
%%%%%%%%%%%%%%%%%%%%%%%%%%%%%%%%%%%%%%%%%%%%%%%%%%%%%%%%%%%%%%%%%%%%%%%%%%%%%%%%%%
\subsection{Controlling for State Tax Policies}

To control for the variation in state-level tax policies we obtained detailed Census Bureau data on quarterly tax revenues for each state.\footnote{http://www.census.gov/govs/qtax/} We consider whether effective total or sales tax rates have co-moved systematically with unemployment benefit durations. We find no support for this hypothesis. The results reported in Table \ref{tab:Benefits_on_unemp_stimulus_taxes} imply that directly controlling for these effective tax rates has virtually no impact on our estimates of the effect of unemployment benefit extensions on unemployment.

Our analysis was based on effective tax rates for two reasons. First, the statutory rates have not changed systematically over our sample period. Despite many states having balanced budget laws, expansions of unemployment benefits have not required changes in tax rates as extensions were mostly federally financed. Second, there are numerous state programs targeted to attract businesses that offer tax deductions to individual firms. For competitive reasons details of such policies are rarely disclosed. We can effectively measure them, however, by focusing on actual tax receipts.


%%%%%%%%%%%%%%%%%%%%%%%%%%%%%%%%%%%%%%%%%%%%%%%%%%%%%%%%%%%%%%%%%%%%%%%%%%%%%%%%%%
%--------------------------------------------------------------------------------%
%%%%%%%%%%%%%%%%%%%%%%%%%%%%%%%%%%%%%%%%%%%%%%%%%%%%%%%%%%%%%%%%%%%%%%%%%%%%%%%%%%
\subsection{Controlling for Other State Policies}


While we found no evidence that the effects of unemployment benefit extensions on unemployment are a proxy for changes in other tax policies, we now consider whether they could be driven by other state policies, such as changes in regulatory or litigation environment. For this purpose we obtain data from three prominent indexes of state policies - U.S. State Business Policy Index (SBSI), State Business Tax Climate Index (SBTCI), and BHI State Competitiveness Index (BHI).\footnote{www.sbecouncil.org, www.taxfoundation.org, www.beaconhill.org, respectively.}  The construction of these indexes is based on a well-documented methodology, the data is available annually over our sample period, and can be made consistent over time. A description of these indexes, the analysis of their predictive performance for state economic outcomes, and references to other academic evaluations can be found in \citet{kolko+neumark+mejia:13}.


\begin{table}[t]
\centering
   \caption{Controlling for Other State Policies}\label{tab:Benefits_on_unemp_other_policies}
\small
\input{Benefits_on_unemp_other_policies.tex}
\vspace{-0.2cm}
\end{table}


%

The motivation for using these broad policy indexes was provided in \citet{holmes:98}, who found that controlling for a similar (but no longer available) index of state policies accounted for the positive relationship between right-to-work laws and manufacturing employment. This suggests that the conclusion about the effects of one policy may be misleading without taking into account other state policies reflected in a broad index. In contrast, the results reported in Table \ref{tab:Benefits_on_unemp_other_policies} imply that controlling for such indexes does not affect the measured impact of unemployment benefit extensions on unemployment.



%%%%%%%%%%%%%%%%%%%%%%%%%%%%%%%%
%------------------------------%
%%%%%%%%%%%%%%%%%%%%%%%%%%%%%%%%
\section{Unemployment Benefit Extensions and Wages: Controlling for Selection}
\label{app:wages}

In this appendix, we revisit the effects of unemployment benefits on wages by controlling for two types of selection. First, we study the effects of benefits on wages of stayers. Second, we measure the effects of benefits on the wages of new hires.

%%%%%%%%%%%%%%%%%%%%%%%%%%%%%%%
%-----------------------------%
%%%%%%%%%%%%%%%%%%%%%%%%%%%%%%%
\subsection{Accounting for Selection: Wages of Job Stayers.} 

The idiosyncratic productivity of workers moving from non-employment to employment or from job to job depend on business cycle conditions \citep{gertler+trigari:09, hagedorn+manovskii:13}. We can decompose the idiosyncratic component $\phi^i$ into permanent worker ability $\mu^i$, job specific productivity $\kappa^i$ and a stochastic component $\epsilon^i$:
\beq
\log(\phi^i_t) = \log(\mu_t^i) + \log(\kappa_t^i) + \log(\epsilon^i_t).
\eeq
The decision of a non-employed to accept a job depends on $z_t$, $\mu_t^i$, the job-specific productivity $\hat{\kappa}$ as well as on benefits $b$. The decision of a worker to switch jobs depends on the worker's current job specific productivity $\kappa_t^i$ and the job-specific productivity in the new job $\hat{\kappa}$.
Productivity $\hat{\kappa}$ is a random draw from a distribution $F$. A worker who has received $N$ offers during a period accepts the highest draw $\kappa$, which is distributed according to $F^N$. Since the $F^N$ are ordered by first-order stochastic dominance, the expected value of $\kappa$ is increasing in $N$ and is thus increasing in the number of vacancies. A more generous unemployment insurance system leads to a drop in vacancy posting and therefore to fewer offers and a lower expected value of $\kappa$.
By the Law of Large Numbers, workers starting a new job in a recession or when benefits are high then have a lower average value of $\kappa$ than workers starting a job when many offers are available such as in a boom or when benefits are low. Thus, if we regress wages on benefits we also pick up the impact of benefits on the average value of $\kappa$.\footnote{Benefits may also affect $\kappa$ by making liquidity constrained workers more selective in the jobs they accept.}
To deal with this issue, we follow \citet{hagedorn+manovskii:13} and consider job stayers, defined as workers who have the same job in periods $t$ and $t+1$ and thus also the same value of $\kappa$. Taking differences across time for a job stayer yields
\beq
\log (w_{t+1}^i) - \log (w_t^i) & = & \beta_z (\log(z_{t+1}^a) - \log(z_t^a))
+ \beta_{\theta} (\log(\theta_{t+1}^a) - \log(\theta_t^a)) \\
& +&  \beta_b (\log(b_{t+1}^a)-log(b_t^a)) + \log(\epsilon^i_{t+1}) -  \log(\epsilon^i_t) + \eta_{t+1}^i - \eta_{t}^i, \notag
\eeq
or, denoting by $\bar{\Delta} x $ the difference between the value of variable $x$ in period $t+1$ and $t$ (not to be confused with $\Delta$ used to denote the difference between border counties),
\beq
\bar{\Delta}log (w^i) =  \beta_z \bar{\Delta} \log(z^a)
+ \beta_{\theta} \bar{\Delta} \log(\theta^a)
+ \beta_b \bar{\Delta} \log(b^a)
+ \bar{\Delta} \log(\epsilon^i) + \bar{\Delta} \eta^i,
\eeq
that is the terms $\mu^i$ and $\kappa^i$ drop out.
We therefore consider a group of workers in county $a$ who worked in period $t$ and $t+1$ for the same employer with average wages $w^a_{t}$ in period $t$ and $w^a_{t+1}$ in period $t+1$. Theory then predicts that regressing the difference in wages, $\bar{\Delta} \log(w^a)$, on the difference in benefits, $\bar{\Delta} \log(b)$, yields a positive coefficient. We again have to control for the endogeneity of policy and to this end we again invoke assumption (\ref{assu:identify_main}) and consider the difference across paired border counties. We then implement the same regression as Eq (\ref{eq:wageequation_estimate}).





The coefficient $\tilde{\beta}_b$ captures the equilibrium wage response which, using (\ref{eq:wage_thetab}), combines the direct effect of benefits on wages, $\beta_b$, and the indirect effect of benefits on market tightness $\theta$, $\beta_{\theta} \beta_{\theta,b}$, where $\beta_{\theta,b}$ is the regression coefficient of market tightness on benefits,
\beq
\tilde{\beta}_b = \beta_b + \beta_{\theta} \beta_{\theta,b}. \notag
\eeq
We, therefore, obtain instead of the potentially large direct effect $\beta_b$ the smaller equilibrium response $\tilde{\beta}_b$, which takes into account the benefit-induced change in market tightness and its effect on wages. It is the latter, the equilibrium response, $\tilde{\beta}_b$, which theory predicts to be positive.

To implement this procedure, we measure wages of job stayers using the QWI. We begin with a measure of full-quarter employment - workers who remained employed at the same firm for the entire quarter - and average wage earnings of full-quarter employees. However, in quarter $t$ the measure of full quarter employment also includes workers who will separate in $t+1$, and in quarter $t$ the measure includes new hires from quarter $t$.  Thus, to isolate the wages of stayers we difference out the average wages of $t+1$ separators from the average wages in $t$ and difference out the average $t$ new hire wages from the average wages in $t+1$. This yields the true average wages of stayers in quarters $t$ and $t+1$. In the regression, to reduce measurement error we accumulate this stayer-wage double difference over the current and the next three quarters, measuring the benefit differential over the corresponding window (from $t-1$ to $t+2$) --- the analogue, for the noisier stayer-wage series, of the two-quarter averaging used for all workers in the main text.

Columns (1) and (2) of Table \ref{tab:app_Benefits_on_Wages} shows the result for raw wages and wages adjusted for the UI payroll tax, respectively. We find that wages statistically significantly increase in response to an increase in benefits. To assess the implied quantitative magnitude, consider a typical county pair in the Great Recession. The estimate in Column (1) implies that a county with $70$ weeks of benefits has a $0.78\%$ higher level of wages than a county with $50$ weeks of benefits, everything else equal.\footnote{$\exp(0.0232*(\log(70)-\log(50)))
= 1.0078.$} As we explain in Appendix \ref{app:wage_channel}, the increase in wages of job stayers indicates that benefit extensions increase the outside option available to these workers when they bargain on the job, consistent with the existing US laws and UI system regulations.

\begin{table}[t]
  \centering
   \caption{Unemployment Benefit Extensions and Wages} \label{tab:app_Benefits_on_Wages}
\small
\input{app_Benefits_on_Wages.tex}
\vspace{-0.2cm}
\end{table}



%%%%%%%%%%%%%%%%%%%%%%%%%%%%%%%%%%%%%%%%%%%%%%%%%%%%%%%%%%%%%%%%%%%%%%%%%%%%%%%%%
%-------------------------------------------------------------------------------%
%%%%%%%%%%%%%%%%%%%%%%%%%%%%%%%%%%%%%%%%%%%%%%%%%%%%%%%%%%%%%%%%%%%%%%%%%%%%%%%%%
\paragraph{Accounting for Selection: Wages of New Hires.}
It is well known that wages of new hires are more cyclically sensitive than wages of job stayers, e.g., \citet{haefke+sonntag+vanrens:12}, \citet{pissarides:09}. A common interpretation of this finding is based on the idea that firms provide long term (implicit) contracts to workers which partially insulate workers from aggregate labor market conditions. This literature interprets the idiosyncratic component $\phi^i$ as being a function of past aggregate labor market variables since the start of the job. For newly hired workers who negotiate new employment contracts, however, wages depend on the current conditions only so that the term $\phi^i$ equals zero. Under this theory, we can obtain an unbiased estimate of the effect of unemployment benefits on wages by restricting the sample to newly hired workers only.

Similar to what we have done above, and re-using the same notation, denote by $\bar{\Delta} \log(w^a)$ the difference between the average wage of workers newly hired in period $t+1$ and the average wage of workers newly hired in period $t$ in county $a$. Theory then predicts that regressing the difference in wages, $\bar{\Delta} \log(w^a)$, on the corresponding difference in benefits, $\bar{\Delta} \log(b)$, yields a positive coefficient. To control for the endogeneity of policy we consider the difference across paired border counties. Taking differences across counties $a$ and $b$ in the same pair $p$ of $\bar{\Delta}\log(w^a)$ and $\bar{\Delta}\log(w^b)$ yields once again Eq. (\ref{eq:wage_thetab}) with the only difference that the average wages of job stayers are replaced with the average wages of new hires. Implementing the regression in Eq. (\ref{eq:wageequation_estimate}) of the double difference of wages of new hires on the double difference in benefits, we obtain the equilibrium response of the wages of newly hired workers which once again combines the direct effect of benefits on wages, $\beta_b$, and the indirect effect of benefits on market tightness, $\beta_{\theta} \beta_{\theta,b}$:
\beq
\tilde{\beta}_b = \beta_b + \beta_{\theta} \beta_{\theta,b}. \notag
\eeq

The results of implementing this regression using the QWI data on the wages of newly hired workers are reported in Columns (3) and (4) of Table \ref{tab:Benefits_on_Wages}, for raw wages and wages adjusted for the UI payroll tax, resp. We find that wages of newly hired workers statistically significantly increase in response to an increase in benefits. Under the implicit contracting theory of the labor market, this is the wage response relevant for the decision to create a new job. As the expected labor costs associated with filling new jobs rise, the number of jobs created can be expected to fall.



%%%%%%%%%%%%%%%%%%%%%%%%%%%%%%%%%%%%%%%%%%%%%%%%%%%%%%%%%%%%%%%%%%%%%%%%%%%%%%%%%%
%--------------------------------------------------------------------------------%
%--------------------------------------------------------------------------------%
%%%%%%%%%%%%%%%%%%%%%%%%%%%%%%%%%%%%%%%%%%%%%%%%%%%%%%%%%%%%%%%%%%%%%%%%%%%%%%%%%%
%\clearpage
\section{Additional Comments on the Wage Channel}
\label{app:wage_channel}

It is clear that unemployment benefit extensions directly affect wages of newly hired workers by raising the value of the outside option of declining a job (there is, of course, the offsetting equilibrium effect that with lower job availability newly hired workers might be hired into lower quality jobs). For the incumbent workers this effect might be less clear because to exercise their outside option they must quit the job and in theory quitters are not eligible for unemployment benefits. This potentially raises two questions. First, if the wage channel does not operate for the incumbent workers, is there a need for our methodology based on quasi-differencing the variables of interest. Second, how to interpret our empirical finding that wages of incumbent workers respond positively to benefit extensions. We address these questions in this Appendix.


First, we explain that the issue of whether unemployment benefit extensions affect wages of incumbent workers is irrelevant for our methodology and findings. We assume that the wage of a new hire will respond to future determinants, including expected changes in unemployment benefits, productivity, demand, etc. This assumption is clearly not controversial. It is the expected productivity, the expected wages and the expected workers' value of unemployment which matter for the firm's decision to post a vacancy. These expected values depend on both current and future values.

Consider two neighboring counties $A$ and $B$ with the same unemployment benefits and the same worker productivity in the current period but where productivity is known to increase in county $A$ next period and to remain unchanged in county $B$. Benefits remain unchanged in both counties. Obviously, the current period incentives to post vacancies (create jobs) are higher in county $A$ than in county $B$ although current period values of benefits and productivity are identical. Our quasi-difference estimator accounts for this formation of expectations.

The same argument applies to benefits. Suppose it is expected that next period productivity remains the same but UI benefits increase in county $A$. This future increase in benefits negatively affects vacancy posting in county $A$ in the current period. This happens because the workers' value of remaining unemployed increases today. In \citet{hall+milgrom:08} this also increases the (newly hired) worker's payoff while bargaining.\footnote{A worker and a firm who start bargaining but do not reach an agreement in period $t$, continue bargaining in period $t+1$ and are in a bargaining situation similar to the one of a worker and a firm who meet and start bargaining in period $t+1$.}
This increases the wage and lowers profits and thus fewer vacancies are posted in county $A$ in the current period. Our quasi-difference estimator accounts for this effect of expectations.\footnote{To put it differently, dropping the quasi-difference estimator would yield a coefficient which is an uninterpretable convolution of the current and future county-differences in UI benefits.}

Note that these arguments do not invoke the assumption that an outside option continues to be available when the incumbent worker bargains on the job. The need for the quasi-difference estimator is independent of that assumption. The presence of continuous bargaining (where the outside option remains available) may make the dependence on future values quantitatively stronger. No matter how strong the expectation effect is, however, it has to be accounted for by the estimation strategy.

Furthermore, we find that wages of job stayers (in the same firm in consecutive periods), wages of new hires, and overall wages do increase when benefits are extended. This is an empirical fact that the empirical strategy to estimate the contemporaneous effects of benefits on unemployment, vacancies, tightness and employment should also be consistent with, and our quasi-differenced specification is.

We now turn to the second question of whether our finding that wages of incumbent workers respond to unemployment benefit extension are consistent with economic theory and its interpretation.

The result that unemployment benefit extensions affect wages of incumbent workers is clearly and firmly rooted in economic theory. Consider, for example, the efficiency wage models that have been the workhorse model of wage setting for a number of recent decades. The central element of these models is that the effort of an employee is not fully observable and is not verifiable by a court. Thus, if workers are dismissed, they will certainly be entitled to unemployment benefits. When the outside option of (incumbent) workers improves, they have to be paid more to exert effort. Thus, any such model has the implication that an increase in the outside option leads to a wage increase. Whether the employer pays a higher piece rate, a larger bonus or just increases the wage is irrelevant for our methodology and our results. %This is so standard that even 30 years ago Lawrence Summers argued (with much consensus)  during the General Discussion of a paper by Gary Burtless at the 1983 Brookings Institution that raising benefit duration puts upwards pressure on wages.

The bargaining models have the same implication. Viewed through the lens of these models, our empirical finding that wages of stayers respond to benefits implies that the outside option continues to be available when the worker bargains on the job. This model of wages is certainly widely accepted as it is used not only in the standard \citet{pissarides:00} textbook but also in prominent papers on the subject, e.g., \citet{hall+milgrom:08} and even \citet{hall:05}. One may question the assumption underlying this literature by arguing that to exercise the outside option, the worker would need to quit, and quitters do not receive UI benefits. However, this assertion does not fully reflect reality in the U.S. labor market. In particular, it is hard to tell apart quits and layoffs and the burden of proof is on the employer. We will provide some examples below, but it is clear that many employers will not be able or willing to contest UI claims of employees. Contesting is costly even in normal times but especially during the Great Recession employers' incentives to engage in providing such evidence have been presumably negligible when benefit extensions are paid by the federal government. Our empirical findings suggest that quitters receive, at least with some probability, UI benefits. This conclusion is supported by our analysis of the legal features of the UI system summarized below.



%%%%%%%%%%%%%%%%%%%%%%%%%%%%%%%%%%%%%%%%%%%%%%%%%%%%%%%%%%%%%%%%%%%%%%%%%%%%%%%%%%
%--------------------------------------------------------------------------------%
%%%%%%%%%%%%%%%%%%%%%%%%%%%%%%%%%%%%%%%%%%%%%%%%%%%%%%%%%%%%%%%%%%%%%%%%%%%%%%%%%%
\subsection{Elements of California Unemployment Insurance Law}
\label{sec:law}

We now discuss some legal details on the eligibility of workers for benefits in the State of California. UI policies and procedures are state-dependent but the general principles are similar. Much of the discussion is copied verbatim from the State of California Benefit Determination Guide, an eight-volume compendium, designed to present definitive discussions on points of unemployment insurance law for the field office determination interviewer.

The basic line of argument in this section is as follows.
\begin{enumerate}
\item As a general rule, voluntary quitters are not entitled to benefits. In Section \ref{app:mutual_misunderstanding} we provide examples illustrating the difficulties in establishing whether a voluntary quit has occurred.

\item In Section \ref{app:eligible_quitters} we explain that even if the quit is voluntary in the sense that the employer had the job available for the worker and had no intention of firing the employee, the quit may not be considered voluntary from the point of view of the Unemployment Insurance laws and regulations. If the employee can argue that he had a good reason for leaving the employer, he will be entitled to benefits. We provide a small subset of such reasons that illustrate the potential for uncertainty on the part of the employer as to whether the separating employee will be able to collect benefits. This is sufficient to explain why employers accede to wage demands of incumbent workers when the generosity of the UI system increases.

\item In Section \ref{app:misconduct} we argue that an improvement in the generosity of the UI system strengthens workers' hand in bargaining with the employer through an additional channel. Instead of the threat of outright quitting, the worker can implicitly threaten the employer to induce a firing. While workers fired for misconduct are not eligible for benefits, establishing misconduct is very difficult, in part due to the necessity of proving that misconduct was willful, and  the burden of proof is on the employer. It seems likely that many employers would have little ability, resources, or economic motivation, to contest such cases in the courts. It may well be cheaper to accept workers' wage demands instead.

\end{enumerate}


\subsubsection{Was the Separation a Quit?}
\label{app:mutual_misunderstanding}
This is not very straightforward to establish. For example, if separation is due to mutual agreement or mutual misunderstanding the worker is eligible for benefits. In particular, when both parties have a reasonable but mistaken belief of the others understanding of the separation, the claimant is not subject to disqualification. In addition, there may be a separation by mutual agreement if the employer and the employee have mutually agreed to separate, either at the time of the termination, or initially, at the time of hire. In such cases the termination is neither a discharge nor a leaving and thus a disqualification cannot arise under Section 1256.

The following Precedent Decisions illustrate:\footnote{Precedent decisions refer to the body of case law that is developed through the adjudicatory process at the California Unemployment Insurance Appeals Board (CUIAB) and contains the Appeals Board's definitive expression on unemployment matters. The Unemployment Insurance Code specifically authorizes CUIAB Board Members to consider, decide and designate as precedent decisions those cases that contain a significant legal or policy determination of general application that is likely to recur. CUIAB, its administrative law judges, and the Employment Development Department Director are controlled by these precedents, except as modified by judicial review.}


In P-B-253, the claimant's attendance became irregular because of poor health. Her union contract provided for a leave of absence for a maximum of two years.
The claimant was carried on the employer's ``absent-sick service payroll'' from January to March. In March the claimant contacted her supervisor, saying she was still ill and didn't know when she would be able to return. During the course of the interview, she and the employer agreed that the claimant's separation ``might be the best thing to do.'' Neither suggested the leave continue. In its decision, the Board said:
\begin{quotation}
...[T]he evidence before us justifies a conclusion that the conversation ... resulted in a mutual agreement between the claimant and her employer that under the circumstances no useful purpose would be served by the indefinite extension of her then existing leave of absence. Under these facts, we hold that the claimant's abandonment of the employer-employee relationship ... was with good cause... .
\end{quotation}

Some separations appear insolvable from the standpoint of a misunderstanding between the claimant and the employer, in which each thinks the other has been the moving party in the separation. In cases such as the following, the Board has considered the separation to be neither a quit nor a discharge.

In P-B-458, the claimant had been counseled concerning his job performance some five weeks prior to the separation. On the last day of his employment, he was called to a meeting with the president and general manager. At the meeting, the claimant remarked that if he were in charge he would place the blame for slow business upon himself. The president felt the claimant had not been working to capacity, and the claimant specifically recalled that the president told him they ``should part company.'' Shortly after that, the claimant announced that he would be leaving. The claimant cleaned out his desk and left. The employer interpreted the conversation and events as a resignation, while the claimant felt he had been discharged. In its decision, the Board stated:
\begin{quotation}
The record does not sufficiently reflect that either the claimant or the employer was the moving party. We hold that where the claimant and the employer are mutually but reasonable mistaken about the other party's understanding of the separation, the claimant is not subject to disqualification under Section 1256 of the Code.
\end{quotation}

\noindent {\bf Important Caveat.} In fairness, we must admit, however, that we could not find a precedent decision clarifying how a separation upon an exogenous separation shock after $78^{th}$ round of Hall-Milgrom bargaining would be treated... Separation by mutual agreement or mutual misunderstanding?



\subsubsection{Can Voluntary Quitters Be Eligible for Benefits?}
\label{app:eligible_quitters}

As a general rule, voluntary quitters are not entitled to benefits. There are many exceptions, though, to which a worker voluntary leaving his or her job may appeal in order to receive benefits. In this section we mention some of many such reasons. The point of this discussion is that there is at least a chance, and perhaps a sizable one, that a worker might be able to collect benefits even in the event of quitting. Even on its own, this is sufficient to explain why a more generous UI system induces a higher equilibrium wage even for incumbent workers.

\begin{enumerate}

\item Section VQ90 A: Conscientious objection.

When directly related to working conditions, a conscientious objection is considered to be a compelling reason for restricting availability for work or for voluntarily leaving work. Conscientious objection means an objection by an individual to performing an act that individual sincerely believes is wrong. The objection may be based on ethical, moral, religious, or philosophical grounds.

Title 22, Section 1256-6 (b), provides:
\begin{quotation} ... If an individual has, or after working a time newly acquires a conscientious objection to the work condition or assigned work based on religious beliefs founded on the tenets or beliefs of a church, sect, denomination, or other religious group, or on ethical or philosophical grounds, an individual's voluntary leaving of the most recent work based on religious beliefs or other grounds is with good cause...
\end{quotation}
The degree to which the claimant's beliefs are commonly held or considered reasonable by others is immaterial.

\item Section 1256, VQ150:

AA. All Reasonable Transportation Alternatives Exhausted

The claimant quit your employment because of a lack of transportation. There is no adequate public transportation available and the claimant had exhausted all alternatives before quitting. Available information shows that the claimant had good cause for leaving work.

BB. Commute Time Excessive

The claimant quit your employment because of the commute time required. Available information shows that the claimant had good cause for leaving work.

CC. Moved - No Transfer Available

The claimant quit your employment to move. He/she could not have transferred to a job site nearer to the new home. Available information shows that the claimant had good cause for leaving work.

EE. Transportation Costs Too High

The claimant quit your employment because the transportation costs were too high. Available information shows the claimant had good cause for leaving work.


\item Section 1256, VQ 155

AA. Compelling Domestic Obligations

The claimant quit your employment because of domestic reasons. Available information shows that the claimant had good cause for leaving work.

BB. Moving After Marriage - Outside Normal Commute Area

The claimant quit your employment to move with his/her spouse to a place outside the normal commute area. Available information shows that the claimant had good cause for leaving work.

CC. Moving After Marriage - No Transfer Available

The claimant quit your employment to move with his/her spouse. He/she was unable to transfer to another worksite nearer the new home. Available information shows that the claimant had good cause for leaving work.

DD. Family Illness or Death - No Leave Available

The claimant quit your employment because of a family illness/death. Available information shows that the claimant had good cause for leaving work

EE. Unemancipated Minor

The claimant quit your employment at the insistence of his/her parents. The claimant is a minor, subject to parental control. Available information shows that the claimant had good cause for leaving work.

FF. Domestic Violence Abuse - No Reasonable Alternative


\item Section 1256, VQ235

AA. Medical Advice to Quit

The claimant quit your employment on his/her doctor's advice. A leave of absence was not available or would not have resolved the problem. Available information shows that the claimant had good cause for leaving work.

BB. Reasonable Concern for Health or Safety

The claimant quit your employment because of a reasonable concern for his/her health or safety. Available information shows that the claimant had good cause for leaving work.

CC. Failure to Take Drug Test - Employer Request Unreasonable

The claimant quit your employment because he/she was asked to take a drug test. The claimant had not previously consented to the test and there was no reasonable suspicion that he/she was under the influence of drugs. Available information shows that the claimant had good cause for leaving work.
\end{enumerate}

There are many, many other reasons the worker can establish that a voluntary leave was for a good cause, including arguing that the workplace represented an intimidating, hostile or offensive working environment as illustrated by the following two Precedent Decisions.

In P-B-300, the claimant did establish real and compelling cause for her action. The claimant worked as a bookkeeper for a small insurance firm. She quit that employment because the employer repeatedly criticized her in a sarcastic manner in front of customers; some of the criticism was caused by errors made by the claimant in her work, but some criticism concerned matters not attributable to the claimant and some concerned matters wholly unrelated to the claimant's work. Occasionally, the claimant left the employer's office in tears. Three witnesses testified on behalf of the claimant. In finding the claimant eligible for benefits, the Board stated:
\begin{quotation}
... the record established that the conduct of the claimant's employer in the instant case was abusive and hostile, moreover, this conduct was repeated on numerous occasions. Under the circumstances this constituted a compelling reason for the claimant to leave her employment...
\end{quotation}

Thus, if undue embarrassment, or harassment is caused by continual criticism, in contrast to a single instance of criticism, good cause for leaving does exist. In P-B-475 the Board ruled that offensive (to the worker) behavior of sexual nature also constitutes a valid reason to leave employment because of offensive working environment.\\




\subsubsection{Was the Discharge for Misconduct?}
\label{app:misconduct}
As a general rule, employees discharged for misconduct are not eligible for benefits. Only those who were discharged not through the fault of their own are. The question we are interested in here is whether an employee can implicitly threaten the employer with misconduct during the wage bargaining. It appears that the answer is at least to some degree affirmative.


\vspace{0.5cm}\noindent{\bf What Constitutes Misconduct?}

\smallskip

For an employee to be discharged for misconduct, it has to be proven \emph{by the employer} that misconduct was willful. Where the element of willfulness is missing, the claimant's actions would generally not be misconduct. For example, according to Section 1256-30(b)(3) of Title 22, misconduct generally does not exist, because willfulness is missing, if the claimant:
\begin{itemize}
\item Has been merely inefficient.
\item Has failed to perform well due to inability or incapacity.
\item Has been inadvertent.
\item Has been ordinarily negligent in isolated instances.
\item Has made good faith errors in judgment or discretion.
\end{itemize}

The following sequence of examples illustrates.

\begin{enumerate}

\item Example - Inefficiency 1:

In P-B-222, the claimant was a pasteurizer for a large creamery. Prior to the claimant's discharge, there had been several discussions between the superintendent and the claimant in connection with the quality of the claimant's services. Although the claimant testified that his work improved after those discussions, his superintendent believed that the claimant had failed to improve sufficiently to warrant keeping the claimant. The principal complaint against the claimant appears to be a failure to pasteurize milk on occasions at proper temperatures and that the claimant at times held milk in the vats an excessive time, resulting in the milk acquiring an undesirable flavor. In one instance, about three hundred gallons of milk were spoiled due to improper pasteurization, resulting in a considerable financial loss to the employer. In finding the claimant eligible, the Board said:
\begin{quotation}
A careful review of the entire evidence in the instant matter does not disclose, in our opinion, more than inefficiency or unsatisfactory performance on the part of the claimant...  The record does not establish that the claimant wilfully or intentionally disregarded the employer's interest or that the occurrences forming the basis for the discharge were deliberate violations of standard good behavior...
\end{quotation}

\item Example - Inefficiency 2:

In P-B-184, the employer hired the claimant as a production worker after the claimant indicated that he had operated drill presses, lathes, punch presses, reamers, and similar equipment. He was assigned to work a drill press and found to be unsatisfactory. He was next assigned to a lathe and was moved from that job when he incorrectly loaded a part and wrecked a fixture which required several hours to rebuild. He was, thereafter, tried on several other jobs but failed to meet the employer's standards on any of them and was discharged about three weeks after being hired. The Board found him eligible and stated:
\begin{quotation}
The record does not establish that the claimant wilfully or intentionally disregarded the employer's interests, or that the occurrences forming the basis for the discharge were deliberate violations of standards of good behavior which the employer had a right to expect of his employee.
\end{quotation}

\item Example - Inability to Perform to Employer's Standard:

In P-B-224, the claimant was employed for four weeks as a bookkeeper, and let go because the employer considered that her work was not ``up to par.'' The Board found her eligible and stated:
\begin{quotation}
We find that the efficient cause of the claimant's discharge was her inability to satisfy the employer's standards in relation to the quality of her work ... mere ineptitude is not misconduct...
\end{quotation}

\item Example - Incapable of Meeting Standard:

The claimant, a tube-bender and assembler for an aircraft manufacturer, was discharged after six years' employment because of his inability to produce an acceptable amount of work on a swaging machine. He had been assigned to this new task for only four hours when he was given a ``correction interview.'' At this interview, he was informed that his production was 50 percent below standard and that he would be discharged unless he showed immediate improvement.

The employer contended that the claimant had deliberately ``stalled'' but was unable to substantiate such a statement. The claimant had performed satisfactorily on other operations, had even been graded ``excellent'' in production on other tasks. When the claimant was again assigned to the swaging machine the next workday, he refused the assignment as he knew that if he did not make the quota he would be fired. He was discharged as a result.

The claimant complied with the employer's orders when he was initially assigned to a new machine and according to the record he made every effort to become proficient in its operation. Because of his age and slight physical stature the claimant could foresee that he would not be able to operate the new machine to the satisfaction of the employer and felt justified in refusing the assignment.

In this case the discharge would not be for misconduct. The claimant was unable to meet the employer's standards because of his age and slight physical stature. It should also be noted that the employer did not give the claimant a sufficient amount of time to meet the standards (only four hours). Likewise, if an employer should fail to provide adequate equipment for doing the work or should set quantity standards so high that only the exceptional few could meet them, a failure to produce the required quantity of work would not be misconduct.


\item Example - Error in Judgment:

In P-B-195 the claimant, a cab driver, was discharged because of a traffic accident. At the time the claimant was hired, he received a course of instructions covering the company's rules and the motor vehicle laws with which he was expected to comply. Shortly after the end of the course, the claimant was involved in a minor accident when he backed into a parked car. He was warned that he would be discharged if involved in one more accident within a year. Several months later, the claimant was en route to pick up a passenger. He was driving approximately 40 feet behind another car, when he was hailed by someone on the left side of the street and glanced toward the person hailing him. He heard the screech of brakes, immediately looked to the front and applied his own brakes when he saw that the traffic in front of him had stopped. He was unable to stop before colliding with the car in front of him. The collision was observed by two police officers and the claimant was cited under Section 22350 of the California Vehicle Code. The Board found the claimant eligible and stated:
\begin{quotation}
In this case, the claimant was cited under Section 22350 of the California Vehicle Code. We do not consider the fact of citation controlling in this case, but only one of the factors which we must consider in arriving at our conclusion. The quoted section of the Vehicle Code is so phrased as to allow the driver of a vehicle to exercise judgment in the operation of such vehicle. Assuming that the claimant was careless as found by the traffic officers involved, his carelessness was, at most, an error of judgment. Admittedly, it was his fault that the collision occurred. However, he was following the vehicle preceding him at a reasonable distance and erred only when he withdrew his attention from the road when he was hailed by a person on the sidewalk. It appears to us that the claimant's action could readily be defined as a reflex action in response to the call, especially since it was the practice of the taxi drivers to seek to identify such a person so that the company could be informed of a possible customer.
\end{quotation}

\item Example - Isolated Incident of Ordinary Negligence 1:

In \emph{Silva v. CUIAB} (First Appellate Court, 1973), the claimant was being trained for new and unfamiliar work; he became nervous and frustrated and either ``blew up'' or felt he was going to blow up. He left work without permission in midafternoon. The employer was aware of some emotional problems the claimant was having. The employer spoke to the claimant the next morning about his unauthorized departure. The claimant's reply was sarcastic and, when told if such action was repeated he would be discharged, he responded with a vulgar remark. He was told if that was the way he felt, he could leave, whereupon he left. He would have been discharged for his attitude and language that morning had he not left. The court held:
\begin{quotation}
Given the tests of fault and wilful or wanton behavior as essential elements of 'misconduct', the single instance of an offensive remark ... uttered in the circumstances disclosed falls within the category of a mere mistake or error in judgment, a 'minor pecadillo' and is not misconduct disqualifying appellant from unemployment insurance benefits.
\end{quotation}

\item Example - Isolated Instance of Ordinary Negligence 2:

The claimant was hired to drive his employer's new cars from a freight depot to the company's storage warehouse. The automobiles were shipped directly from the factory and were serviced as they were unloaded. The employer testified that oral warnings had been given all employees to check oil and water levels before driving the cars and that any driver who subsequently caused damage to a car would be discharged.

One of the automobiles the claimant was driving incurred engine damage because the car was driven with no oil in it. The claimant denied that he had been warned to check the oil and water levels before driving the vehicles. Additionally, there was dispute as to whether the oil gage was operating correctly.

The claimant's contented that this was an ``isolated'' incident and that he had acted unknowingly and without evil design or intent. Because of the dispute as to the employer's warning to check oil and water levels and the working condition of the oil gage, it cannot be shown that there was wilful negligence. The discharge would not be for misconduct.


\item Example - Action Not Willful:

In \emph{Maywood Glass Co. v. Stewart} (1959), the claimant was discharged because she packed defective glassware on several occasions. The employer testified that she had been warned several times she would be discharged if she persisted.

The claimant denied such warnings were given. The claimant stated she packed bad glassware because of the rapidity in which they were working. She also had a headache. The court held her discharge was not for misconduct and stated:
\begin{quotation}
Moreover, even if the claimant had been warned, the evidence does not compel a finding that she was guilty of 'misconduct' within the meaning of the statute. Although (claimant) admitted packing defective bottles, she denied that she had intentionally done so. (Claimant) worked the 'graveyard shift' from midnight to 8 o'clock in the morning. She testified that on the night in question she was suffering from a headache and that there was a high percentage of defective glassware coming down the line. In these circumstances the trier of fact could reasonably conclude that her conduct did not constitute 'misconduct' within the meaning of the statute...
\end{quotation}


\end{enumerate}



\vspace{0.5cm}\noindent{\bf Burden of Proof and Presumption of Eligibility}

\smallskip

Section 1256 of the UI Code provides in part:
\begin{quotation}
An individual is presumed to have been discharged for reasons other than misconduct in connection with his or her work and not to have voluntarily left his or her work without good cause unless his or her employer has given written notice to the contrary to the department as provided in Section 1327, setting forth facts sufficient to overcome the presumption. The presumption provided by this section is rebuttable.
\end{quotation}

In \emph{Perales v. California Department of Human Resources Development} (1973), the Appellate court held that because the presumption in Section 1256 was established to implement the public policy of prompt payment of benefits to the unemployed so as to reduce the suffering caused thereby (Section 100 of the UI Code), the presumption affects the burden of proof. To overcome the presumption the employer and the Department must prove that the claimant was discharged for misconduct in connection with his or her work by a preponderance of the evidence.

This is also the position in the following court decisions:
\begin{itemize}
\item In \emph{Maywood Glass Co. v. Stewart} (1959), the Court stated that the employer has the burden of establishing ``misconduct'' to protect its reserve fund.
\item In \emph{Prescod v. California Unemployment Insurance Appeals Board} (1976), the Court held that the burden of disqualification is on the employer or the Department, and not the claimant.
\end{itemize}

\noindent {\bf The punch line:}  Proving misconduct is costly for the employer. And the employee can probably make it not worthwhile for many employers by threatening to drag the process out through multiple appeal procedures involving testimony and provision of documentation and witnesses by the employer.



%%%%%%%%%%%%%%%%%%%%%%%%%%%%%%%%%%%%%%%%%%%%%%%%%%%%%%%%%%%%%%%%%%%%%%%%%%%%%%%%%%
%--------------------------------------------------------------------------------%
%--------------------------------------------------------------------------------%
%%%%%%%%%%%%%%%%%%%%%%%%%%%%%%%%%%%%%%%%%%%%%%%%%%%%%%%%%%%%%%%%%%%%%%%%%%%%%%%%%%
\clearpage
\section{Appendix Figures}


\begin{figure}[ht]
    \begin{center}
  	\vspace{0.0cm}
        \scalebox{0.32}{\includegraphics{countymap2.png}}
    \vspace{0.0cm}
     \caption{Map of U.S.A. with state and county outlines.}
        \label{fig:Counties_Map}
\end{center}
\end{figure}


\clearpage
  \thispagestyle{empty}
%\begin{changemargin}{0in}{-2cm}
\begin{figure}[ht]
\makebox[\textwidth]{\parbox{1.5\textwidth}{
\centering
 %   \begin{center}
	\vspace{-1.5cm}
%	\hspace{-1in}
\subfigure[June 2008.]
    {

        \scalebox{0.38}{\includegraphics{AllMaps42.jpg}}
        \label{fig:Benefit_Maps_42}
    }
\hspace{0.0cm}
\vspace{-0.25cm}
\subfigure[December 2008.]
    {

        \scalebox{0.38}{\includegraphics{AllMaps48.jpg}}
        \label{fig:Benefit_Maps_48}

    }
\hspace{0.0cm}
%\vspace{-0.25cm}
\subfigure[June 2009.]
    {
        \scalebox{0.38}{\includegraphics{AllMaps54.jpg}}
        \label{fig:Benefit_Maps_54}
    }
\hspace{0.0cm}
%\vspace{-0.25cm}
\subfigure[December 2009.]
    {
        \scalebox{0.38}{\includegraphics{AllMaps60.jpg}}
        \label{fig:Benefit_Maps_60}
    }


\vspace{-0.25cm}
  \subfigure[June 2010.]
    {
        \scalebox{0.38}{\includegraphics{AllMaps66.jpg}}
        \label{fig:Benefit_Maps_66}
    }
\vspace{-0.25cm}
   \subfigure[December 2010.]
    {
        \scalebox{0.38}{\includegraphics{AllMaps72.jpg}}
        \label{fig:Benefit_Maps_72}
    }
\vspace{-0.25cm}
  \subfigure[June 2011.]
    {
        \scalebox{0.38}{\includegraphics{AllMaps78.jpg}}
        \label{fig:Benefit_Maps_78}
    }
\vspace{-0.25cm}
   \subfigure[December 2011.]
    {
        \scalebox{0.38}{\includegraphics{AllMaps84.jpg}}
        \label{fig:Benefit_Maps_84}
    }

\vspace{0.5cm}
%\hspace{-1in}
\subfigure[June 2012.]
    {
        \scalebox{0.38}{\includegraphics{AllMaps90.jpg}}
        \label{fig:Benefit_Maps_90}
    }
\vspace{0.0cm}
\subfigure[December 2012.]
    {
        \scalebox{0.38}{\includegraphics{AllMaps96.jpg}}
        \label{fig:Benefit_Maps_96}
    }
    \vspace{-0.25cm}
     \caption{Unemployment benefit duration across U.S. states during the Great Recession. Selected months.}
    \label{fig:Benefit_Maps}
%\end{center}
}}
\end{figure}
%\end{changemargin}
\clearpage



%\end{appendix}




%%%%%%%%%%%%%%%%%%%%%%%%%%%%%%%%%%%%%%%%%%%%%%%%%%%%%%%%%%%%%%%%%%%%%%%%%%%%%%%%%%
%--------------------------------------------------------------------------------%
%--------------------------------------------------------------------------------%
%%%%%%%%%%%%%%%%%%%%%%%%%%%%%%%%%%%%%%%%%%%%%%%%%%%%%%%%%%%%%%%%%%%%%%%%%%%%%%%%%%
% \clearpage
% \section{TO DO}

% \begin{enumerate}
% \item Re-write introduction
% \item Decide what we want the baseline specification to be
% \item Clean up the section on the empirical methodology
% \item Trim anything that we can from the main text
% \end{enumerate}


\end{appendix}

\end{document}




